% Generated mechanically from frozen input p12/ko/build/ch111/exercises_full.tex.
% Do not edit this generated file; edit the bound locale source instead.

% OK, start here.
%

\setcounter{chapter}{110}
\chapter{연습문제}



\phantomsection
\label{exercises-section-phantom}


\section{대수}
\label{exercises-section-algebra}

\noindent
이 첫 절에는 여러 가지 잡다한 문제를 모아 두었다.

\begin{exercise}
\label{exercises-exercise-isomorphism-localizations}
$A$를 환, ${\mathfrak m}$을 극대 아이디얼이라 하자. $A[X]$에서
$\tilde {\mathfrak m}_1 = ({\mathfrak m}, X)$ 및
$\tilde {\mathfrak m}_2 = ({\mathfrak m}, X-1)$이라 하자. 다음을 보여라.
$$
A[X]_{\tilde {\mathfrak m}_1} \cong A[X]_{\tilde {\mathfrak m}_2}.
$$
\end{exercise}

\begin{exercise}
\label{exercises-exercise-coherent}
뇌터 환이 아닌 환 $R$로서, $R$의 모든 유한 생성 아이디얼이
$R$-가군으로서 유한 표시되는 예를 찾아라.
(마지막 성질이 성립하는 환을 {\it 코히어런트}라고 한다.)
\end{exercise}

\begin{exercise}
\label{exercises-exercise-flat-ideals-pid}
$(A, {\mathfrak m}, k)$가 뇌터 국소환이라고 하자. 임의의
유한 $A$-가군 $M$에 대하여 $r(M)$을 $A$-가군으로서 $M$을 생성하는
생성원의 최소 개수라 정의하자. NAK에 의하여 이 수는
$\dim_k M/{\mathfrak m} M = \dim_k M \otimes_A k$와 같다.
\begin{enumerate}
\item $r(M \otimes_A N) = r(M)r(N)$임을 보여라.
\item $r(I) > 1$인 아이디얼 $I\subset A $를 잡자. 다음을 보여라.
$r(I^2) < r(I)^2$.
\item $A$의 모든 아이디얼이 평탄 가군이면 $A$가 주 아이디얼 정역
(또는 체)임을 결론지어라.
\end{enumerate}
\end{exercise}

\begin{exercise}
\label{exercises-exercise-not-isomorphic}
$k$를 체라 하자. 다음 $k$-대수의 각 쌍이 동형이 아님을 보여라.
\begin{enumerate}
\item 임의의 $n\geq 1$에 대하여 $k[x_1, \ldots, x_n]$과
$k[x_1, \ldots, x_{n + 1}]$.
\item $k[a, b, c, d, e, f]/(ab + cd + ef)$와 $k[x_1, \ldots, x_n]$,
여기서 $n = 5$.
\item $k[a, b, c, d, e, f]/(ab + cd + ef)$와 $k[x_1, \ldots, x_n]$,
여기서 $n = 6$.
\end{enumerate}
\end{exercise}

\begin{remark}
\label{exercises-remark-simple-geometric}
물론 이 연습문제의 취지는 각 경우에 ``큰'' 정리를 적용하기보다
간단한 논증을 찾는 데 있다. 그럼에도 일반 원리를 지침으로 삼는 것은 좋다.
\end{remark}

\begin{exercise}
\label{exercises-exercise-silly}
대수. (사소하며 쉬워야 한다.)
\begin{enumerate}
\item 환 $A$와, 분할되지 않는 다음과 같은 $A$-가군의 짧은 완전열의
예를 들어라.
$$
0 \to M_1 \to M_2 \to M_3 \to 0.
$$
\item 위와 같이 분할되지 않는 $A$-가군의 열과 충실히 평탄한
$A \to B$의 예로서, 다음 열이 $B$-가군의 열로서는 분할되는 예를 들어라.
$$
0 \to M_1\otimes_A B \to M_2\otimes_A B \to M_3\otimes_A B \to 0.
$$
\end{enumerate}
\end{exercise}

\begin{exercise}
\label{exercises-exercise-field-kummer}
$k$가 1의 원시 $n$제곱근 $\zeta$를 갖는 체라고 하자. 이는
$\zeta^n = 1$이지만 $0 < m < n$이면 $\zeta^m\not = 1$이라는 뜻이다.
\begin{enumerate}
\item $k$의 표수가 $n$과 서로소임을 보여라.
\item $a \in k$가 $k$의 원소이고, $n \geq d > 1$인 $n$의 임의의 약수
$d$에 대하여 $k$에서 $d$제곱이 아니라고 하자. $k[x]/(x^n-a)$가
체임을 보여라. (힌트: $x^n-a$의 분해체를 생각하고 Galois 이론을 사용하라.)
\end{enumerate}
\end{exercise}

\begin{exercise}
\label{exercises-exercise-valuation}
$\nu : k[x]\setminus \{0\}  \to {\mathbf Z}$가 다음 성질을 갖는 사상이라고
하자. $f$, $g$가 영이 아니면 $\nu(fg) = \nu(f) + \nu(g)$이고,
$f$, $g$, $f + g$가 영이 아니면
$\nu(f + g) \geq min(\nu(f), \nu(g))$이며, 모든 $c\in k^*$에 대하여
$\nu(c) = 0$이다.
\begin{enumerate}
\item $f$, $g$, $f + g$가 영이 아니고 $\nu(f) \not = \nu(g)$이면
$\nu(f + g) = min(\nu(f), \nu(g))$임을 보여라.
\item $f = \sum a_i x^i$, $f\not = 0$이면
$\nu(f) \geq min(\{i\nu(x)\}_{a_i\not = 0})$임을 보여라. 언제 등호가
성립하는가?
\item $\nu$가 음의 값을 취하면 어떤 $n\in {\mathbf N}$에 대하여
$\nu(f) = -n \deg(f)$임을 보여라.
\item $\nu(x) \geq 0$이라고 하자. 다음 집합이 $k[x]$의 소 아이디얼임을
보여라.
$\{f \mid f = 0, \ or\ \nu(f) > 0\}$.
\item 가능한 모든 $\nu$를 기술하라.
\end{enumerate}
\end{exercise}


\noindent
$A$를 환이라 하자. {\it 멱등원}은 $e^2 = e$를 만족하는 원소
$e \in A$이다. 원소 $1$과 $0$은 언제나 멱등원이다.
{\it 비자명한 멱등원}은 영과 같지 않은 멱등원이다. 두 멱등원
$e, e' \in A$가 {\it 직교한다}는 것은 $ee' = 0$이라는 뜻이다.

\begin{exercise}
\label{exercises-exercise-product}
$A$를 환이라 하자. $A$가 두 영이 아닌 환의 곱일 필요충분조건은
$A$가 비자명한 멱등원을 갖는 것임을 보여라.
\end{exercise}

\begin{exercise}
\label{exercises-exercise-lift-idempotents}
$A$를 환, $I \subset A$를 국소 멱영 아이디얼이라 하자. 사상
$A \to A/I$가 멱등원들 사이의 전단사를 유도함을 보여라.
(힌트: 먼저 $I$가 멱영인 경우를 증명하는 편이 더 쉬울 수 있다.
그런 다음 ``절대 뇌터 축소''를 사용하여 멱영인 경우로 귀착하라.)
\end{exercise}








\section{여극한} \label{exercises-section-colimits}

\begin{definition} \label{exercises-definition-directed-poset} {\it 유향 집합}이란 준순서 $\leq$를 갖춘 공집합이 아닌 집합 $I$으로서, 임의의 쌍 $i, j \in I$에 대해, $i \leq k$ 및 $j \leq k$를 만족하는 $k \in I$가 존재하는 것을 말한다. $I$ 위의 {\it 환의 계}란, 각 $i \in I$에 대한 환 $A_i$과, $i \leq j$일 때의 환 준동형 $\varphi_{ij} : A_i \to A_j$로서, $i \leq j \leq k$일 때 합성 $A_i \to A_j \to A_k$가 $A_i \to A_k$와 같아지는 것들로 이루어진 것이다. \end{definition}

\noindent 마찬가지로 군의 계, 고정된 환 위의 가군의 계, 체 위의 벡터 공간의 계 등을 정의한다.

\begin{exercise} \label{exercises-exercise-directed-colimit} $I$를 유향 집합, $(A_i, \varphi_{ij})$를 $I$ 위의 환의 계라고 한다. 환 $A$과 사상 $\varphi_i : A_i \to A$가 있으며, 모든 $i \leq j$에 대해 $\varphi_j \circ \varphi_{ij} = \varphi_i$가 되고, 또한 다음과 같은 보편성을 가지는 것이 존재함을 보이라: 임의의 환 $B$과, 모든 $i \leq j$에 대해 $\psi_j \circ \varphi_{ij} = \psi_i$를 만족하는 사상 $\psi_i : A_i \to B$이 주어졌을 때, $\psi_i = \psi \circ \varphi_i$를 만족하는 유일한 환 준동형 $\psi : A \to B$가 존재한다. \end{exercise}

\begin{definition} \label{exercises-definition-colimit} 연습 \ref{exercises-exercise-directed-colimit} 으로 구성한 환 $A$을, 이 계의 {\it 여극한} 이라고 한다. 표기는 $\colim A_i$로 한다. \end{definition}

\begin{exercise} \label{exercises-exercise-prime-in-colimit} % P12-ERR-0419: first frozen reversed-order declaration; see ../control/ERRATA_CANDIDATES.jsonl.
$(I, \geq)$을 유향 집합, $(A_i, \varphi_{ij})$을 $I$ 위의 환의 계로 하고, 그 여극한을 $A$ 라 하자. 다음 전단사가 존재함을 증명하라: $$
\Spec(A) = \{(\mathfrak p_i)_{i \in I} \mid
\mathfrak p_i \subset A_i \text{ 이고 }
\mathfrak p_i = \varphi_{ij}^{-1}(\mathfrak p_j)\ \forall i \leq j\}
\subset \prod\nolimits_{i \in I} \Spec(A_i)
$$ 오른쪽 집합은 집합들 $\Spec(A_i)$의 역극한이다. 표기는 $\lim \Spec(A_i)$로 한다. \end{exercise}

\begin{exercise} \label{exercises-exercise-colimit-surjective} % P12-ERR-0419: second frozen reversed-order declaration; see ../control/ERRATA_CANDIDATES.jsonl.
$(I, \geq)$를 유향 집합으로 하고, $(A_i, \varphi_{ij})$를 $I$ 위의 환의 계로 하며, 그 여극한을 $A$라고 하자. 모든 $i \leq j$에 대해 $\Spec(A_j) \to \Spec(A_i)$가 전사라고 가정하자. 이때, 모든 $i$에 대해 $\Spec(A) \to \Spec(A_i)$가 전사임을 보이라. (힌트: Tychonoff 정리를 사용해도 되지만, 대수, 보조정리 \ref{algebra-lemma-in-image}에 근거한 본질적으로 자명한 직접적인 대수적 증명도 있다.) \end{exercise}

\begin{exercise} \label{exercises-exercise-integral-colimit-finite} % P12-ERR-0420: frozen broken English parallel syntax; see ../control/ERRATA_CANDIDATES.jsonl.
$A \subset B$를 정수적 환 확대라 하자. $\Spec(B) \to \Spec(A)$가 전사임을 증명하라. 위 연습문제와, 유한 환 확대의 경우에 이 명제가 성립한다는 강의의 결과, 그리고 $B = \colim B_i$가 유한 확대 $A \subset B_i$들의 유향 여극한이라는 사실을 사용하라. \end{exercise}

\begin{exercise} \label{exercises-exercise-colimit-tensor} $(I, \geq)$를 유향 집합으로 하자. $A$를 환으로 하고, $(N_i, \varphi_{i, i'})$를 $I$로 첨자 붙은 $A$-가군의 유향계로 하자. $M$를 다른 $A$-가군으로 하자. 다음을 증명하라: $$
\colim_{i\in I} M \otimes_A N_i\cong
M \otimes_A \Big( \colim_{i\in I} N_i\Big).
$$ \end{exercise}

\begin{definition} \label{exercises-definition-finite-presentation} $R$ 위의 가군 $M$가 $R$ 위에서 {\it 유한 표현}이라는 것은 유한 자유 가군의 사상 $ R^{\oplus n} \to R^{\oplus m}$의 여핵과 동형임을 의미한다. \end{definition}

\begin{exercise} \label{exercises-exercise-colimit-modules} 임의의 환 위의 임의의 가군이 다음과 같이 표현될 수 있음을 증명하라: \begin{enumerate} \item 그 유한 생성 부분 가군들의 여극한, \item 어떤 방식으로든 유한 표현 가군들의 여극한. \end{enumerate} \end{exercise}


\section{가법 범주와 Abel 범주} \label{exercises-section-additive}

\begin{exercise} \label{exercises-exercise-filtered-vector-spaces} $k$을 체로 한다.$\mathcal{C}$를 $k$ 위의 여과 벡터 공간의 범주로 한다. 임의의 범주의 여과 대상의 정의에 대해서는 호몰로지, 정의 \ref{homology-definition-filtered} 를 참조하라. \begin{enumerate} \item 이것이 가법 범주임을 보여라(두 대상의 직접 합이 무엇인지를 주의 깊게 설명하라). \item $f : (V, F) \to (W, F)$를 $\mathcal{C}$의 사상으로 하자. $f$가 핵과 여핵을 가짐을 보여라($f$의 핵과 여핵이 무엇인지를 정확히 설명하라). \item 표준 사상 $\Coim(f) \to \Im(f)$가 동형이 아닌 $\mathcal{C}$의 사상의 예를 들어라. \end{enumerate} \end{exercise}

\begin{exercise} \label{exercises-exercise-torsion-free} $R$를 뇌터 정역으로 하자. $\mathcal{C}$를 유한 생성 비틀림 없는 $R$-가군의 범주로 하자. \begin{enumerate} \item 이것이 가법 범주임을 증명하라. \item $f : N \to M$를 $\mathcal{C}$의 사상으로 하자. $f$가 핵과 여핵을 가진다는 것을 보이라($f$의 핵과 여핵을 반드시 정확히 정의하라). \item 표준 사상 $\Coim(f) \to \Im(f)$이 동형이 아닌 뇌터 정역 $R$와 $\mathcal{C}$의 사상의 예를 들라. \end{enumerate} \end{exercise}

\begin{exercise} \label{exercises-exercise-other} 가법적이며 핵과 여핵을 가지지만, 연습 \ref{exercises-exercise-filtered-vector-spaces} 및 \ref{exercises-exercise-torsion-free} 의 것과는 다른 범주의 예를 들어라. \end{exercise}


\section{텐서곱} \label{exercises-section-tensor-product}

\noindent 텐서곱은 대수, 절 \ref{algebra-section-tensor-product}에서 도입하였다. $R$를 환으로 하고, $\text{Mod}_R$를 $R$-가군의 범주로 한다. 함자 $F : \text{Mod}_R \to \text{Mod}_R$에 대해 다음과 같이 말한다: \begin{enumerate} \item 임의의 $R$-가군 $M, N$에 대해 $F : \Hom_R(M, N) \to \Hom_R(F(M), F(N))$가 아벨 군의 준동형일 때, 가법적이라고 한다. 호몰로지, 정의 \ref{homology-definition-preadditive}를 참조하라. % P12-ERR-0421: frozen missing-predicate defect; see ../control/ERRATA_CANDIDATES.jsonl.
\item 임의의 $R$-가군 $M, N$에 대하여 $F : \Hom_R(M, N) \to \Hom_R(F(M), F(N))$가 $R$-선형일 때, $R$-선형이라고 한다. \item 임의의 짧은 완전열 $0 \to M_1 \to M_2 \to M_3 \to 0$에 대하여 열 $F(M_1) \to F(M_2) \to F(M_3) \to 0$가 완전할 때, 오른쪽 완전하다고 한다. \item 임의의 짧은 완전열 $0 \to M_1 \to M_2 \to M_3 \to 0$에 대해 열 $0 \to F(M_1) \to F(M_2) \to F(M_3)$가 완전할 때, 왼쪽 완전하다고 한다. \item 집합 $I$과 $R$-가군 $M_i$이 주어졌을 때, 사상 $F(M_i) \to F(\bigoplus M_i)$이 동형 $\bigoplus F(M_i) = F(\bigoplus M_i)$을 유도한다면, 직합과 가환한다고 한다. \end{enumerate}

\begin{exercise} \label{exercises-exercise-characterize-tensor-functor} $R$를 환으로 하고, 위의 표기법을 사용한다. \begin{enumerate} \item $R$-선형이 아닌 가법적 함자를 $F : \text{Mod}_R \to \text{Mod}_R$ 갖는 환 $R$의 예를 들어라. \item $N$를 $R$-가군이라고 하자. 함자 $F(M) = M \otimes_R N$가 $R$-선형이며 오른쪽 완전임을 보이고, 직합과 가환함을 증명하라. \item 반대로, $R$-선형이며 오른쪽 완전이고, 직합과 가환하는 임의의 함자 $F : \text{Mod}_R \to \text{Mod}_R$는, 어떤 $R$-가군 $N$에 의해 $F(M) = M \otimes_R N$의 형태로 나타낼 수 있음을 증명하라. \item (3)에서, $F$가 직합과 가환한다는 가정을 제거하면, 결론이 더 이상 성립하지 않음을 보여라. \end{enumerate} \end{exercise}


\section{평탄한 환 사상} \label{exercises-section-flat}

\begin{exercise} \label{exercises-exercise-localization-flat} $S$를 환 $A$의 곱셈적 부분집합으로 하자. \begin{enumerate} \item $A$-가군 $M$에 대해서 $S^{-1}M = S^{-1}A \otimes_A M$를 증명하라. \item $S^{-1}A$가 $A$ 위에서 평탄함을 증명하라. \end{enumerate} \end{exercise}

\begin{exercise} \label{exercises-exercise-examples-not-flat} 다음 각 경우에, $A$-가군의 단사 $M_1 \to M_2$ 이면서 $M_1\otimes N \to M_2 \otimes N$가 단사가 아닌 것을 찾아라: \begin{enumerate} \item $A = k[x, y]$ 그리고 $N = (x, y) \subset A$. (여기서도, 이하에서도 $k$는 체로 한다.) \item $A = k[x, y]$ 그리고 $N = A/(x, y)$. \end{enumerate} \end{exercise}

\begin{exercise} \label{exercises-exercise-flat-not-projective} % P12-ERR-0422: frozen mismatched-article coordination; see ../control/ERRATA_CANDIDATES.jsonl.
환 $A$과 유한 $A$-가군 $M$ 중 평탄하지만 사영적이지 않은 $A$-가군의 예를 들어라. \end{exercise}

\begin{remark} \label{exercises-remark-flat-not-projective} $M$가 유한 표현을 가지며, $A$ 위에서 평탄하다면, $M$는 $A$ 위에서 사영적이다. 따라서, 요구하는 예에서는 환 $A$은 뇌터 환이 아닐 필요가 있다. 나는 ${\mathbf R}$ 위의 ${\mathcal C}^\infty$-함수의 환을 $A$로 하는 예를 알고 있다. \end{remark}

\begin{exercise} \label{exercises-exercise-flat-not-free-dvr} ${\mathbf Z}_{(2)}$ 위의 평탄하지만 자유롭지 않은 가군을 찾아라. \end{exercise}

\begin{exercise} \label{exercises-exercise-flat-deformations} 평탄 변형. \begin{enumerate} \item $k$를 체로 하고, $k[\epsilon]$를 쌍대수 환 $k[\epsilon] = k[x]/(x^2)$로 하고, $\epsilon = \bar x$로 한다. 임의의 $k$-대수 $A$에 대해, $A$가 $B/\epsilon B$와 동형이 되도록 하는 평탄 $k[\epsilon]$-대수 $B$가 존재함을 보여라. \item $k = {\mathbf F}_p = {\mathbf Z}/p{\mathbf Z}$로 하고, $$
A = k[x_1, x_2, x_3, x_4, x_5, x_6]/
(x_1^p, x_2^p, x_3^p, x_4^p, x_5^p, x_6^p).
$$ $B/pB$가 $A$와 동형이 되도록 하는 평탄 ${\mathbf Z}/p^2{\mathbf Z}$-대수 $B$가 존재함을 보여라. (여기서 $p$는 $\epsilon$의 역할을 한다.) \item 이번에는 $p = 2$로 하고, $k = {\mathbf F}_2 = {\mathbf Z}/2{\mathbf Z}$ 및 $$
A = k[x_1, x_2, x_3, x_4, x_5, x_6]/
(x_1^2, x_2^2, x_3^2, x_4^2, x_5^2, x_6^2, x_1x_2 + x_3x_4 + x_5x_6).
$$에 대해 같은 문제를 생각한다. 그러나 이 경우에는, $B/2B$가 $A$와 동형이 되도록 하는 평탄 ${\mathbf Z}/4{\mathbf Z}$-대수 그러한 $B$는 {\it 존재하지 않음}을 보여라. (장치를 간파하라! 같은 예는 임의의 표수 $p > 0$에서 성립하지만, 계산은 더 어려워진다.) \end{enumerate} \end{exercise}

\begin{exercise} \label{exercises-exercise-flat-given-residue-field-extension} $(A, {\mathfrak m}, k)$를 국소환으로 하고, $k'/k$를 유한 차수 체 확대으로 한다. 국소환의 평탄한 국소 준동형 $A \to B$이 존재함을 보여라. ${\mathfrak m}_B = {\mathfrak m} B$가 되며, 또한 $B/{\mathfrak m} B$가 $k$-대수로서 $k'$와 동형이 되는 경우이다. (힌트: 먼저 $k \subset k'$가 하나의 원소로 생성되는 경우를 다루어라.) \end{exercise}

\begin{remark} \label{exercises-remark-flat-given-residue-field-extension-general} 동일한 결과는 임의의 체 확장 $K/k$에 대하여 성립한다. \end{remark}


\section{환의 스펙트럼} \label{exercises-section-spectrum-ring}

\begin{exercise} \label{exercises-exercise-spec-Z} $\Spec(\mathbf{Z})$을 집합으로 계산하고, 그 위상을 서술하시오. \end{exercise}

\begin{exercise} \label{exercises-exercise-basis-opens-standard} $A$을 임의의 환으로 하자. $f\in A$에 대해 $D(f):= \{\mathfrak p \subset A \mid f \not \in \mathfrak p\}$로 정한다. 열린 집합 $D(f)$가 $\Spec(A)$의 위상의 기저를 이룬다는 것을 증명하시오. \end{exercise}

\begin{exercise} \label{exercises-exercise-radical-ideals-closed} 함수 $I\mapsto V(I)$가 자연스러운 전단사 $$
\{I\subset A\text{ 이고 }I = \sqrt{I}\}
\longrightarrow
\{T\subset \Spec(A)\text{ 는 닫힌 집합}\}
$$를 정의한다는 것을 증명하시오. \end{exercise}

\begin{definition} \label{exercises-definition-quasi-compact} 위상 공간 $X$가 {\it 준콤팩트} 이라는 것은, 임의의 열린 덮개 $X = \bigcup_{i\in I} U_i$에 대해, 유한한 부분 집합 $\{i_1, \ldots, i_n\}\subset I$가 존재하여 $X = U_{i_1}\cup\ldots U_{i_n}$가 되도록 하는 것을 말한다. \end{definition}

\begin{exercise} \label{exercises-exercise-spec-quasi-compact} 임의의 환 $A$에 대해 $\Spec(A)$가 준콤팩트임을 증명하시오. \end{exercise}

\begin{definition} \label{exercises-definition-Hausdorff} 위상 공간 $X$가 분리공리 $T_0$를 만족한다는 것은, 임의의 서로 다른 두 점 $x, y\in X$, $x\not = y$에 대해, 그 중 한 점을 포함하고 다른 점을 포함하지 않는 $X$의 열린 부분 집합이 존재하는 것을 말한다. $X$가 {\it Hausdorff} 라는 것은, 임의의 서로 다른 두 점 $x, y\in X$, $x\not = y$에 대해, $x\in U$ 및 $y\in V$를 만족하는 서로소인 열린 부분집합 $U, V$가 존재함을 말한다. \end{definition}

\begin{exercise} \label{exercises-exercise-not-hausdorff} $\Spec(A)$가 일반적으로 {\bf Hausdorff 가 아님} 을 보여라. $\Spec(A)$가 $T_0$ 임을 증명하라. $T_0$가 아닌 위상 공간 $X$의 예를 들어라. \end{exercise}

\begin{remark} \label{exercises-remark-not-hausdorff} 보통, '콤팩트'라는 용어는 준콤팩트이면서 Hausdorff인 공간에 대해 사용한다. \end{remark}

\begin{definition} \label{exercises-definition-irreducible} 위상 공간 $X$가 {\it 기약} 이라는 것은, $X$가 비어있지 않고, 닫힌 부분집합 $Z_1, Z_2\subset X$에 의해 $X = Z_1\cup Z_2$로 쓸 수 있을 때, $Z_1 = X$ 또는 $Z_2 = X$ 중 하나가 성립함을 말한다. 위상 공간의 부분집합 $T\subset X$가 {\it 기약} 이라는 것은, $X$에서 유도된 위상에 관해 기약 위상 공간임을 말한다. 이 정의로부터, $T$가 기약인 필요충분조건은 $X$에서 $T$의 닫힘 $\bar T$가 기약인 것이다. \end{definition}

\begin{exercise} \label{exercises-exercise-irreducible-spec} $\Spec(A)$가 기약인 필요충분조건은 $Nil(A)$가 소 아이디얼인 것이며, 이 경우 그것이 $A$의 유일한 극소 소 아이디얼임을 증명하라. \end{exercise}

\begin{exercise} \label{exercises-exercise-irreducible-prime} 닫힌 부분집합 $T\subset \Spec(A)$가 기약인 필요충분조건은 어떤 소 아이디얼 ${\mathfrak p}\subset A$에 대하여 $T = V({\mathfrak p})$로 쓸 수 있다는 것을 증명하라. \end{exercise}

\begin{definition} \label{exercises-definition-generic-point} 기약 위상공간 $X$의 점 $x$가 $X$의 {\it 일반점} 이란, $X$가 부분집합 $\{x\}$의 닫힘과 같음을 의미한다. \end{definition}

\begin{exercise} \label{exercises-exercise-irreducible-T0-at-most-one-generic} $T_0$ 공간 $X$에서 임의의 기약 닫힌 부분집합이 많아야 하나의 일반점만 가짐을 보여라. \end{exercise}

\begin{exercise} \label{exercises-exercise-spec-sober} % P12-ERR-0423: frozen undefined-X codomain; see ../control/ERRATA_CANDIDATES.jsonl.
$\Spec(A)$ 에서는 임의의 기약 닫힌 부분집합이 일반점을 {\it 실제로} 갖는다는 것을 증명하라. 실제로, 사상 ${\mathfrak p} \mapsto \overline{\{{\mathfrak p}\}}$이 $\Spec(A)$와 $X$의 기약 닫힌 부분집합 전체 사이의 전단사임을 보이시오. \end{exercise}

\begin{exercise} \label{exercises-exercise-irreducible-subset-not-generic} $\Spec(\mathbf{Z})$의 기약 부분집합 중 일반점을 갖지 않는 예를 들어라. \end{exercise}

\begin{definition} \label{exercises-definition-Noetherian-space} 위상공간 $X$가 {\it 뇌터} 라는 것은, $X$의 닫힌 부분집합의 임의 감소열 $Z_1\supset Z_2 \supset Z_3\supset \ldots$이 정지함을 의미한다. (닫힌 부분집합의 임의 증가열이 정지할 때, {\it Artin} 이라고 한다.) \end{definition}

\begin{exercise} \label{exercises-exercise-Noetherian-spec} 환 $A$가 뇌터 이라면, 위상공간 $\Spec(A)$가 뇌터 임을 보이시오. 반대가 거짓임을 보이는 예를 들시오. (원하면 Artin의 경우에도 마찬가지로 하시오.) \end{exercise}

\begin{definition} \label{exercises-definition-irreducible-component} 극대 기약 부분집합 $T\subset X$을, 공간 $X$의 {\it 기약 성분} 라 한다. $X$의 그러한 기약 성분은 자동으로 $X$의 닫힌 부분집합이 된다. \end{definition}

\begin{exercise} \label{exercises-exercise-irreducible-in-irreducible} $X$의 임의 기약 부분집합이 $X$의 어떤 기약 성분에 포함됨을 증명하시오. \end{exercise}

\begin{exercise} \label{exercises-exercise-Noetherian-finite-nr-irreducible} 뇌터 위상 공간 $X$가 유한 개의 기약 성분, 예를 들어 $X_1, \ldots, X_n$ 만을 가지며, $X = X_1\cup X_2\cup\ldots\cup X_n$가 됨을 증명하라. (임의의 $X$는 항상 그 기약 성분의 합집합임에 주의하라. 그러나 예를 들어 $X = {\mathbf R}$에 일반적인 위상을 넣으면, $X$의 기약 성분은 한 점 부분집합이다. 이는 별로 흥미롭지 않다.) \end{exercise}

\begin{exercise} \label{exercises-exercise-irreducible-components-minimal-primes} $\Spec(A)$의 기약 성분이 $A$의 극소 소 아이디얼에 대응함을 보여라. \end{exercise}

\begin{definition} \label{exercises-definition-closed} 점 $x\in X$이 {\it 닫힌 점} 이라는 것은, $\overline{\{x\}} = \{ x\}$ 임을 의미한다. $x, y$를 $X$의 점으로 하자. $x\in \overline{\{y\}}$ 일 때, $x$은 $y$의 {\it 특수화} 이거나, 또는 $y$은 $x$의 {\it 일반화} 라고 한다. \end{definition}

\begin{exercise} \label{exercises-exercise-closed-maximal} $\Spec(A)$의 닫힌 점이 $A$의 극대 아이디얼에 대응함을 보여라. \end{exercise}

\begin{exercise} \label{exercises-exercise-generalization} ${\mathfrak p}$가 $\Spec(A)$ 에서의 ${\mathfrak q}$의 일반화가 되기 위한 필요충분조건은 ${\mathfrak p}\subset {\mathfrak q}$ 임을 보여라. 닫힌 점, 극대 아이디얼, 일반점, 극소 소아이디얼을 일반화와 특수화를 이용하여 특징지어라. % P12-ERR-0424: frozen singular/plural generic-point mismatch; see ../control/ERRATA_CANDIDATES.jsonl.
(여기서는, 기약이 아닐 수도 있는 위상공간 $X$의 점이 $X$의 기약 성분 중 하나의 일반점일 때, 그것을 일반점이라고 부르는 용어를 채택한다.) \end{exercise}

\begin{exercise} \label{exercises-exercise-disjoint-closed-spec} $I$와 $J$을 $A$의 아이디얼로 하자. $V(I)$와 $V(J)$가 서로 소가 되기 위한 조건은 무엇인가. \end{exercise}

\begin{definition} \label{exercises-definition-connected-component} 위상공간 $X$가 {\it 연결} 이라는 것은, 공집합이 아니고 두 개의 공집합이 아닌 서로 소인 열린 부분집합의 합집합이 아님을 의미한다. $X$의 극대 연결 부분집합을 {\it 연결 성분} 이라고 한다. $X$의 임의의 점은 $X$의 어떤 연결 성분에 포함되며, $X$의 임의의 연결 성분은 $X$에서 닫혀 있다. (그러나 일반적으로 연결 성분은 $X$에서 열린 것은 아니다.) \end{definition}

\begin{exercise} \label{exercises-exercise-disconnected-spec} $A$를 0이 아닌 환으로 하자. $\Spec(A)$가 비연결이기 위한 필요충분조건은, 어떤 영이 아닌 환 $B, C$에 대해 $A\cong B \times C$가 되는 것임을 보여라. \end{exercise}

\begin{exercise} \label{exercises-exercise-connected-component-stable-generalization} $T$를 $\Spec(A)$의 연결 성분으로 하자. $T$가 일반화 아래에서 안정적임을 증명하라. $A$가 뇌터라면, $T$가 $\Spec(A)$의 열린 부분집합임을 증명하라. (주의: 예를 들어 $A$가 ${\mathbf F}_2$의 복사본의 무한 곱일 경우, 이는 잘못된 것임. 이 환의 스펙트럼은 무한 개의 닫힌 점으로 이루어져 있음.) \end{exercise}

\begin{exercise} \label{exercises-exercise-primes-kx} $\Spec(k[x])$를 계산하라. 즉, 이 환의 소 아이디얼, 가능한 특수화, 그리고 위상을 기술하라. ($k$가 대수적으로 닫힌 체인 경우뿐만 아니라, $k$가 대수적으로 닫혀 있지 않은 경우도 구하라.) \end{exercise}

\begin{exercise} \label{exercises-exercise-primes-kxy} $k$를 대수적으로 닫힌 체로 하고, $\Spec(k[x, y])$를 계산하라. [힌트: 사상 $\varphi : \Spec(k[x, y]) \to \Spec(k[x])$를 사용하라. $\varphi({\mathfrak p}) = (0)$라면, $S = \{f\in k[x] \mid f \not = 0\}$에 관해서 국소화하고, 국소화와 $\Spec$에 관한 강의 결과를 사용하라.] (왜 대수기하학자들은 이것을 아핀 2-공간이라고 부르는가?) \end{exercise}

\begin{exercise} \label{exercises-exercise-primes-Zy} $\Spec(\mathbf{Z}[y])$를 계산하라. ［힌트: 위와 동일하다.］($\mathbf{Z}$ 위의 아핀 1-공간.) \end{exercise}

\section{국소화} \label{exercises-section-localization}

\begin{exercise} \label{exercises-exercise-submodule-localization} $A$를 환으로, $S \subset A$를 곱셈적 부분집합으로 하자. $M$를 $A$-가군, $N \subset S^{-1}M$를 $S^{-1}A$-부분가군으로 하자. $N = S^{-1}N'$가 되는 $A$-부분가군 $N' \subset M$가 존재함을 보이라. (이 유용한 결과는, 특히 $S^{-1}A$의 아이디얼에 적용할 수 있다.) \end{exercise}

\begin{exercise} \label{exercises-exercise-localize-zero} $A$를 환, $M$를 $A$-가군, $m \in M$로 하자. \begin{enumerate} \item $I = \{a \in A \mid am = 0\}$가 $A$의 아이디얼임을 보이라. \item $A$의 소 아이디얼 $\mathfrak p$에 대하여, $M_\mathfrak p$에서 $m$의 상이 0이 되기 위한 필요충분조건이 $I \not \subset \mathfrak p$임을 보이시오. \item $m$이 영일 필요충분조건은 모든 $A$의 소 아이디얼 $\mathfrak p$에 대하여 $M_\mathfrak p$에서 $m$의 상이 영인 것임을 보여라. \item $m$이 영일 필요충분조건은 모든 $A$의 극대 아이디얼 $\mathfrak m$에 대하여 $M_\mathfrak m$에서 $m$의 상이 영인 것임을 보여라. \item $M = 0$일 필요충분조건은 모든 극대 아이디얼 $\mathfrak m$에 대하여 $M_{\mathfrak m}$이 영인 것임을 보여라. \end{enumerate} \end{exercise}

\begin{exercise} \label{exercises-exercise-localization-is-field} $A$가 서로 다른 소 아이디얼을 세 개 이상 가진 정역이며, $f \in A$가 주 국소화 $A_f = \{1, f, f^2, \ldots \}^{-1}A$를 체로 만드는 원소인 조 $(A, f)$를 찾아라. \end{exercise}

\begin{exercise} \label{exercises-exercise-localize-finite-module-zero} $A$을 환이라 하고, $M$을 유한 $A$-가군, $S \subset A$를 곱셈 집합으로 하자. $S^{-1}M = 0$라고 가정하자. 주 국소화 $M_f = \{1, f, f^2, \ldots \}^{-1}M$가 0이 되는 $f \in S$가 존재함을 보여라. \end{exercise}

\begin{exercise} \label{exercises-exercise-localization-is-quotient} $A$가 환이고, $0 \not = I \not = A$가 0이 아닌 진 아이디얼이며, $S \subset A$가 곱셈적 부분집합이고, $A$-대수로서 $A/I \cong S^{-1}A$이 되는 세 쌍 $(A, I, S)$의 예를 들어라. \end{exercise}

\section{Nakayama의 보조정리} \label{exercises-section-nakayama}

\begin{exercise} \label{exercises-exercise-nakayama} $A$를 환으로 한다. $I$를 $A$의 아이디얼로 한다. $M$를 $A$-가군으로 한다. $x_1, \ldots, x_n \in M$로 한다. 다음을 가정하자: \begin{enumerate} \item $M/IM$은 $x_1, \ldots, x_n$에서 생성되며, \item $M$은 유한 $A$-가군이다, \item $I$은 $A$의 모든 극대 아이디얼에 포함된다. \end{enumerate} $x_1, \ldots, x_n$가 $M$를 생성함을 보이라. (해법 제안: 연습 \ref{exercises-exercise-localize-zero}과 국소화의 완전성을 이용하여, $A$의 극대 아이디얼에서의 국소화로 귀결시켜라. 다음으로, $M$를 $x_1, \ldots, x_n$로 생성되는 부분가군으로 나눈 몫을 고려하고, 강의에서의 Nakayama 보조정리의 주장으로 귀결시켜라.) \end{exercise}



\section{길이} \label{exercises-section-length}

\begin{definition} \label{exercises-definition-length} % P12-ERR-0425: frozen undefined-R module designation; see ../control/ERRATA_CANDIDATES.jsonl.
$A$를 환, $M$를 $A$-가군으로 하자. $M$의 $R$-가군으로서의 {\it 길이} 란, $$
\text{length}_A(M)
=
\sup
\{
n
\mid
\exists\ 0 = M_0 \subset M_1 \subset \ldots \subset M_n = M,
\text{ }M_i \not = M_{i + 1}
\}.
$$를 의미한다. 다시 말하면, 부분가군 사슬의 길이의 상한이다. \end{definition}

\begin{exercise} \label{exercises-exercise-length-is-one} 환 $A$ 위의 가군 $M$의 길이가 $1$인 필요충분조건은, $A$의 어떤 극대 아이디얼 $\mathfrak m$에 대해 $A/\mathfrak m$와 동형임을 보이는 것이다. \end{exercise}

\begin{exercise} \label{exercises-exercise-length-easy} 다음 각 환 위의 각 가군의 길이를 계산하라. 답은 간결하게 (!) 설명하라. (짧은 완전열에서의 길이 함수의 덧셈성을 자유롭게 사용해도 좋다. 대수, 보조정리 \ref{algebra-lemma-length-additive} 참조.) \begin{enumerate} \item $\mathbf{Z}$ 위의 $\mathbf{Z}/120\mathbf{Z}$의 길이. \item $\mathbf{C}[x]$ 위의 $\mathbf{C}[x]/(x^{100} + x + 1)$의 길이. \item $\mathbf{R}[x]$ 위의 $\mathbf{R}[x]/(x^4 + 2x^2 + 1)$의 길이. \end{enumerate} \end{exercise}

\begin{exercise} \label{exercises-exercise-compute-length} $A = k[x, y]_{(x, y)}$를 원점에서의 아핀 평면의 국소환으로 하자. 체 $k$에 대해서는 원하는 가정을 둘 수 있다. $f = x^3 + x^2y^2 + y^{100}$ 및 $g = y^3 - x^{999}$로 한다. $A$-가군으로서 $A/(f, g)$의 길이는 얼마인가. (진행 방법의 한 예: 다양한 $n$에 대해, $A/{\mathfrak m}_A^n=
k[x, y]/(x, y)^n$ 형태의 몫에서 $f$와 $g$가 생성하는 아이디얼을 생각하라. % P12-ERR-0426: frozen ungrouped quotient-plus-power formula; see ../control/ERRATA_CANDIDATES.jsonl.
$A/(f, g)+{\mathfrak m}_A^n \cong A/(f, g)+{\mathfrak m}_A^{n + 1}$가 되는 $n$를 찾아 NAK를 사용하라.) \end{exercise}


\section{수반 소 아이디얼} \label{exercises-section-ass}

\noindent 수반 소 아이디얼은 대수, 절 \ref{algebra-section-ass}에서 논의한다.

\begin{exercise} \label{exercises-exercise-compute-ass} 다음 각 가군에 대해, 수반 소 아이디얼의 집합을 계산하라. \begin{enumerate} \item $R = k[x, y]$ 및 $M = R/(xy(x + y))$, \item $R = \mathbf{Z}[x]$ 및 $M = R/(300x + 75)$, \item $R = k[x, y, z]$ 및 $M = R/(x^3, x^2y, xz)$. \end{enumerate} 여기서도 통상대로 $k$는 체로 한다. \end{exercise}

\begin{exercise} \label{exercises-exercise-prime-power-not-primary} 뇌터 환 $R$와 소 아이디얼 $\mathfrak p$ 중, $\mathfrak p$가 $R/\mathfrak p^2$의 유일한 수반 소 아이디얼이 아닌 것을 예로 들라. \end{exercise}

\begin{exercise} \label{exercises-exercise-product-primes-not-only-associated} $R$을, 포함 관계가 없는 소 아이디얼 $\mathfrak p$, $\mathfrak q$를 가진 뇌터 환으로 하자. 즉, $\mathfrak p \not \subset \mathfrak q$ 및 $\mathfrak q \not \subset \mathfrak p$로 하자. \begin{enumerate} \item $N = R/(\mathfrak p \cap \mathfrak q)$에 대해 $\text{Ass}(N) = \{\mathfrak p, \mathfrak q\}$가 됨을 보이시오. \item 가군 $M = R/\mathfrak p \mathfrak q$가 $\mathfrak p$와도 $\mathfrak q$와도 다른 수반 소 아이디얼을 가질 수 있음을, 예를 들어 보이시오. \end{enumerate} \end{exercise}


\section{Ext 군} \label{exercises-section-ext}

\noindent Ext 군은 대수, 절 \ref{algebra-section-ext}에서 정의한다.

\begin{exercise} \label{exercises-exercise-compute-ext-abelian-groups} 주어진 가군의 모든 Ext 군 $\Ext^i(M, N)$을 $\mathbf{Z}$-가군의 범주(Abel 군의 범주라고도 불림)에서 계산하라. \begin{enumerate} \item $M = \mathbf{Z}$ 그리고 $N = \mathbf{Z}$, \item $M = \mathbf{Z}/4\mathbf{Z}$ 그리고 $N = \mathbf{Z}/8\mathbf{Z}$, \item $M = \mathbf{Q}$ 그리고 $N = \mathbf{Z}/2\mathbf{Z}$, \item $M = \mathbf{Z}/2\mathbf{Z}$ 그리고 $N = \mathbf{Q}/\mathbf{Z}$. \end{enumerate} \end{exercise}

\begin{exercise} \label{exercises-exercise-compute-in-regular} $R = k[x, y]$으로 하고, $k$를 체로 한다. \begin{enumerate} \item Koszul 복합체 $$
0 \to R
\xrightarrow{
\left(
\begin{matrix}
y \\
-x
\end{matrix}
\right)
}
R^{\oplus 2} \xrightarrow{(x, y)} R \xrightarrow{f \mapsto f(0, 0)} k \to 0
$$가 완전함을 손으로 계산해 보이라. \item $k = R/(x, y)$를 $R$-가군으로 생각하고, $\Ext^i_R(k, k)$를 계산하라. \end{enumerate} \end{exercise}

\begin{exercise} \label{exercises-exercise-infinitely-many-nonzero-ext} 모든 $i \geq 0$에 대해 $\Ext^i_R(M, N)$가 0이 아닌, 뇌터 환 $R$과 유한 가군 $M$, $N$의 예를 들어라. \end{exercise}

\begin{exercise} \label{exercises-exercise-infinite-ext} $\Ext^1_R(R/I, R/I)$가 유한 $R$-가군이 아닌 그런 환 $R$과 아이디얼 $I$의 예를 들어라. ($R$가 뇌터라면, 대수, 보조정리 \ref{algebra-lemma-ext-noetherian}에 의해, 이것은 일어날 수 없는 것으로 알려져 있다.) \end{exercise}


\section{깊이} \label{exercises-section-depth}

\noindent 깊이는 대수학, 절 \ref{algebra-section-depth}에서 정의하고, 쌍대화 복합체, 절 \ref{dualizing-section-depth}에서 더욱 연구한다.

\begin{exercise} \label{exercises-exercise-compute-depth} $R$를 환, $I \subset R$를 아이디얼, $M$를 $R$-가군으로 한다. 다음 각 경우에 $\text{depth}_I(M)$를 계산하라. \begin{enumerate} \item $R = \mathbf{Z}$, $I = (30)$, $M = \mathbf{Z}$, \item $R = \mathbf{Z}$, $I = (30)$, $M = \mathbf{Z}/(300)$, \item $R = \mathbf{Z}$, $I = (30)$, $M = \mathbf{Z}/(7)$, \item $R = k[x, y, z]/(x^2 + y^2 + z^2)$, $I = (x, y, z)$, $M = R$, \item $R = k[x, y, z, w]/(xz, xw, yz, yw)$, $I = (x, y, z, w)$, $M = R$. \end{enumerate} 여기서 $k$는 체라 한다. 마지막 두 경우에는, 극대 아이디얼 $I$에서 자유롭게 국소화해도 좋다. \end{exercise}

\begin{exercise} \label{exercises-exercise-depth-not-inherited-localization} 깊이 $\geq 1$의 뇌터 국소환 $(R, \mathfrak m, \kappa)$과, 다음을 만족하는 소 아이디얼 $\mathfrak p$의 예를 들라: \begin{enumerate} \item $\text{depth}_\mathfrak m(R) \geq 1$, \item $\text{depth}_\mathfrak p(R_\mathfrak p) = 0$, \item $\dim(R_\mathfrak p) \geq 1$. \end{enumerate} (3)을 요구하지 않으면, 이 연습문제는 너무 쉽다. 왜인가? \end{exercise}

\begin{exercise} \label{exercises-exercise-depth-torsion-free} $(R, \mathfrak m)$를 뇌터 국소 정역으로 하고, $M$를 유한 $R$-가군으로 하자. \begin{enumerate} \item $M$가 비틀림이 없다면, $M$의 $R$ 위의 깊이가 적어도 $1$ 임을 보이시오. \item 깊이가 $1$인 예를 들어라. \end{enumerate} \end{exercise}

\begin{exercise} \label{exercises-exercise-depth-examples} 임의의 $m \geq n \geq 0$에 대해, $\dim(R) = m$이고 $\text{depth}(R) = n$가 되는 뇌터 국소환 $R$의 예를 들어라. \end{exercise}

\begin{exercise} \label{exercises-exercise-make-depth-1} $(R, \mathfrak m)$를 뇌터 국소환으로 하고, $M$를 유한 $R$-가군으로 한다. 다음 성질을 가진 표준적인 짧은 완전열 $$
0 \to K \to M \to Q \to 0
$$가 존재함을 보여라: \begin{enumerate} \item $\text{depth}(Q) \geq 1$, \item $K$는 0이거나, $\text{Supp}(K) = \{\mathfrak m\}$, \item $\text{length}_R(K) < \infty$. \end{enumerate} 힌트: 뇌터 성질을 사용하여 (2)와 같은 극대 부분 가군 $K$가 존재함을 보여주고, 그 다음 $Q = M/K$가 (1)을 만족하고, $K$가 (3)을 만족함을 보여라. \end{exercise}

\begin{exercise} \label{exercises-exercise-make-depth-2} $(R, \mathfrak m)$를 뇌터 국소환, $M$를 깊이 $\geq 2$의 유한 $R$-가군, $N \subset M$를 0이 아닌 부분가군으로 두자. \begin{enumerate} \item $\text{depth}(N) \geq 1$ 임을 보이라. \item $\text{depth}(N) = 1$ 임을 위한 필요충분조건은, 몫가군 $M/N$가 $\text{depth}(M/N) = 0$를 만족함임을 보이라. \item $N \subset N'$이 유한 여길이, 즉 $\text{length}_R(N'/N) < \infty$이며, 또한 $N'$의 깊이가 $\geq 2$가 되는 부분 가군 $N' \subset M$이 존재함을 보여라. 힌트: 연습 \ref{exercises-exercise-make-depth-1}를 $M/N$에 적용하고, $N'$를 $K$의 역상으로 선택하라. \end{enumerate} \end{exercise}

\begin{exercise} \label{exercises-exercise-Hartshorne-reduced} $(R, \mathfrak m)$를 뇌터 국소환으로 하자. $R$가 축소환이라고 가정하자. 즉, $R$는 영이 아닌 멱영원을 가지지 않는다고 가정한다. 또한, $R$는 서로 다른 두 개의 극소 소 아이디얼 $\mathfrak p$와 $\mathfrak q$를 가진다고 가정한다. \begin{enumerate} \item % P12-ERR-0427: frozen ungrouped quotient-ideal expressions; see ../control/ERRATA_CANDIDATES.jsonl.
$R$-가군의 열 $$
0 \to R \to R/\mathfrak p \oplus R/\mathfrak q \to
R/\mathfrak p + \mathfrak q \to 0
$$가 완전함을 보이자(모든 경우에 확인하라). 사상은 $x \mapsto (x \bmod \mathfrak p, x \bmod \mathfrak q)$ 및 $(y \bmod \mathfrak p, z \bmod \mathfrak q) \mapsto
(y - z \bmod \mathfrak p + \mathfrak q)$이다. \item $\text{depth}(R) \geq 2$라면, $\dim(R/\mathfrak p + \mathfrak q) \geq 1$임을 보이자. \item $\text{depth}(R) \geq 2$라면, $U = \Spec(R) \setminus \{\mathfrak m\}$가 연결 위상 공간임을 보이라. \end{enumerate} 이는 Hartshorne의 연결성 정리의 매우 특별한 경우를 증명하고 있다. 같은 정리는 $\text{depth} \geq 2$의 뇌터 국소환의 뚫린 스펙트럼 $U$가 연결되어 있음을 말한다. \end{exercise}

\begin{exercise} \label{exercises-exercise-depth-2} $(R, \mathfrak m)$를 뇌터 국소환으로 하고, $x, y \in \mathfrak m$를 길이 $2$의 정칙열로 하자. 임의의 $n \geq 2$에 대해, $$
x^{n - 1}y^{n - 1} = a x^n + b y^n
$$를 만족하는 $a, b \in R$가 존재하지 않음을 보이라. 제안: 먼저 $n = 2$의 경우를 시도하고, 논의하는 방법을 찾아라. 주의: 이 결과에는 단항식 추측이라고 불리는 대폭적인 일반화가 있다. \end{exercise}


\section{Cohen--Macaulay 가군과 환} \label{exercises-section-CM}

\noindent Cohen--Macaulay 가군은 대수, 장 \ref{algebra-section-CM}에서 연구하고, Cohen--Macaulay 환은 대수, 장 \ref{algebra-section-CM-ring}에서 연구한다.

\begin{exercise} \label{exercises-exercise-examples-CM} 다음 각 경우에 대해, 예 또는 아니오로 답하시오. 설명이나 증명은 필요 없다. \begin{enumerate} \item $p$를 소수라 하자. 국소환 $\mathbf{Z}_{(p)}$은 Cohen--Macaulay 국소환인가? \item $p$을 소수라 하자. 국소환 $\mathbf{Z}_{(p)}$은 정칙 국소환인가? \item $k$를 체라 하자. 국소환 $k[x]_{(x)}$은 Cohen--Macaulay 국소환인가? \item $k$을 체라 하자. 국소환 $k[x]_{(x)}$은 정칙 국소환인가? \item $k$을 체라 하자. 국소환 $(k[x, y]/(y^2 - x^3))_{(x, y)} =
k[x, y]_{(x, y)}/(y^2 - x^3)$은 Cohen--Macaulay 국소환인가? \item $k$을 체라 하자. 국소환 $(k[x, y]/(y^2, xy))_{(x, y)} =
k[x, y]_{(x, y)}/(y^2, xy)$은 Cohen--Macaulay 국소환인가? \end{enumerate} \end{exercise}


\section{특이점} \label{exercises-section-singularities}

\begin{exercise} \label{exercises-exercise-singularities} $k$를 임의의 체로 한다. $A = k[[x, y]]/(f)$ 및 $B = k[[u, v]]/(g)$ 라 하고, 여기서 $f = xy$, $g = uv + \delta$, $\delta \in (u, v)^3$ 라 하자. $A$와 $B$가 동형인 환임을 보여라. \end{exercise}

\begin{remark} \label{exercises-remark-singularities} 체 $k$ 위의 곡선의 특이점이 보통 이중점이라는 것은, 그 점에서의 곡선의 완비 국소환이 $k'[[x, y]]/(f)$의 형태이고, 다음을 만족함을 의미한다: (a) $k'$는 $k$의 유한 차 분리 확장이다; (b) $f$의 초기항은 차수 2, 즉 모두 영인 것은 아닌 $a, b, c\in k'$에 대하여 $q = ax^2 + bxy + cy^2$의 형태를 가진다; (c) $q$는 $k'$ 위의 비퇴화 2차 형식이다 (특성 2에서는, 이는 $b$가 0이 아님을 의미한다). 일반적으로, 조합 $(k', q)$의 각 동형류마다, 그러한 환의 동형류가 하나 존재한다. \end{remark}

\begin{exercise} \label{exercises-exercise-periodic-resolution} $R$를 환, $n \geq 1$ 라 하자. $A$, $B$를, 계수가 $R$에 속하는 $n \times n$ 행렬이라 하고, $R$의 어떤 영인자가 아닌 원소 $f$에 대해 $AB = f 1_{n \times n}$를 만족한다고 하자. $S = R/(f)$라 하자. $$
\ldots \to
S^{\oplus n} \xrightarrow{B}
S^{\oplus n} \xrightarrow{A}
S^{\oplus n} \xrightarrow{B}
S^{\oplus n} \to \ldots
$$가 완전함을 보여라. \end{exercise}


\section{구성가능집합} \label{exercises-section-constructible}

\noindent $k$를 대수폐쇄체, 예를 들어 복소수체 $\mathbf{C}$로 하자. $n \geq 0$로 둔다. 다항식 $f \in k[x_1, \ldots, x_n]$는 값을 대입함으로써 함수 $f : k^n \to k$를 정한다. 부분집합 $Z \subset k^n$가 {\it 대수집합} 이라는 것은, 다항식족의 공통 영점집합임을 의미한다.

\begin{exercise} \label{exercises-exercise-finite-nr-equations} 대수집합이 항상 유한 개의 다항식의 영점집합으로 쓸 수 있음을 증명하라. \end{exercise}

\noindent 위의 표기법 아래에서, 부분집합 $E \subset k^n$가 {\it 구성가능} 이라는 것은, 다항식 $f$에 의한 $Z \cap \{f \not = 0\}$의 형태의 집합의 유한합집합임을 의미한다.

\begin{exercise} \label{exercises-exercise-constructible-classical} 다음을 보이라: \begin{enumerate} \item 구성가능집합의 여집합은 구성가능집합이다, \item 구성가능집합의 유한합집합은 구성가능집합이다, \item 구성가능집합의 유한교집합은 구성가능집합이다, \item 임의의 구성가능집합 $E$는 유한한 서로소 합집합 $E = \coprod E_i$로 쓸 수 있다. 여기서 각 $E_i$는 $Z \cap \{f \not = 0\}$의 형태이며, $Z$는 대수집합, $f$는 다항식이다. \end{enumerate} \end{exercise}

\begin{exercise} \label{exercises-exercise-division-with-remainder} $R$를 환으로 하고, $f = a_d x^d + a_{d - 1} x^{d - 1} + \ldots + a_0 \in R[x]$로 한다. (보통대로, 이 표기는 $a_0, \ldots, a_d \in R$를 의미한다.) $g \in R[x]$로 한다. $$
a_d^N g = q f + r
$$이고 $\deg(r) < d$, 즉 어떤 $c_i \in R$에 대해 $r = c_0 + c_1 x + \ldots + c_{d - 1}x^{d - 1}$가 되는 $N \geq 0$ 및 $r, q \in R[x]$를 찾을 수 있음을 증명하라. \end{exercise}


\section{힐베르트의 영점 정리} \label{exercises-section-Hilbert-Nullstellensatz}

\begin{exercise} \label{exercises-exercise-uncountable} {\it 복소수를 사용한 사소한 논의!} ${\mathbf C}$를 복소수체로 하고, $V$를 ${\mathbf C}$ 위의 벡터 공간으로 하자. 선형 작용소 $T : V \to V$의 스펙트럼이란, 작용소 $T - \lambda \text{id}_V$가 가역이 아닌 복소수 $\lambda \in {\mathbf C}$의 집합이다. \begin{enumerate} \item $\mathbf{C}(X)$의 ${\mathbf C}$ 위의 차원이 비가산임을 보이라. \item $V$의 차원이 유한하거나 가산이면, $V$ 위의 임의의 선형 작용소가 공집합이 아닌 스펙트럼을 가짐을 보이라. \item 유한 생성 ${\mathbf C}$-대수 $R$가 체라면, 사상 ${\mathbf C}\to R$이 동형임을 보이라. \item ${\mathbf C}[x_1, \ldots, x_n]$의 임의의 극대 아이디얼 ${\mathfrak m}$은, 어떤 $\alpha_i \in {\mathbf C}$에 의해 $(x_1-\alpha_1, \ldots, x_n-\alpha_n)$의 형태로 쓸 수 있음을 보이라. \end{enumerate} \end{exercise}

\begin{remark} \label{exercises-remark-HNSS} $k$를 체로 하자. 이때, 임의의 정수 $n\in {\mathbf N}$와 임의의 극대 아이디얼 ${\mathfrak m} \subset k[x_1, \ldots, x_n]$에 대해, 몫 $k[x_1, \ldots, x_n]/{\mathfrak m}$는 $k$의 유한 차수 체 확장이다. 이는 강의 후반에 보일 것이다. 물론, (확인해 보라), 여기서 유한 생성 $k$-대수의 극대 아이디얼에 대해서도 유사한 주장이 도출된다. 위의 연습은 $k = {\mathbf C}$의 경우에 이를 증명한다. \end{remark}

\begin{exercise} \label{exercises-exercise-Hilbert-Nullstellensatz} $k$를 체라 하자. 참고: \ref{exercises-remark-HNSS}를 사용하라. \begin{enumerate} \item $R$를 $k$-대수라고 하자. $\dim_k R < \infty$이고 $R$가 정역이라고 가정하자. $R$가 체임을 보여라. \item $R$가 유한 생성 $k$-대수이고 $f\in R$가 멱영원이 아니라고 하자. 극대 아이디얼 ${\mathfrak m} \subset R$가 존재하여 $f\not\in {\mathfrak m}$가 됨을 보여라. \item $R$가 체 위의 유한형 대수가 아닐 때 이 명제가 거짓임을 예로 보여라. \item 모든 근기 아이디얼 $I \subset {\mathbf C}[x_1, \ldots, x_n]$가 이를 포함하는 극대 아이디얼들의 교집합임을 보여라. \end{enumerate} \end{exercise}

\begin{remark} \label{exercises-remark-Hilbert-Nullstellensatz} 이것이 Hilbert 영점정리이다. 즉, $\Spec(k[x_1, \ldots, x_n])$의 닫힌 부분집합(앞서 언급한 근기 아이디얼에 해당)은 그 안에 포함된 닫힌 점들에 의해 결정된다. \end{remark}

\begin{exercise} \label{exercises-exercise-product-matrices-ring} $A =
{\mathbf C}[x_{11}, x_{12}, x_{21}, x_{22}, y_{11}, y_{12}, y_{21}, y_{22}]$ $I$을 $A$의 아이디얼로, 행렬 $XY$의 성분으로 생성되는 것으로 하자. 여기서 $$
X = \left(
\begin{matrix}
x_{11} & x_{12}\\
x_{21} & x_{22}
\end{matrix}
\right)
\quad\text{그리고}\quad
Y = \left(
\begin{matrix}
y_{11} & y_{12}\\
y_{21} & y_{22}
\end{matrix}
\right).
$$ $\Spec(A)$의 닫힌 부분집합 $V(I)$의 기약 성분을 찾아라. (그것들을 서술하고, 각각을 정하는 방정식을 제시하라는 의미이다. 쓴 방정식이 소아이디얼을 정함을 증명할 필요는 없다.) 힌트: \begin{enumerate} \item Hilbert의 영점 정리를 사용해도 된다. $V(I)$의 닫힌 점 전체를 덮는 기약 국소 닫힌 부분집합을 찾으면 충분하다. \item 쉬운 성분이 두 개 있다. \item 연속사상에 의한 기약 집합의 상은 약분이다. \end{enumerate} \end{exercise}

\section{차원} \label{exercises-section-dimension}

\begin{exercise} \label{exercises-exercise-dimension-bigger-one-finite-nr-primes} 소 아이디얼이 유한 개이고, 차원이 $> 1$인 환 $A$를 구성하라. \end{exercise}

\begin{exercise} \label{exercises-exercise-hypersurface-in-A2-dimension-one} $f \in \mathbf{C}[x, y]$를 상수가 아닌 다항식이라고 하자. $\mathbf{C}[x, y]/(f)$의 차원이 $1$임을 보여라. \end{exercise}

\begin{exercise} \label{exercises-exercise-dimension-polynomial-ring} $(R, \mathfrak m)$를 뇌터 국소환으로 하고, $n \geq 1$로 하자. 다항식환 $R[x_1, \ldots, x_n]$에서 $\mathfrak m' = (\mathfrak m, x_1, \ldots, x_n)$로 두자. 다음을 보여라: $$
\dim(R[x_1, \ldots, x_n]_{\mathfrak m'}) = \dim(R) + n.
$$ \end{exercise}


\section{사슬형 환} \label{exercises-section-catenary}

\begin{definition} \label{exercises-definition-catenary} 뇌터 환 $A$가 {\it 사슬형}이라는 것은, 임의의 세 개의 소 아이디얼 ${\mathfrak p}_1 \subset {\mathfrak p}_2 \subset {\mathfrak p}_3$에 대해 $$
ht({\mathfrak p}_3 / {\mathfrak p}_1) = ht({\mathfrak p}_3/{\mathfrak p}_2) +
ht({\mathfrak p}_2/{\mathfrak p}_1).
$$가 성립함을 의미한다. 여기서 $ht(\mathfrak p/\mathfrak q)$는, 환 $A/\mathfrak q$에서 $\mathfrak p/\mathfrak q$의 높이를 뜻한다. 식으로 쓰면, $$
ht(\mathfrak p/\mathfrak q) =
\dim(A_\mathfrak p/\mathfrak qA_\mathfrak p) =
\dim((A/\mathfrak q)_\mathfrak p) =
\dim((A/\mathfrak q)_{\mathfrak p/\mathfrak q})
$$ 위상 공간 $X$가 {\it 사슬형}이라는 것은, $T$과 $T'$가 닫혀 있고 기약인 $T \subset T' \subset X$가 주어졌을 때, 기약 닫힌 부분 집합의 극대 사슬 $$
T = T_0 \subset T_1 \subset \ldots \subset T_n = T'
$$가 존재하며, 그러한 모든 사슬이 동일한 (유한한) 길이를 가진다는 것을 말한다. \end{definition}

\begin{exercise} \label{exercises-exercise-catenary-the-same} 대수, 정의 \ref{algebra-definition-catenary} 로 정의한 사슬형 개념이, 뇌터 환에 대해서는 정의 \ref{exercises-definition-catenary} 의 개념과 일치함을 보이라. \end{exercise}

\begin{exercise} \label{exercises-exercise-Noetherian-local-domain-dim-2-catenary} 차원 $2$의 뇌터 국소 정역이 사슬형임을 보여라. \end{exercise}

\begin{exercise} \label{exercises-exercise-finite-type-over-field-catenary} $k$를 체로 한다. 유한형 $k$-대수가 사슬형임을 보여라. \end{exercise}

\begin{exercise} \label{exercises-exercise-example-no-dim-function} 차원 함수 $\delta : X \to \mathbf{Z}$를 가지지 않는, 유한하고, sober하며, 사슬형인 위상 공간 $X$의 예를 들어라. 여기서 $\delta : X \to \mathbf{Z}$가 차원 함수라는 것은, $x, y \in X$에 대해 다음이 성립함을 의미한다: \begin{enumerate} \item $x \leadsto y$ 그리고 $x \not = y$ 라면 $\delta(x) > \delta(y)$, \item $x \leadsto y$ 그리고 $\delta(x) \geq \delta(y) + 2$ 라면, $x \leadsto z \leadsto y$ 및 $\delta(x) > \delta(z) > \delta(y)$를 만족하는 $z \in X$가 존재한다. \end{enumerate} 공간을 명확히 기술하고, 차원 함수가 존재할 수 없는 이유를 간단히 설명하라. \end{exercise}



\section{분수체} \label{exercises-section-fraction-fields}

\begin{exercise} \label{exercises-exercise-find-fraction-field} 정역 $$
{\mathbf Q}[r, s, t]/(s^2-(r-1)(r-2)(r-3), t^2-(r + 1)(r + 2)(r + 3)).
$$를 생각한다. 이것과 동형인 분수체를 가지는 ${\mathbf Q}[x, y]/(f)$ 형태의 정역을 찾아라. \end{exercise}



\section{초월 차수} \label{exercises-section-transcendence}

\begin{exercise} \label{exercises-exercise-algebraic-extension} $K'/K/k$를 체 확장으로 하고, $K'$는 $K$ 위에서 대수적이라고 하자. $\text{trdeg}_k(K) = \text{trdeg}_k(K')$를 증명하라. (힌트: $x_1, \ldots, x_d \in K$가 $k$ 위에서 대수적 독립이며, $d < \text{trdeg}_k(K')$ 라면, $k(x_1, \ldots, x_d) \subset K$는 대수 확장이 될 수 없음을 보여라.) \end{exercise}

\begin{exercise} \label{exercises-exercise-growth-powers-subvector-space} $k$를 체로 하고, $K/k$를 초월 차수 $d$의 유한 생성 확장으로 하자. $V, W \subset K$가 유한 차원의 $k$-부분 벡터 공간일 때, $$
VW = \{f \in K \mid f = \sum\nolimits_{i = 1, \ldots, n} v_i w_i
\text{인 어떤 }n\text{과 }v_i \in V, w_i \in W\}
$$로 쓴다. 이는 유한 차원의 $k$-부분 벡터 공간이다. $V^2 = VV$, $V^3 = V V^2$ 등으로 둔다. \begin{enumerate} \item 모든 $n \geq 1$에 대해 $\dim V^n \geq \epsilon n^d$가 되도록 하는 $V \subset K$ 및 $\epsilon > 0$를 찾을 수 있음을 보이라. \item 반대로, 임의의 유한 차원 $V \subset K$에 대해, 모든 $n \geq 1$에 대해 $\dim V^n \leq C n^d$가 되도록 하는 $C > 0$가 존재함을 보이라. (진행 방법의 한 예: 먼저 $k[x_1, \ldots, x_d]$의 부분 벡터 공간에 대해 보이라. 다음으로 $k(x_1, \ldots, x_d)$의 부분 벡터 공간에 대해 보이라. 마지막으로, $K/k(x_1, \ldots, x_d)$가 유한 확장인 경우, $K$의 $k(x_1, \ldots, x_d)$ 위의 기저를 선택하고, 이 기저에 의한 전개를 사용하여 논의하라.) \item $K$의 유한 차원 부분 벡터 공간의 거듭제곱 증가를 통해, 초월 차수를 재정의할 수 있음을 유도하라. \end{enumerate} 이것은 $k$ 위의 (비가환) 대수의 Gelfand--Kirillov 차원과 관계가 있다. \end{exercise}


\section{올의 차원} \label{exercises-section-dimension-fibres}

\noindent 차원 공식에 관한 문제를 몇 가지 제시한다. 대수, 절 \ref{algebra-section-dimension-formula}를 참조하라.

\begin{exercise} \label{exercises-exercise-nr-components-fibre} $k$를 임의의 대수적으로 닫힌 체로 하자. 이하, $k[x]$와 $k[x, y]$는 다항식환을 나타낸다. \begin{enumerate} \item 임의의 정수 $n \geq 0$에 대해, $A/xA$의 스펙트럼이 정확히 $n$ 개의 기약 성분을 갖도록 하는, 정역의 유한형 확장 $k[x] \subset A$를 찾아라. \item 모든 $\alpha \in k$에 대해 $A/(x - \alpha)A$의 스펙트럼이 공집합이 아니고 기약이 되도록 하는, 정역의 유한형 확장 $k[x] \subset A$의 예를 만들어라. \item 모든 $(\alpha, \beta) \in k^2 \setminus \{(0, 0)\}$에 대해 $A/(x - \alpha, y - \beta)A$의 스펙트럼이 기약\footnote{가법이라면 비어있지 않음을 기억하라. }이며, $A/(x, y)A$의 스펙트럼이 공집합이 아니고 기약이 아니게이 되도록 하는 정역의 유한형 확장 $k[x, y] \subset A$의 예를 만들어라. \end{enumerate} \end{exercise}

\begin{exercise} \label{exercises-exercise-codim-1} $k$를 원하는 대수폐쇄체로 한다. $n \geq 1$로 한다. $k[x_1, \ldots, x_n]$를 다항식환으로 하고, $\mathfrak m = (x_1, \ldots, x_n)$로 둔다. $k[x_1, \ldots, x_n] \subset A$를 정역의 유한형 확장으로 한다. $d = \dim(A)$로 둔다. \begin{enumerate} \item $A/\mathfrak mA \not = 0$이면 $d - 1 \geq \dim(A/\mathfrak m A) \geq d - n$ 임을 보여라. \item 각 값이 실제로 나타날 수 있음을 예를 들어 보여라. \item $\Spec(A/\mathfrak m A)$이 서로 다른 차원의 기약 성분을 가질 수 있음을 예를 들어 보여라. \end{enumerate} \end{exercise}


\section{유한 국소 자유 가군} \label{exercises-section-finite-locally-free}

\begin{definition} \label{exercises-definition-finite-locally-free} $A$를 환으로 한다. {\it 유한 국소 자유} $A$-가군 $M$이란, 임의의 ${\mathfrak p} \in \Spec(A)$에 대하여, $M_f$가 유한 자유 $A_f$-가군이 되는 $f\in A$,$f \not \in {\mathfrak p}$이 존재하는 그러한 가군임을 기억하자. $M$이 계수 $1$의 유한 국소 자유 가군일 때, $M$를 {\it 가역 가군}이라고 한다. 즉, 임의의 ${\mathfrak p} \in \Spec(A)$에 대해, $A_f$-가군으로서 $M_f \cong A_f$가 되는 $f\in A$, $f \not \in \mathfrak p$가 존재한다. \end{definition}

\begin{exercise} \label{exercises-exercise-tensor-finite-locally-free} 유한 국소 자유 가군의 텐서곱이 유한 국소 자유임을 증명하라. 두 가역 가군의 텐서곱이 가역임을 증명하라. \end{exercise}

\begin{definition} \label{exercises-definition-class-group} $A$를 환으로 한다. {\it $A$의 유군}, 때로는 {\it $A$의 Picard 군}이라고도 불리는 것은, 가역 $A$-가군의 동형류 전체의 집합 $\Pic(A)$에 텐서곱으로 정의되는 군 연산을 가진 것이다 (연습 \ref{exercises-exercise-tensor-finite-locally-free} 참조). \end{definition}

\noindent $A$의 유군이 자명하기 위한 필요충분조건은 모든 가역 가군이 계수 1짜리 자유 가군과 동형이라는 것임에 주의하라.

\begin{exercise} \label{exercises-exercise-class-group-trivial} 다음 각 환의 군이 자명함을 보여라: \begin{enumerate} \item 체 $k$ 위의 다항식환 $A = k[x]$, \item 정수환 $A = \mathbf{Z}$, \item 체 $k$ 위의 다항식환 $A = k[x, y]$, \item 체 $k$ 위의 몫 $k[x, y]/(xy)$. \end{enumerate} \end{exercise}

\begin{exercise} \label{exercises-exercise-class-group-not-trivial} $k$를 표수가 $2$가 아닌 체로 하고, $t_1, \ldots, t_n \in k$는 서로 다르며, $n \geq 3$는 홀수이고, $f(x) = (x - t_1) \ldots (x - t_n)$라고 하자. 환 $A = k[x, y]/(y^2 - f(x))$의 유군이 자명하지 않음을 보이라. (힌트: 아이디얼 $(y, x - t_1)$이 $\Pic(A)$의 비자명한 원을 정함을 보이라.) \end{exercise}

\begin{exercise} \label{exercises-exercise-trace-det} $A$을 환으로 한다. \begin{enumerate} \item $M$를 유한 국소 자유 $A$-가군으로 하고, $\varphi : M \to M$를 자기준동형으로 한다. $\varphi$의 {\it 트레이스}와 {\it 행렬식}을 정의·구성하고, 이 구성이 '삼중항 $(A, M, \varphi)$에 관하여 함자적'임을 증명하시오. \item % P12-ERR-0428: frozen determinant assertion without equal-rank hypothesis; see ../control/ERRATA_CANDIDATES.jsonl.
$M, N$가 유한 국소 자유 $A$-가군이며, $\varphi : M \to N$ 및 $\psi : N \to M$라면, $\text{Trace}(\varphi \circ \psi) = \text{Trace}(\psi \circ \varphi)$ 및 $\det(\varphi \circ \psi) = \det(\psi \circ \varphi)$ 임을 보이라. \item $M$가 유한 국소 자유인 경우, $\text{Trace}$가 $A$-선형 사상 $\text{End}_A(M) \to A$를 정의하고, $\det$가 곱셈 사상 $\text{End}_A(M) \to A$를 정의함을 보이라. \end{enumerate} \end{exercise}

\begin{exercise} \label{exercises-exercise-trace-det-rings} 이번에는 $B$를, $A$-가군으로서 유한 국소 자유인 $A$-대수로 한다. 다시 말하면, $B$은 유한 국소 자유 $A$-대수이다. \begin{enumerate} \item 연습 \ref{exercises-exercise-trace-det} 의 $\text{Trace}$와 $\det$을 사용하여, $\text{Trace}_{B/A}$와 $\text{Norm}_{B/A}$를 정의하라. \item $b\in B$로 하고, $\pi : \Spec(B) \to \Spec(A)$를 유도되는 사상으로 하자. $\pi(V(b)) = V(\text{Norm}_{B/A}(b))$를 증명하라. ($V(f) = \{ {\mathfrak p} \mid f \in {\mathfrak p}\}$를 떠올려라.) \item (기저 변환. )$i : A \to A'$를 환 준동형이라 하고, $B' = B \otimes_A A'$라고 둔다. $i(\text{Norm}_{B/A}(b))$가 $\text{Norm}_{B'/A'}(b \otimes 1)$와 같은 이유를 설명하시오. \item $B = A \times A \times A \times \ldots \times A$ 및 $b = (a_1, \ldots, a_n)$의 경우, $\text{Norm}_{B/A}(b)$를 계산하시오. \item 유한 평탄 사상 ${\mathbf Q}[x] \to {\mathbf Q}[y]$, $x \to y^n$에 관한 $y-y^3$의 놈을 계산하시오. (힌트: "기저 변환" $A = {\mathbf Q}[x] \subset A' = {\mathbf Q}(\zeta_n)(x^{1/n})$을 사용하라.) \end{enumerate} \end{exercise}


\section{이어붙이기} \label{exercises-section-glueing}

\begin{exercise} \label{exercises-exercise-cover} $A$를 환으로 하고, $M$를 $A$-가군으로 한다. $f_i$, $i \in I$를, 다음을 만족하는 $A$의 원소들의 집합으로 한다: $$
\Spec(A) = \bigcup D(f_i).
$$ \begin{enumerate} \item $M_{f_i}$가 유한 $A_{f_i}$-가군이면, $M$가 유한 $A$-가군임을 보이라. \item $M_{f_i}$가 평탄 $A_{f_i}$-가군이면, $M$가 평탄 $A$-가군임을 보여라. (올바른 방식으로 생각하면, 이것은 다소 사소한 일이다.) \end{enumerate} \end{exercise}

\begin{remark} \label{exercises-remark-cover} 대수기하학의 언어로 말하면, 이것은 '유한 생성'이거나 '평탄'이라는 성질이 (적절한 의미에서) Zariski 위상에 대해 국소적임을 의미한다. '유한 표현'이라는 성질에 대해서도 같은 것을 보일 수 있다. \end{remark}

\begin{exercise} \label{exercises-exercise-cover-ring-map} $A \to B$를 환 준동형이라 하자. $f_i \in A$, $i \in I$ 및 $g_j \in B$, $j \in J$를 다음을 만족하는 원들의 집합으로 하자: $$
\Spec(A) = \bigcup D(f_i)
\quad\text{그리고}\quad
\Spec(B) = \bigcup D(g_j).
$$ 모든 $i, j$에 대해 $A_{f_i} \to B_{f_ig_j}$가 유한형이면, $A \to B$가 유한형임을 보이라. \end{exercise}


\section{상승정리와 하강정리} \label{exercises-section-going-up}

\begin{definition} \label{exercises-definition-GU-GD} $\phi : A \to B$를 환 준동형이라고 하자. 다음 조건을 만족할 때, $\phi$에 대해 {\it 상승정리} 가 성립한다고 한다: \begin{itemize} \item[(GU)] ${\mathfrak p} \subset {\mathfrak p}'$를 만족하는 임의의 ${\mathfrak p}, {\mathfrak p}' \in \Spec(A)$와, ${\mathfrak p}$ 위에 있는 임의의 $P \in \Spec(B)$에 대해, $P \subset P'$를 만족하고 ${\mathfrak p}'$ 위에 있는 $P'\in \Spec(B)$가 존재한다. \end{itemize} 마찬가지로, 다음 조건을 만족할 때, $\phi$에 대해 {\it 하강정리} 가 성립한다고 한다: \begin{itemize} \item[(GD)] ${\mathfrak p} \subset {\mathfrak p}'$를 만족하는 임의의 ${\mathfrak p}, {\mathfrak p}' \in \Spec(A)$와, ${\mathfrak p}'$ 위에 있는 임의의 $P' \in \Spec(B)$에 대해, $P \subset P'$를 만족하고 ${\mathfrak p}$ 위에 있는 $P\in \Spec(B)$가 존재한다. \end{itemize} \end{definition}

\begin{exercise} \label{exercises-exercise-GU-GD} 다음 각 경우에 (GU), (GD) 가 성립하는지 판단하고 그 이유를 설명하라. (찾을 수 있는 어떤 명제·정리·보조정리를 사용해도 좋지만, 각 경우에 가정을 확인하라.) \begin{enumerate} \item $k$는 체이며, $A = k$, $B = k[x]$. \item $k$는 체이며, $A = k[x]$, $B = k[x, y]$. \item $A = {\mathbf Z}$, $B = {\mathbf Z}[1/11]$. \item $k$는 대수적으로 닫힌 체이며, $A = k[x, y]$, $B = k[x, y, z]/(x^2-y, z^2-x)$. \item $A = {\mathbf Z}$, $B = {\mathbf Z}[i, 1/(2 + i)]$. \item $A = {\mathbf Z}$, $B = {\mathbf Z}[i, 1/(14 + 7i)]$. \item $k$는 대수적으로 닫힌 체이며, $A = k[x]$, $B = k[x, y, 1/(xy-1)]/(y^2-y)$. \end{enumerate} \end{exercise}

\begin{exercise} \label{exercises-exercise-image} $A$를 환으로 하고, $B = A[x]$를 $A$ 상의 일변수 다항식 대수로 한다. $f = a_0 + a_1 x + \ldots + a_r x^r \in B = A[x]$로 한다. $\Spec(A)$ 에서의 $D(f)$의 상이 $D(a_0) \cup \ldots \cup D(a_r)$와 같음을 자세히 증명하라. \end{exercise}

\begin{exercise} \label{exercises-exercise-images} $k$를 대수적으로 닫힌 체로 한다. 다음 각 사상의 $\Spec(k[x, y])$ 에서의 상을 계산하라: \begin{enumerate} \item $\Spec(k[x, yx^{-1}]) \to \Spec(k[x, y])$, 여기서 $k[x, y] \subset k[x, yx^{-1}] \subset k[x, y, x^{-1}]$. (힌트: 혼동을 피하기 위해, 원소 $yx^{-1}$에 다른 이름을 붙여라.) \item $\Spec(k[x, y, a, b]/(ax-by-1))\to \Spec(k[x, y])$. \item $x \mapsto t^2$ 및 $y \mapsto t^3$에 의해 유도되는 $\Spec(k[t, 1/(t-1)]) \to \Spec(k[x, y])$. \item $k = {\mathbf C}$(복소수), $\Spec(k[s, t]/(s^3 + t^3-1)) \to \Spec(k[x, y])$, 여기서 $x\mapsto s^2$, $y \mapsto t^2$. \end{enumerate} \end{exercise}

\begin{remark} \label{exercises-remark-elimination-theory} 위와 같이 상을 구하려면, 일반적으로 소거법을 사용한다. \end{remark}



\section{Fitting 아이디얼} \label{exercises-section-fitting-ideals}

\begin{exercise} \label{exercises-exercise-fitting} $R$를 환, $M$를 유한 $R$-가군으로 하자. $M$의 표현 $$
\bigoplus\nolimits_{j \in J} R \longrightarrow R^{\oplus n}
\longrightarrow M \longrightarrow 0.
$$를 선택하라. 사상 $R^{\oplus n} \to M$는 $M$의 원소 열 $x_1, \ldots, x_n$에 의해 주어진다는 것에 주의하라. 원소 $x_i$는 $M$의 {\it 생성원}이다. 사상 $\bigoplus_{j \in J} R \to R^{\oplus n}$은 계수가 $R$에 속하는 $n \times J$ 행렬 $A$에 의해 주어진다. 다시 말해, $A = (a_{ij})_{i = 1, \ldots, n, j \in J}$이다. $A$의 열 $(a_{1j}, \ldots, a_{nj})$,$j \in J$을 {\it 관계식}이라고 한다. $\sum r_i x_i = 0$을 만족하는 임의의 벡터 $(r_i) \in R^{\oplus n}$는 $A$의 열의 선형 결합이다. 물론, 임의의 유한 $R$-가군은 많은 서로 다른 표현을 가진다. \begin{enumerate} \item $A$의 $(n - k) \times (n - k)$ 소행렬식으로 생성되는 아이디얼이, 표현의 선택에 의존하지 않음을 보이라. 이 아이디얼을 {\it $M$의 제 $k$ Fitting 아이디얼}이라고 하며, $Fit_k(M)$라고 쓴다. \item $Fit_0(M) \subset Fit_1(M) \subset Fit_2(M) \subset \ldots$를 보이라. (힌트: 행렬식은 한 열을 따라 전개하여 계산할 수 있다는 점을 이용하라.) \item 다음이 동치임을 보이라: \begin{enumerate} \item $Fit_{r - 1}(M) = (0)$ 그리고 $Fit_r(M) = R$, \item $M$는 계수 $r$의 국소 자유 가군이다. \end{enumerate} \end{enumerate} \end{exercise}


\section{Hilbert 함수} \label{exercises-section-hilbert}

\begin{definition} \label{exercises-definition-numerical-polynomial} {\it 정수값 다항식}이란, 임의의 정수 $n$에 대해 $f(n) \in {\mathbf Z}$가 되는 다항식 $f(x) \in {\mathbf Q}[x]$을 말한다. \end{definition}

\begin{definition} \label{exercises-definition-graded-module} 환 $A$ 위의 {\it 등급 가군} $M$ 이란, $A$-부분가군으로의 직합 분해 $
\bigoplus\nolimits_{n \in {\mathbf Z}} M_n
$를 갖춘 $A$-가군 $M$을 말한다. 모든 $M_n$가 유한 $A$-가군일 때, $M$는 {\it 국소유한}이라고 한다. $A$를 뇌터 환이라고 하고, $\varphi$를 유한 $A$-가군 위의 {\it Euler--Poincar\'e 함수}라고 한다. 이것은, 임의의 유한 생성 $A$-가군 $M$에 정수 $\varphi(M) \in {\mathbf Z}$가 주어지고, 임의의 짧은 완전 열 $$
0
\longrightarrow
M'
\longrightarrow
M
\longrightarrow
M''
\longrightarrow
0
$$에 대하여 $\varphi(M) = \varphi(M') + \varphi(M'')$이 됨을 의미한다. 국소 유한 등급 가군 $M$의 ($\varphi$에 관한) {\it Hilbert 함수}란, 함수 $\chi_\varphi(M, n) = \varphi(M_n)$을 뜻한다. 충분히 큰 모든 정수 $n$에 대해 $\chi_\varphi(M, n) = P_\varphi(n)$가 되는 정수값 다항식 $P_\varphi$가 존재할 때, $M$은 {\it Hilbert 다항식} 을 가진다고 한다. \end{definition}

\begin{definition} \label{exercises-definition-graded-algebra} {\it 차수 있는 $A$-대수}란, 차수 있는 $A$-가군 $B = \bigoplus_{n \geq 0} B_n$과, $A$-쌍선형 사상 $$
B \times B \longrightarrow B, \ (b, b') \longmapsto bb'
$$으로서, $B$를 $A$-대수로 만들고, $B_n \cdot B_m \subset B_{n + m}$를 만족하는 것을 말한다. 마지막으로, {\it 등급 $A$-대수 $B$ 위의 등급 모듈 $M$}란, 등급 $A$-모듈 $M$에 (양립하는) $B$-모듈 구조로 $B_n \cdot M_d \subset M_{n + d}$를 만족시키는 것을 부여한 것이다. 이것으로, {\it 등급 가군·환의 준동형}, {\it 등급 부분가군}, {\it 등급 아이디얼}, {\it 등급 가군의 완전열} 등등을 정의할 수 있다. \end{definition}

\begin{exercise} \label{exercises-exercise-Euler-Poincare-field} $A = k$를 체로 한다. 이 경우, 유한 $A$-가군 위의 Euler--Poincar\'e 함수로서 어떤 것이 있는지, 모두 구하라. \end{exercise}

\begin{exercise} \label{exercises-exercise-Euler-Poincare-Z} $A ={\mathbf Z}$라고 하자. 이 경우, 유한 $A$-가군 위의 Euler--Poincar\'e 함수로서 어떤 것이 있는지, 모두 구하라. \end{exercise}

\begin{exercise} \label{exercises-exercise-Euler-Poincare-node} $A = k[x, y]/(xy)$로 하고, $k$를 대수적으로 닫힌 체로 하자. 이 경우, 유한 $A$-가군 위의 Euler--Poincar\'e 함수로서 어떤 것이 있는지 모두 구하시오. \end{exercise}

\begin{exercise} \label{exercises-exercise-kernel-locally-finite} $A$를 뇌터 환이라고 하자. 국소 유한한 $A$-가군 사상의 핵이 국소 유한임을 보이라. \end{exercise}

\begin{exercise} \label{exercises-exercise-no-hilbert} $k$를 체라고 하고, $A = k$ 및 $B = k[x, y]$라고 하자. 등급은 $\deg(x) = 2$와 $\deg(y) = 3$에 의해 정해진다. $\varphi(M) = \dim_k(M)$로 한다. 등급 $k$-모듈로서 $B$의 Hilbert 함수를 계산하라. 이 경우, Hilbert 다항식이 존재하는가. \end{exercise}

\begin{exercise} \label{exercises-exercise-no-hilbert-or-is-there} $k$를 체로 하고, $A = k$ 및 $B = k[x, y]/(x^2, xy)$로 둔다. 등급은 $\deg(x) = 2$와 $\deg(y) = 3$에 의해 정해진다. $\varphi(M) = \dim_k(M)$라고 하자. 등급 $k$-가군으로서 $B$의 Hilbert 함수를 계산하라. 이 경우, Hilbert 다항식이 존재하는가. \end{exercise}

\begin{exercise} \label{exercises-exercise-hilbert-to-compute} $k$를 체로 하고, $A = k$라고 하자. $\varphi(M) = \dim_k(M)$라고 한다. $d\in {\mathbf N}$을 고정한다. 등급 $A$-대수 $B = k[x, y, z]/(x^d + y^d + z^d)$를 생각한다. 여기서 $x, y, z$는 각각 차수 $1$를 가진다. $B$의 Hilbert 함수를 계산하라. 이 경우, Hilbert 다항식은 존재하는가. \end{exercise}


\section{환의 Proj}
\label{exercises-section-proj-ring}

\begin{definition}
\label{exercises-definition-homogeneous-ideal}
$R$을 등급환이라 하자. {\it 동차} 아이디얼이란 $R$의 등급 부분가군이기도 한 아이디얼 $I \subset R$을 말한다. 동치로, 동차 원소들로 생성되는 아이디얼이다. 또 동치로, $f \in I$이고
$$
f = f_0 + f_1 + \ldots + f_n
$$
가 $R$에서 $f$를 동차 성분들로 분해한 것이면 각 $i$에 대해 $f_i \in I$이다.
\end{definition}

\begin{definition}
\label{exercises-definition-Proj-R}
등급환 $R$의 {\it 동차 스펙트럼 $\text{Proj}(R)$}을 $R_{+} \not \subset {\mathfrak p}$를 만족하는 $R$의 동차 소 아이디얼 ${\mathfrak p}$들의 집합으로 정의한다. $\text{Proj}(R)$는 $\Spec(R)$의 부분집합이므로 자연스러운 유도 위상을 갖는다.
\end{definition}

\begin{definition}
\label{exercises-definition-Dplus-Vplus}
$R = \oplus_{d \geq 0} R_d$를 등급환이라 하고, $f\in R_d$이며 $d \geq 1$이라고 하자. {\it $R_{(f)}$}를, $r$가 동차이고 $\deg(r) = nd$인 $r/f^n$ 꼴의 원소들로 이루어진 $R_f$의 부분환으로 정의한다. 또한
$$
D_{+}(f) = \{ {\mathfrak p} \in \text{Proj}(R) | f \not\in {\mathfrak p} \}.
$$
로 정의한다. 마지막으로 동차 아이디얼 $I \subset R$에 대해 $V_{+}(I) = V(I) \cap \text{Proj}(R)$로 정의한다.
\end{definition}

\begin{exercise}
\label{exercises-exercise-topology-proj}
$\text{Proj}(R)$의 위상에 관하여. 위의 정의와 표기 아래 다음 명제들을 증명하라.
\begin{enumerate}
\item $D_{+}(f)$가 $\text{Proj}(R)$에서 열린집합임을 보여라.
\item $D_{+}(ff') = D_{+}(f) \cap D_{+}(f')$임을 보여라.
\item $g = g_0 + \ldots + g_m$를 $g_i \in R_i$인 $R$의 원소라 하자. $D(g) \cap \text{Proj}(R)$를 $D_{+}(g_i)$, $i \geq 1$과 $D(g_0) \cap \text{Proj}(R)$를 사용하여 나타내라. 증명은 필요 없다.
\item $g\in R_0$를 차수 $0$의 동차 원소라 하자. $D(g) \cap \text{Proj}(R)$를 양의 차수를 갖는 적절한 동차 원소족 $f_\alpha \in R$에 대한 $D_{+}(f_\alpha)$들로 나타내라.
\item 열린집합들의 모임 $\{D_{+}(f)\}$가 $\text{Proj}(R)$의 위상의 기저를 이룸을 보여라.
\item
\label{exercises-item-bijection}
표준적인 전단사 $D_{+}(f) \to \Spec(R_{(f)})$가 존재함을 보여라.
% P12-ERR-0429: frozen imperative typo; see ../control/ERRATA_CANDIDATES.jsonl.
(힌트: $\Spec$의 경우의 증명을 모방하되, 어느 지점에서 아이디얼의 근기를 끼워 넣어라.)
\item (\ref{exercises-item-bijection})의 사상이 위상동형임을 보여라.
\item $\text{Proj}(R)$가 준콤팩트가 아닌 $R$의 예를 들어라. 증명은 필요 없다.
\item 임의의 닫힌 부분집합 $T \subset \text{Proj}(R)$가 어떤 동차 아이디얼 $I \subset R$에 대한 $V_{+}(I)$ 꼴임을 보여라.
\end{enumerate}
\end{exercise}

\begin{remark}
\label{exercises-remark-continuous-proj-spec}
연속사상 $ \text{Proj}(R) \longrightarrow \Spec(R_0) $가 존재한다.
\end{remark}

\begin{exercise}
\label{exercises-exercise-iso-polynomial-ring-one-variable}
$R = A[X]$이고 $\deg(X) = 1$이면 자연스러운 사상 $\text{Proj}(R) \to \Spec(A)$가 전단사이고, 실제로 위상동형임을 보여라.
\end{exercise}

\begin{exercise}
\label{exercises-exercise-blowing-up-I}
블로업: 제I부. 이 연습에서는 $R = Bl_I(A) = A \oplus I \oplus I^2 \oplus \ldots$로 둔다. 자연스러운 사상 $b : \text{Proj}(R) \to \Spec(A)$를 생각하자. $U = \Spec(A) - V(I)$로 두어라.
$$
b : b^{-1}(U) \longrightarrow U
$$
가 위상동형임을 보여라. 따라서 $U$를 $\text{Proj}(R)$의 열린 부분집합으로 생각할 수 있다. $Z \subset \Spec(A)$를 일반점 $\xi \in Z$를 갖는 기약 닫힌 부분스킴이라 하자. $\xi \not\in V(I)$, 다시 말해 $Z \not\subset V(I)$, 다시 말해 $\xi \in U$, 다시 말해 $Z\cap U \not = \emptyset$이라고 가정하자. $Z$의 {\it 고유 변환} $Z'$을 $\xi$ 위에 놓이는 유일한 점 $\xi'$의 폐포로 정의한다. 달리 말하면, $Z'$은 국소 닫힌 부분집합 $Z\cap U \subset U \subset \text{Proj}(R)$의 $\text{Proj}(R)$ 안에서의 폐포이다.
\end{exercise}

\begin{exercise}
\label{exercises-exercise-blowing-up-II}
블로업: 제II부. $A = k[x, y]$이고 $k$가 체라 하며, $I = (x, y)$라 하자. $R$을 $A$와 $I$의 블로업 대수라 하자.
\begin{enumerate}
\item $Z_1 = V(\{x\})$와 $Z_2 = V(\{y\})$의 고유 변환들이 서로소임을 보여라.
\item $Z_1 = V(\{x\})$와 $Z_2 = V(\{x-y^2\})$의 고유 변환들이 서로소가 아님을 보여라.
\item $V(J) = V(I)$를 만족하고, $J$를 따라 블로업했을 때 $Z_1 = V(\{x\})$와 $Z_2 = V(\{x-y^2\})$의 고유 변환들이 서로소가 되게 하는 아이디얼 $J \subset A$를 찾아라.
\end{enumerate}
\end{exercise}

\begin{exercise}
\label{exercises-exercise-proj-when-empty}
$R$을 등급환이라 하자.
\begin{enumerate}
\item
% P12-ERR-0430: frozen double-greater-than tokens; see ../control/ERRATA_CANDIDATES.jsonl.
모든 $n >> 0$에 대해 $R_n = (0)$이면 $\text{Proj}(R)$가 공집합임을 보여라.
\item $R$이 정역이고 $R_{+}$가 영이 아니면 $\text{Proj}(R)$가 기약 위상 공간임을 보여라. (공집합인 위상 공간은 기약이 아님을 상기하라.)
\end{enumerate}
\end{exercise}

\begin{exercise}
\label{exercises-exercise-blowing-up-III}
블로업: 제III부. 고유 변환의 정의에서와 같은 $A$, $I$, $U$, $Z$를 생각하자. 어떤 소 아이디얼 ${\mathfrak p}$에 대해 $Z = V({\mathfrak p})$라 하자. $\bar A = A/{\mathfrak p}$라 하고, $\bar I$를 $\bar A$에서 $I$의 상이라 하자.
\begin{enumerate}
\item
% P12-ERR-0431: frozen spaced definitional-equality tokens; see ../control/ERRATA_CANDIDATES.jsonl.
전사 환 준동형 $R: = Bl_I(A) \to \bar R: = Bl_{\bar I}(\bar A)$가 존재함을 보여라.
\item
% P12-ERR-0432: frozen manual locator 5(b); see ../control/ERRATA_CANDIDATES.jsonl.
위의 환 준동형이 $\text{Proj}(\bar R)$에서 $Z$의 고유 변환 $Z'$ 위로 가는 전단사 사상을 유도함을 보여라. (이는 그리 쉽지 않다. 힌트: 위의 5(b)를 사용하라.)
\item 고유 변환이 $Z' = V_{+}(P)$임을 결론지어라. 여기서 $P \subset R$는 $P_d = I^d \cap {\mathfrak p}$로 정의되는 동차 아이디얼이다.
\item $Z_1 = V({\mathfrak p})$와 $Z_2 = V({\mathfrak q})$가 소 아이디얼들로 정의된 기약 닫힌 부분집합들이고 $Z_1 \not \subset Z_2$이며 $Z_2 \not \subset Z_1$이라고 하자. 아이디얼 $I = {\mathfrak p} + {\mathfrak q}$의 블로업이 $Z_1$과 $Z_2$의 고유 변환들을 분리함을, 즉 $Z_1' \cap Z_2' = \emptyset$임을 보여라. (힌트: (c)의 동차 아이디얼 $P$와 $Q$를 생각하고 $V(P + Q)$를 생각하라.)
\end{enumerate}
\end{exercise}

\section{차원 1인 Cohen--Macaulay 환}
\label{exercises-section-CM-dim-1}

\begin{definition}
\label{exercises-definition-CM}
뇌터 국소환 $A$가 차원 $d$인 {\it Cohen--Macaulay} 환이라는 것은, 그 차원이 $d$이고 $A$의 매개변수계 $x_1, \ldots, x_d$가 존재하여 각 $i = 1, \ldots, d$에 대해 $x_i$가 $A/(x_1, \ldots, x_{i-1})$에서 영인자가 아닌 원소가 되는 것을 말한다.
\end{definition}

\begin{exercise}
\label{exercises-exercise-CM-dim-1-I}
차원 1인 Cohen--Macaulay 환. 제I부: 이론.
\begin{enumerate}
\item
% P12-ERR-0433: source omits the head noun "ring".
$(A, {\mathfrak m})$을 $\dim A = 1$인 뇌터 국소환이라 하자. $x\in {\mathfrak m}$이 영인자가 아니면 다음을 보여라.
\begin{enumerate}
\item $\dim A/xA = 0$, 다시 말해 $A/xA$는 아르틴 환이고, 다시 말해 $\{x\}$는 $A$의 매개변수계이다.
\item
% P12-ERR-0434: duplicated finite verb in the source.
$A$에는 매장 소 아이디얼이 없다.
\end{enumerate}
\item 반대로, $(A, {\mathfrak m})$을 차원 $1$인 뇌터 국소환이라 하자. $A$에 매장 소 아이디얼이 없으면 ${\mathfrak m}$ 안에 영인자가 아닌 원소가 존재함을 보여라.
\end{enumerate}
\end{exercise}

\begin{exercise}
\label{exercises-exercise-CM-dim-1-II}
차원 1인 Cohen--Macaulay 환. 제II부: 예.
\begin{enumerate}
\item $A$를 $k[x, y]/(x^2, xy)$의 $(x, y)$에서의 국소환이라 하자.
\begin{enumerate}
\item $A$의 차원이 1임을 보여라.
\item ${\mathfrak m}\subset A$의 모든 원소가 영인자임을 증명하라.
\item $\dim A/zA = 0$을 만족하는 $z\in {\mathfrak m}$을 찾아라(증명은 필요 없다).
\end{enumerate}
\item $A$를 $k[x, y]/(x^2)$의 $(x, y)$에서의 국소환이라 하자. ${\mathfrak m}$에서 영인자가 아닌 원소를 찾아라(증명은 필요 없다).
\end{enumerate}
\end{exercise}

\begin{exercise}
\label{exercises-exercise-embedding-dim-1}
매장 차원 $1$인 국소환. $(A, {\mathfrak m}, k)$를 매장 차원 $1$인 뇌터 국소환, 즉
$$
\dim_k {\mathfrak m}/{\mathfrak m}^2 = 1.
$$
을 만족하는 환이라 하자. 함수 $f(n) = \dim_k {\mathfrak m}^n/{\mathfrak m}^{n + 1}$가 값 $1$인 상수함수이거나 그 값들이
$$
1, 1, \ldots, 1, 0, 0, 0, 0, 0, \ldots
$$
임을 보여라.
\end{exercise}

\begin{exercise}
\label{exercises-exercise-regular-local-dim-1}
차원 $1$인 정칙 국소환. $(A, {\mathfrak m}, k)$를 차원 $1$인 정칙 뇌터 국소환이라 하자. 이는 $A$의 차원이 $1$이고 매장 차원도 $1$, 즉
$$
\dim_k {\mathfrak m}/{\mathfrak m}^2 = 1.
$$
임을 뜻한다는 것을 상기하라. $x\in{\mathfrak m}$를 ${\mathfrak m}/{\mathfrak
m}^2$에서의 동치류가 영이 아닌 임의의 원소라 하자.
\begin{enumerate}
\item
% P12-ERR-0543: preserve the universal quantifier, including zero.
${\mathfrak m}$의 모든 원소 $y$에 대해 정수 $n$이 존재하여 $y$를 $y = ux^n$로 쓸 수 있음을 보여라. 여기서 $u\in A^\ast$는 단위원이다.
\item $x$가 $A$에서 영인자가 아닌 원소임을 보여라.
\item $A$가 정역임을 결론지어라.
\end{enumerate}
\end{exercise}

\begin{exercise}
\label{exercises-exercise-nonzerodivisor-graded}
$(A, {\mathfrak m}, k)$를 수반 등급환 $Gr_{\mathfrak m}(A)$을 갖는 뇌터 국소환이라 하자.
\begin{enumerate}
\item $x\in {\mathfrak m}^d$가 $Gr_{\mathfrak m}(A)$의 차수 $d$에서 영인자가 아닌 원소 $\bar x \in {\mathfrak m}^d/{\mathfrak m}^{d + 1}$로 사상된다고 하자. $x$가 영인자가 아님을 보여라.
\item $A$의 깊이가 적어도 $1$이라고 하자. 즉, 영인자가 아닌 원소 $y \in {\mathfrak m}$이 존재한다고 하자. 이 경우에는 더 강한 결론을 얻을 수 있다. $x\in {\mathfrak m}^d$가 $Gr_{\mathfrak m}(A)$의 차수 $d$인 원소 $\bar x \in {\mathfrak m}^d/{\mathfrak m}^{d + 1}$로 사상되고, 이것이 충분히 높은 차수들에서 영인자가 아니라고만 가정하자. 즉, $\exists N$이 있어서 모든 $n \geq N$에 대해 $\bar x$의 곱셈 사상
$$
{\mathfrak m}^n/{\mathfrak m}^{n + 1} \longrightarrow
{\mathfrak m}^{n + d}/{\mathfrak m}^{n + d + 1}
$$
이 단사라고 하자. 그러면 $x$가 영인자가 아님을 보여라.
\end{enumerate}
\end{exercise}

\begin{exercise}
\label{exercises-exercise-embedding-2-dim-1}
$(A, {\mathfrak m}, k)$를 차원 $1$인 뇌터 국소환이라 하자. 또한 $A$의 매장 차원이 $2$, 즉
$$
\dim_k {\mathfrak m}/{\mathfrak m}^2 = 2.
$$
라고 가정하자. 표기: $f(n) = \dim_k {\mathfrak m}^n/{\mathfrak m}^{n + 1}$. 생성원 $x, y \in {\mathfrak m}$을 택하고, 어떤 동차 아이디얼 $I$에 대해 $Gr_{\mathfrak m}(A) = k[\bar x, \bar y]/I$로 쓰자.
\begin{enumerate}
\item 동차 원소 $F\in k[\bar x, \bar y]$가 존재하여 $I \subset (F)$이고 충분히 높은 모든 차수에서는 등호가 성립함을 보여라.
\item $f(n) \leq n + 1$임을 보여라.
\item $f(n) < n + 1$이면 $n \geq \deg(F)$임을 보여라.
\item $f(n) < n + 1$이면 $f(n + 1) \leq f(n)$임을 보여라.
\item
% P12-ERR-0443: preserve literal >> in the frozen formula.
모든 $n >> 0$에 대해 $f(n) = \deg(F)$임을 보여라.
\end{enumerate}
\end{exercise}

\begin{exercise}
\label{exercises-exercise-CM-dim-1-embedding-dim-2}
차원 1, 매장 차원 2인 Cohen--Macaulay 환. $(A, {\mathfrak m}, k)$를 차원 $1$인 Cohen--Macaulay 뇌터 국소환이라 하자. 또한 $A$의 매장 차원이 $2$, 즉
$$
\dim_k {\mathfrak m}/{\mathfrak m}^2 = 2.
$$
라고 가정하자. 표기: $f$, $F$, $x, y\in {\mathfrak m}$, $I$는 연습문제 \ref{exercises-exercise-embedding-2-dim-1}에서와 같다. 위 문제들의 결과는 어느 것이든 사용해도 좋다.
\begin{enumerate}
\item $z\in {\mathfrak m}$의 ${\mathfrak m}/{\mathfrak m}^2$에서의 동치류가 $F$와 서로소인 일차형식 $\alpha \bar x + \beta \bar y \in k[\bar x, \bar y]$이라고 하자.
\begin{enumerate}
\item $z$가 $A$에서 영인자가 아닌 원소임을 보여라.
\item $d = \deg(F)$라 하자. 충분히 큰 모든 $n$에 대해 ${\mathfrak m}^n = z^{n + 1-d}{\mathfrak m}^{d-1}$임을 보여라. (힌트: 먼저 $Gr_{\mathfrak m}(A)$에 관해 아는 사실을 이용하여 $z^{n + 1-d}{\mathfrak m}^{d-1} \to {\mathfrak m}^n/{\mathfrak m}^{n + 1}$가 전사임을 보여라. 그런 다음 NAK를 사용하라.)
\end{enumerate}
\item 이러한 $z$의 존재를 보장하는 $k$에 관한 조건은 무엇인가? (증명은 필요 없다. 너무 쉽기 때문이다.)

\noindent
이제 위와 같은 $z$가 존재한다고 가정하겠다. 이는 해롭지 않은 가정임이 드러난다(아래의 (d)와 (e)의 결과를 얻기 위해 이 조건이 성립하는 상황으로 귀착할 수 있다는 뜻이다).
\item 모든 $\ell \geq d$에 대해 ${\mathfrak m}^\ell = z^{\ell - d + 1} {\mathfrak m}^{d-1}$임을 보여라.
\item $I = (F)$임을 결론지어라.
\item
% P12-ERR-0435: the frozen sequence starts at 2 without an explicit starting index.
함수 $f$의 값들이
$$
2, 3, 4, \ldots, d-1, d, d, d, d, d, d, d, \ldots
$$
임을 결론지어라.
\end{enumerate}
\end{exercise}

\begin{remark}
\label{exercises-remark-CM-dim-1-embedding-dim-2}
이는 차원 1, 매장 차원 2인 뇌터 국소 Cohen--Macaulay 환이 $B/FB$ 꼴임을 시사한다. 여기서 $B$는 2차원 정칙 국소환이다. 이는 (환에 적절한 ``좋은 성질''을 가정하면) 거의 참이다.
\end{remark}

\section{무한히 많은 소수}
\label{exercises-section-many-primes}

\noindent
무한히 많은 소수가 가역원이 아닌 환들에 관한 기묘한 문제들을 모은 절이다.

\begin{exercise}
\label{exercises-exercise-not-in-Q}
다음 두 성질을 갖는 유한형 ${\mathbf Z}$-대수 $R$의 예를 들어라.
\begin{enumerate}
\item 환 준동형 $R \to {\mathbf Q}$는 존재하지 않는다.
\item 모든 소수 $p$에 대해 $R/{\mathfrak m} \cong {\mathbf F}_p$를 만족하는 극대 아이디얼 ${\mathfrak m} \subset R$이 존재한다.
\end{enumerate}
\end{exercise}

\begin{exercise}
\label{exercises-exercise-strange-fp-1}
$f \in {\mathbf Z}[x, u]$에 대해 $f_p(x) = f(x, x^p) \bmod p \in {\mathbf F}_p[x]$로 정의한다. 다음 두 성질을 만족하는 $f \in {\mathbf Z}[x, u]$의 예를 들어라.
\begin{enumerate}
\item $f_p$가 ${\mathbf F}_p$에서 영점을 갖지 않는 $p$가 무한히 많이 존재한다.
\item 모든 $p >> 0$에 대해 다항식 $f_p$는 일차 또는 이차 인수를 갖는다.
\end{enumerate}
\end{exercise}

\begin{exercise}
\label{exercises-exercise-strange-fp-2}
$f \in {\mathbf Z}[x, y, u, v]$에 대해 $f_p(x, y) = f(x, y, x^p, y^p) \bmod p \in {\mathbf F}_p[x, y]$로 정의한다. 모든 $p >> 0$에 대해 $f_p$가 가약인 ``흥미로운'' $f$의 예를 들어라. 예를 들어 $f = xv-yu$이면 $f_p = xy^p-x^py = xy(x^{p-1}-y^{p-1})$이므로 ``흥미롭지 않다''. $x, u$에만 의존하는 임의의 $f$도 ``흥미롭지 않다'' 등등.
\end{exercise}

\begin{remark}
\label{exercises-remark-strange-fp}
$h \in {\mathbf Z}[y]$를 차수 $d$인 일계수 다항식이라 하자. 그러면 다음이 성립한다.
\begin{enumerate}
\item 사상 $A = {\mathbf Z}[x] \to B ={\mathbf Z}[y]$, $x \mapsto h$는 계수 $d$인 유한 국소 자유 사상이다.
\item 모든 소수 $p$에 대해 사상 $A_p = {\mathbf F}_p[x]\to B_p = {\mathbf F}_p[y]$, $y \mapsto h(y) \bmod p$는 계수 $d$인 유한 국소 자유 사상이다.
\end{enumerate}
\end{remark}

\begin{exercise}
\label{exercises-exercise-strange-fp-3}
$h, A, B, A_p, B_p$를 위 주석에서와 같이 두자. $f \in {\mathbf Z}[x, u]$에 대해 $f_p(x) = f(x, x^p) \bmod p \in {\mathbf F}_p[x]$로 정의한다. $g \in {\mathbf Z}[y, v]$에 대해서는 $g_p(y) = g(y, y^p) \bmod p \in {\mathbf F}_p[y]$로 정의한다.
\begin{enumerate}
\item 다음 성질을 갖는 $f$가 존재하지 않게 하는 $h$와 $g$의 예를 들어라.
$$
f_p  =  Norm_{B_p/A_p}(g_p).
$$
\item 위와 같은 $h$와 $g$를 어떻게 택하더라도, 영이 아닌 $f$가 존재하여 모든 $p$에 대해
$$
Norm_{B_p/A_p}(g_p)\quad\text{가 다음을 나눈다}\quad f_p .
$$
가 성립함을 보여라. 원한다면 $h = y^n$, 심지어 $n = 2$인 경우로 제한해도 되지만, 이는 일반적으로 참이다.
\item 이것이 앞 문제 모음의 연습문제 6과 7에 어떤 관련이 있는지 논하라.
\end{enumerate}
\end{exercise}

\begin{exercise}
\label{exercises-exercise-strange-fp-unsolved}
미해결 문제. 매우 어려울 수도 있고 쉬울 수도 있다. 나도 모른다.
\begin{enumerate}
\item 무한히 많은 $p$에 대해 $f_p$가 기약이 되는 $f \in {\mathbf Z}[x, u]$가 존재하는가? (힌트: 그렇다. $f(x, u) = u - x - 1$과 $f(x, u) = u^2 - x^2 + 1$이 그 예이다.)
\item $f \in {\mathbf Z}[x, u]$가 영이 아니고, 충분히 큰 모든 $p$에 대해 $\deg_x(f_p) = dp + d'$라고 하자. (다시 말해 $\deg_u(f) = d$이고, $f$에서 $u^d$의 계수 $c$가 $\deg_x(c) = d'$를 만족한다.) $d_1, d_2 > 0$이고 $d'_1, d'_2 \geq 0$인 수들을 택하여 $d = d_1 + d_2$와 $d' = d'_1 + d'_2$로 쓸 수 있으며, 충분히 큰 모든 $p$에 대해 인수분해
$$
f_p = f_{1, p} f_{2, p}
$$
가 존재하여 $\deg_x(f_{1, p}) = d_1p + d'_1$가 된다고 하자. 그렇다면 $f$는 연습문제 4에서와 같은 노름 구성에서 얻어지는가? (더 정확히 말해, 모든 $p >> 0$에 대해 $Norm_{B_p/A_p}(g_p)$가 $f_p$를 나누게 하는 $h$와 $g$가 존재하는가?)
\item (b)와 비슷한 문제를 생각하되, 이제 $f \in {\mathbf Z}[x_1, x_2, u_1, u_2]$가 기약이고 모든 $p >> 0$에 대해 $f_p(x_1, x_2) = f(x_1, x_2, x_1^p, x_2^p) \bmod p$가 인수분해된다고만 가정하라.
\end{enumerate}
\end{exercise}

\section{여과 유도 범주}
\label{exercises-section-filtered-derived}

\noindent
이 절의 연습문제를 풀기 위해서는 호몰로지, 제\ref{homology-section-filtrations}절의 내용을 읽어라. $A$가 $\mathcal{A}$의 여과 대상이라고 말할 때에는, 표기에서는 생략한 여과 $F$가 $A$에 주어져 있다는 뜻으로 사용하겠다.

\begin{exercise}
\label{exercises-exercise-split-injective}
$\mathcal{A}$를 아벨 범주라 하자. $I$를 $\mathcal{A}$의 여과 대상이라 하자. $I$의 여과가 유한하고 각 $\text{gr}^p(I)$가 $\mathcal{A}$의 단사 대상이라고 가정하자. 여과 $F^p(I)$가 $\bigoplus_{p' \geq p} \text{gr}^p(I)$에 대응하는 동형사상 $I \cong \bigoplus \text{gr}^p(I)$이 존재함을 보여라.
\end{exercise}

\begin{exercise}
\label{exercises-exercise-filtered-injective}
$\mathcal{A}$를 아벨 범주라 하자. $I$를 $\mathcal{A}$의 여과 대상이라 하자. $I$의 여과가 유한하다고 가정하자. 다음 두 조건이 동치임을 보여라.
\begin{enumerate}
\item 여과 대상들의 임의의 실선 도식
$$
\xymatrix{
A \ar[r]_\alpha \ar[d] & B \ar@{-->}[ld] \\
I &
}
$$
에서 (\romannumeral1) $A$와 $B$의 여과가 유한하고, (\romannumeral2) $\text{gr}(\alpha)$가 단사이면 도식을 가환하게 만드는 점선 화살표가 존재한다.
\item 각 $\text{gr}^p I$가 단사 대상이다.
\end{enumerate}
\end{exercise}

\noindent
유한 여과를 갖는 여과 대상들의 사상 $\alpha : A \to B$가 주어졌을 때 $\text{gr}(\alpha)$가 단사라는 것은 $\alpha$가 범주 $\text{Fil}(\mathcal{A})$에서 {\it 엄밀 단사 사상}이라는 것과 같다. 실제로 단사 사상이라는 것은 $\Ker(\alpha) = 0$이라는 뜻이고, 엄밀하다는 것은 이것이 $\Ker(\text{gr}(\alpha)) = 0$도 함의한다는 뜻이다. 호몰로지, 보조정리 \ref{homology-lemma-characterize-strict}를 보라. (단사라는 말은 아벨 범주의 사상에 대해서만 사용하겠다. 다만 핵을 갖는 임의의 가법 범주에서도 이 말은 의미가 있다.) 위 연습문제들은 다음 정의를 정당화한다.

\begin{definition}
\label{exercises-definition-injective-filtered}
$\mathcal{A}$를 아벨 범주라 하자. $I$를 $\mathcal{A}$의 여과 대상이라 하자. $I$의 여과가 유한하다고 가정하자. 각 $\text{gr}^p(I)$가 $\mathcal{A}$의 단사 대상이면 $I$를 {\it 여과 단사 대상}이라고 한다.
\end{definition}

\noindent
매번 ``유한 여과를 갖는''이라고 말하지 않아도 되도록 다음을 정의한다.

\begin{definition}
\label{exercises-definition-finite-filtration-category}
$\mathcal{A}$를 아벨 범주라 하자. $\text{Fil}(\mathcal{A})$의 대상 중 여과가 유한한 $A \in \Ob(\text{Fil}(\mathcal{A}))$들로 이루어진 충만 부분범주를 {\it $\text{Fil}^f(\mathcal{A})$}로 나타낸다.
\end{definition}

\begin{exercise}
\label{exercises-exercise-inject-into-injective}
$\mathcal{A}$를 아벨 범주라 하자. $\mathcal{A}$에 단사 대상이 충분히 많다고 가정하자. $A$를 $\text{Fil}^f(\mathcal{A})$의 대상이라 하자. $A$에서 $\text{Fil}^f(\mathcal{A})$의 여과 단사 대상 $I$로 가는 엄밀 단사 사상 $\alpha : A \to I$가 존재함을 보여라.
\end{exercise}

\begin{definition}
\label{exercises-definition-filtered-quasi-isomorphism}
$\mathcal{A}$를 아벨 범주라 하자. $\alpha : K^\bullet \to L^\bullet$를 $\text{Fil}(\mathcal{A})$의 복합체들의 사상이라 하자. 각 $p \in \mathbf{Z}$에 대해 사상 $\text{gr}^p(K^\bullet) \to \text{gr}^p(L^\bullet)$가 준동형이면 $\alpha$를 {\it 여과 준동형}이라고 한다.
\end{definition}

\begin{definition}
\label{exercises-definition-filtered-acyclic}
$\mathcal{A}$를 아벨 범주라 하자. $K^\bullet$을 $\text{Fil}^f(\mathcal{A})$의 복합체라 하자. 각 $p \in \mathbf{Z}$에 대해 복합체 $\text{gr}^p(K^\bullet)$이 비순환이면 $K^\bullet$을 {\it 여과 비순환}이라고 한다.
\end{definition}

\begin{exercise}
\label{exercises-exercise-filtered-quasi-isomorphism}
$\mathcal{A}$를 아벨 범주라 하자. $\alpha : K^\bullet \to L^\bullet$를 $\text{Fil}^f(\mathcal{A})$의 아래로 유계인 복합체들의 사상이라 하자. (위 첨자 $f$에 주의하라.) 다음 조건들이 동치임을 보여라.
\begin{enumerate}
\item $\alpha$는 여과 준동형이다.
\item 각 $p \in \mathbf{Z}$에 대해 사상 $\alpha : F^pK^\bullet \to F^pL^\bullet$은 준동형이다.
\item 각 $p \in \mathbf{Z}$에 대해 사상 $\alpha : K^\bullet/F^pK^\bullet \to L^\bullet/F^pL^\bullet$은 준동형이다.
\item $\alpha$의 원뿔(유도 범주, 정의 \ref{derived-definition-cone} 참조)은 여과 비순환 복합체이다.
\end{enumerate}
또한 $\alpha$가 여과 준동형이면 $\alpha$가 보통 의미에서도 준동형임을 보여라.
\end{exercise}

\begin{exercise}
\label{exercises-exercise-injective-resolution}
$\mathcal{A}$를 아벨 범주라 하자. $\mathcal{A}$에 단사 대상이 충분히 많다고 가정하자. $A$를 $\text{Fil}^f(\mathcal{A})$의 대상이라 하자. $\text{Fil}^f(\mathcal{A})$의 복합체 $I^\bullet$과 사상 $A[0] \to I^\bullet$이 존재하여 다음을 만족함을 보여라.
\begin{enumerate}
\item 각 $I^p$는 여과 단사 대상이다.
\item $p < 0$이면 $I^p = 0$이다.
\item $A[0] \to I^\bullet$은 여과 준동형이다.
\end{enumerate}
\end{exercise}

\begin{exercise}
\label{exercises-exercise-injective-resolution-complex}
$\mathcal{A}$를 아벨 범주라 하자. $\mathcal{A}$에 단사 대상이 충분히 많다고 가정하자. $K^\bullet$을 $\text{Fil}^f(\mathcal{A})$의 대상들로 이루어진 아래로 유계인 복합체라 하자. 여과 준동형 $\alpha : K^\bullet \to I^\bullet$이 존재하되, $I^\bullet$은 여과 단사 항 $I^n$을 갖는 $\text{Fil}^f(\mathcal{A})$의 아래로 유계인 복합체임을 보여라. 실제로 각 $\alpha^n$이 엄밀 단사 사상이 되도록 $\alpha$를 택할 수 있다.
\end{exercise}

\begin{exercise}
\label{exercises-exercise-morphisms-lift}
$\mathcal{A}$를 아벨 범주라 하자. $\text{Fil}^f(\mathcal{A})$의 복합체들로 이루어진 실선 도식
$$
\xymatrix{
K^\bullet \ar[r]_\alpha \ar[d]_\gamma & L^\bullet \ar@{-->}[dl]^\beta \\
I^\bullet
}
$$
을 생각하자. $K^\bullet$, $L^\bullet$, $I^\bullet$이 아래로 유계이고 각 $I^n$이 여과 단사 대상이라고 가정하자. 또한 $\alpha$가 여과 준동형이라고 가정하자.
\begin{enumerate}
\item 호모토피까지 도식을 가환하게 만드는 복합체 사상 $\beta$가 존재한다.
\item $\alpha$가 각 차수에서 엄밀 단사 사상이면 도식을 가환하게 만드는 $\beta$를 찾을 수 있다.
\end{enumerate}
\end{exercise}

\begin{exercise}
\label{exercises-exercise-acyclic-is-zero}
$\mathcal{A}$를 아벨 범주라 하자.
% Frozen source repeats K instead of introducing I here.
$K^\bullet$, $K^\bullet$을 $\text{Fil}^f(\mathcal{A})$의 복합체들이라 하자. 다음을 가정하자.
\begin{enumerate}
\item $K^\bullet$은 아래로 유계이고 여과 비순환이다.
\item $I^\bullet$은 아래로 유계이고 여과 단사 대상들로 이루어져 있다.
\end{enumerate}
그러면 임의의 사상 $K^\bullet \to I^\bullet$은 영 사상과 호모토픽이다.
\end{exercise}

\begin{exercise}
\label{exercises-exercise-morphisms-equal-up-to-homotopy}
$\mathcal{A}$를 아벨 범주라 하자. $\text{Fil}^f(\mathcal{A})$의 복합체들로 이루어진 실선 도식
$$
\xymatrix{
K^\bullet \ar[r]_\alpha \ar[d]_\gamma & L^\bullet \ar@{-->}[dl]^{\beta_i} \\
I^\bullet
}
$$
을 생각하자. $K^\bullet$, $L^\bullet$, $I^\bullet$이 아래로 유계이고 각 $I^n$이 여과 단사 대상이라고 가정하자. 또한 $\alpha$가 여과 준동형이라고 가정하자. 호모토피까지 도식을 가환하게 만드는 임의의 두 사상 $\beta_1, \beta_2$는 서로 호모토픽이다.
\end{exercise}

% Frozen source: exercises.tex lines 2572--2649.
\section{정칙 함수} \label{exercises-section-regular-functions}

\begin{exercise} \label{exercises-exercise-extra-function} 좌표 $s, t$를 갖는 $\mathbf{C}^2$에서 방정식 $t^2 = s^5 + 8$로 주어지는 아핀 곡선 $X$를 생각하자. $x \in X$를 좌표가 $(1, 3)$인 점이라 하고 $U = X \setminus \{x\}$로 두자. $U$ 위의 정칙 함수 가운데 $\mathbf{C}^2$ 위의 정칙 함수의 제한이 아닌 것, 즉 $s$와 $t$에 관한 다항식의 $U$로의 제한이 아닌 것이 존재함을 증명하라. \end{exercise}

\begin{exercise} \label{exercises-exercise-no-extra-function} $n \geq 2$라고 하자. $E \subset \mathbf{C}^n$를 유한 부분집합으로 하자. $\mathbf{C}^n \setminus E$ 위의 임의의 정칙 함수가 다항식임을 보이라. \end{exercise}

\begin{exercise} \label{exercises-exercise-cone} $X \subset \mathbf{C}^n$을 아핀 다양체라고 한다. $x = (a_1, \ldots, a_n) \in X$ 및 $\lambda \in \mathbf{C}$로부터 $(\lambda a_1, \ldots, \lambda a_n) \in X$가 따를 때, $X$는 {\it 원뿔}이라고 한다. 물론, $\mathfrak p \subset \mathbf{C}[x_1, \ldots, x_n]$가 $x_1, \ldots, x_n$의 동차 다항식으로 생성된 소 아이디얼이라면, 아핀 다양체 $X = V(\mathfrak p) \subset \mathbf{C}^n$는 원뿔이다. 반대로, 원뿔에 대응하는 소아이디얼 $\mathfrak p \subset \mathbf{C}[x_1, \ldots, x_n]$은 $x_1, \ldots, x_n$의 동차 다항식으로 생성할 수 있음을 보여라. \end{exercise}

\begin{exercise} \label{exercises-exercise-extra-function-cone} 원뿔인 아핀 다양체 $X \subset \mathbf{C}^n$(연습문제 \ref{exercises-exercise-cone} 참조)와, $U = X \setminus \{(0, \ldots, 0)\}$ 위의 정칙 함수 $f$ 중 $\mathbf{C}^n$ 위의 다항식 함수의 제한이 아닌 것의 예를 제시하라. \end{exercise}

\begin{exercise} \label{exercises-exercise-regular-functions} 이 연습에서는, 대수적으로 닫히지 않은 체 위의 정칙 함수에 무슨 일이 일어나는지 조사한다. $k$을 체로 하자. $Z \subset k^n$를 Zariski 국소 닫힌 부분집합, 즉 아이디얼 $I \subset J \subset k[x_1, \ldots, x_n]$가 존재하여 $$
Z = \{a \in k^n \mid
f(a) = 0\ \forall\ f \in I,\ \exists\ g \in J,\ g(a) \not = 0\}.
$$가 되도록 하는 것으로 하자. 함수 $\varphi : Z \to k$가 {\it 정칙}이라는 것은, 각 $z \in Z$에 대해 Zariski 열린근방 $z \in U \subset Z$과 다항식 $f, g \in k[x_1, \ldots, x_n]$이 존재하여, 모든 $u \in U$에 대해 $g(u) \not = 0$이며, 모든 $u \in U$에 대해 $\varphi(u) = f(u)/g(u)$가 됨을 말한다. \begin{enumerate} \item $k = \bar k$이고 $Z = k^n$라면, 정칙 함수는 다항식으로 주어짐을 보이시오. (이 논의를 이전에 본 적이 없다면 시도하면 된다.) \item $k$가 유한하다면, (a) 모든 함수 $\varphi$는 정칙 함수이며, (b) 정칙 함수의 환은 $k$ 위에서 유한 차원임을 보이라. (원한다면 $Z = k^n$, 나아가 $n = 1$로도 좋다.) \item $k = \mathbf{R}$의 경우, $Z = \mathbf{R}$ 위의 정칙 함수 중 다항식으로 주어지지 않은 예를 제시하라. \item $k = \mathbf{Q}_p$의 경우, $Z = \mathbf{Q}_p$ 위의 정칙 함수로, 다항식으로 주어지지 않는 것의 예를 들어라. \end{enumerate} \end{exercise}


% Frozen source: exercises.tex lines 2650--2810.
\section{층} \label{exercises-section-sheaves}

\noindent 범주 $\mathcal{C}$의 사상 $f : X \to Y$가 {\it 단사 사상}이라는 것은, 임의의 두 사상 $a, b : W \to X$에 대해 $f \circ a = f \circ b \Rightarrow a = b$가 성립하는 것을 말한다. 집합의 범주에서 단사 사상은 집합 사이의 단사 함수이다.

\begin{exercise} \label{exercises-exercise-mono-sheaves-sets} % P12-ERR-0445: source indefinite-article error.
집합층의 사상이 (집합층의 범주에서) 단사 사상인 것과 모든 줄기에 유도되는 사상이 단사인 것이 서로 동치임을, 자세히 증명하라. \end{exercise}

\noindent 범주 $\mathcal{C}$의 사상 $f : X \to Y$가 {\it 동형사상}이라는 것은, 사상 $g : Y \to X$가 존재하여 $f \circ g = \text{id}_Y$ 및 $g \circ f = \text{id}_X$가 되는 것을 말한다. 집합의 범주에서 동형사상은 집합 사이의 전단사 함수이다.

\begin{exercise} \label{exercises-exercise-isomorphism-sheaves-sets} 집합층의 사상이 (집합층의 범주에서) 동형인 것과, 모든 줄기(stalk)에 유도되는 사상이 전단사인 것이 동치임을, 정성껏 증명하시오. \end{exercise}

\noindent 범주 $\mathcal{C}$의 사상 $f : X \to Y$가 {\it 전사 사상}이라는 것은, 임의의 두 사상 $a, b : Y \to Z$에 대해 $a \circ f = b \circ f \Rightarrow a = b$가 성립하는 것을 말한다. 집합의 범주에서 전사 사상은 집합 사이의 전사 함수이다.

\begin{exercise} \label{exercises-exercise-epi-sheaves-sets} 집합의 층의 사상이 (집합의 층의 범주에서) 전사 사상인 것과 모든 줄기에 유도되는 사상이 전사인 것이 서로 동치임을 자세히 증명하라. \end{exercise}

\begin{exercise} \label{exercises-exercise-adjoint-push-pull} $f : X \to Y$를 위상 공간들의 사상이라 하자. {\bf 집합}의 층에 대한 직접상 $f_\ast$와 역상 $f^{-1}$이 한 쌍의 수반 함자를 이룸을 증명하라. \end{exercise}

\begin{exercise} \label{exercises-exercise-j-shriek} $j : U \to X$를 열 포함으로 한다. 다음을 보이라. \begin{enumerate} \item 역상 $j^{-1} : \Sh(X) \to \Sh(U)$은, {\it 공집합에 의한 확장}이라고 불리는 좌수반 $j_{!} : \Sh(U) \to \Sh(X)$를 가진다. \item $\mathcal{G} \in \Sh(U)$에 대해 $j_{!}({\mathcal G})$의 줄기를 특징지어라. \item 역상 $j^{-1} : \textit{Ab}(X) \to \textit{Ab}(U)$은, {\it 영에 의한 연장}이라고 불리는 좌수반 $j_{!} : \textit{Ab}(U) \to \textit{Ab}(X)$을 가진다. \item $\mathcal{G} \in \textit{Ab}(U)$에 대해 $j_{!}({\mathcal G})$의 줄기를 특징짓는다. \end{enumerate} 영에 의한 연장은 공집합에 의한 연장과 다르다는 점을 주의하라! \end{exercise}

\begin{exercise} \label{exercises-exercise-not-locally-generated-by-sections} $X = \mathbf{R}$에 보통 위상을 주자. $\mathcal{O}_X = \underline{\mathbf{Z}/2\mathbf{Z}}_X$로 두자. $i : Z = \{0\} \to X$를 포함사상이라 하고, $\mathcal{O}_Z = \underline{\mathbf{Z}/2\mathbf{Z}}_Z$로 두자. 다음을 증명하라(처음 세 명제는 정의에서 따르지만, 정의가 분명하지 않다면 이 명제들을 통해 분명히 해야 한다). \begin{enumerate} \item $i_*\mathcal{O}_Z$는 마천루층이다. \item 표준적인 전사 사상 $\underline{\mathbf{Z}/2\mathbf{Z}}_X \to
i_*\underline{\mathbf{Z}/2\mathbf{Z}}_Z$가 존재한다. 그 핵을 $\mathcal{I} \subset \mathcal{O}_X$로 나타내자. \item $\mathcal{I}$은 $\mathcal{O}_X$의 아이디얼층이다. \item $X$ 위의 층 $\mathcal{I}$은 (가군, 정의 \ref{modules-definition-locally-generated}의 의미에서) 단면들로 국소 생성될 수 없다. \end{enumerate} \end{exercise}

\begin{exercise} \label{exercises-exercise-quotient-j-shriek-Z} $X$를 위상 공간으로 하자. ${\mathcal F}$를 $X$ 위의 아벨 군 층으로 하자. ${\mathcal F}$가 각 항이 $$
j_{!}(\underline{{\mathbf Z}}_U)
$$ 형태인 층들의 (매우 클 수도 있는) 직합의 몫임을 보여라. 여기서 $U \subset X$는 열린집합이고, $\underline{{\mathbf Z}}_U$는 $U$ 위에서 값이 ${\mathbf Z}$인 상수층을 나타낸다. \end{exercise}

\begin{remark} \label{exercises-remark-direct-sum-stalk-abelian} $X$를 위상 공간으로 한다. 아벨 군의 층 범주에서 층들의 가족 $\{{\mathcal F}_i\}_{i\in I}$의 직합은 % P12-ERR-0446: preserve the unindexed direct sum in the frozen presheaf.
전층 $U \mapsto \oplus {\mathcal F}_i(U)$에 수반된 층이다. 따라서 직합의 점 $x$에서의 줄기는 ${\mathcal F}_i$의 $x$에서의 줄기의 직합이다. \end{remark}

\begin{exercise} \label{exercises-exercise-product-over-points} $X$를 위상 공간으로 한다. $x \in X$에서 첨자가 붙은 아벨 군들의 모임 $A_x$가 주어졌다고 가정하자. 대응 $U \mapsto \prod_{x \in U} A_x$와 명백한 제한 사상에 의해 아벨 군의 층 $\mathcal{G}$가 정의됨을 보이라. 일반적으로 $x \in X$에 대해 $\mathcal{G}_x = A_x$이 되는 것은 아님을 예를 들어 보이라. \end{exercise}

\begin{exercise} \label{exercises-exercise-modified-product-over-points} $X$, $A_x$, $\mathcal{G}$는 연습문제 \ref{exercises-exercise-product-over-points}에서와 같이 두자. $\mathcal{B}$를 $X$의 위상의 기저라 하자. 위상, 정의 \ref{topology-definition-base}를 참조하라. $U \in \mathcal{B}$에 대해 $A_U$를 부분군 $A_U \subset \mathcal{G}(U) = \prod_{x \in U} A_x$라 하자. $U, V \in \mathcal{B}$이고 $U \subset V$이면 제한 사상이 $A_V$를 $A_U$로 보낸다고 가정하자. 열린집합 $U \subset X$에 대해 $$
\mathcal{F}(U) =
\left\{
(s_x)_{x \in U}
\middle|
\begin{matrix}
\text{모든 }x\text{에 대해, }U\text{의 원소라면 다음과 같은 원소가 존재한다: } V \in \mathcal{B} \\
x \in V \subset U\text{이고 } (s_y)_{y \in V} \in A_V
\end{matrix}
\right\}
$$라고 하자. $\mathcal{F}$가 $X$ 위의 아벨 군 층을 정의함을 보이라. 보통은 $U \in \mathcal{B}$에 대해 $\mathcal{F}(U) = A_U$가 되지 않는다는 것을 예로 보여주라. \end{exercise}

\begin{exercise} \label{exercises-exercise-exact-but-not-a-stalk-functor} 위상 공간 $X$와 함자 $$
F : \Sh(X) \longrightarrow \textit{Sets}
$$를 찾아라. 이 함자는 완전(유한곱 및 등화자와 가환하고 유한 쌍대곱 및 쌍대등화자와 가환함; 범주, 제\ref{categories-section-exact-functor}절 참조)이지만, $F$가 줄기 함자 $\mathcal{F} \mapsto \mathcal{F}_x$와 동형이 되는 점 $x \in X$는 존재하지 않아야 한다. \end{exercise}

% Frozen source: exercises.tex lines 2811--2969.
\section{스킴} \label{exercises-section-schemes}

\noindent $LRS$를 국소환 달린 공간들의 범주라 하자. 아핀 스킴이란 $LRS$에서 $(\Spec(A), {\mathcal O}_{\Spec(A)})$ 꼴의 쌍과 동형인 $LRS$의 대상을 말한다. 스킴이란 $LRS$의 대상 $(X, {\mathcal O}_X)$으로서, 모든 점 $x\in X$가 쌍 $(U, {\mathcal O}_X|_U)$이 아핀 스킴이 되는 열린 근방 $U \subset X$를 갖는 것을 말한다.

\begin{exercise} \label{exercises-exercise-one-point} 스킴이 아닌 $1$ 점의 국소환을 가진 공간을 찾아라. \end{exercise}

\begin{exercise} \label{exercises-exercise-two-points} $X$을, 그 기저위상 공간이 2개의 점을 가진 스킴이라고 하자. $X$가 아핀 스킴임을 보여라. \end{exercise}

\begin{exercise} \label{exercises-exercise-discrete-finite-set-points} $X$을, 그 기저위상 공간이 유한 이산 집합인 스킴이라고 하자. $X$가 아핀 스킴임을 보이라. \end{exercise}

\begin{exercise} \label{exercises-exercise-three-points} 세 점을 가지는 비아핀 스킴이 존재함을 보이라. \end{exercise}

\begin{exercise} \label{exercises-exercise-quasi-compact-closed-point} $X$을 공집합이 아닌 준콤팩트 스킴이라고 하자. $X$가 닫힌 점을 가짐을 보이라. \end{exercise}

\begin{remark} \label{exercises-remark-open-immersion} $(X, {\mathcal O}_X)$는 환 구조를 가진 공간이고, $U \subset X$가 열린 부분집합이라면, $(U, {\mathcal O}_X|_U)$은 환 구조를 가진 공간이다. 표기법: ${\mathcal O}_U = {\mathcal O}_X|_U$. % P12-ERR-0447: source singular/plural grammar.
환 구조를 가진 공간의 표준 사상 $$
j : (U, {\mathcal O}_U) \longrightarrow (X, {\mathcal O}_X).
$$가 존재한다. $(X, {\mathcal O}_X)$가 국소 환 구조를 가진 공간이라면, $(U, {\mathcal O}_U)$도 그렇고, $j$는 국소 환 구조를 가진 공간의 사상이다. $(X, {\mathcal O}_X)$가 스킴이라면, $(U, {\mathcal O}_U)$도 그렇고, $j$는 스킴의 사상이다. $(U, {\mathcal O}_U)$는 $(X, {\mathcal O}_X)$의 {\it 열린 부분 스킴}이라고 하며, $j$는 {\it 열린 포함}이라고 한다. 보다 일반적으로, 위와 같은 사상 $j : (U, {\mathcal O}_U) \to (X, {\mathcal O}_X)$과 {\it 동형}인 임의의 사상 $j' : (V, {\mathcal O}_V) \to (X, {\mathcal O}_X)$을 열린 포함이라고 부른다. \end{remark}

\begin{exercise} \label{exercises-exercise-open-affine-not-affine} % P12-ERR-0448: preserve the frozen missing restriction subscript.
아핀 스킴 $(X, {\mathcal O}_X)$과 열린 집합 $U \subset X$에 대해, $(U, {\mathcal O}_X|U)$이 아핀 스킴이 아닌 경우의 예를 제시하라. \end{exercise}

\begin{exercise} \label{exercises-exercise-morphism-does-not-extend} % P12-ERR-0449: source imperative grammar.
아핀 스킴 쌍 $(X, {\mathcal O}_X)$,$(Y, {\mathcal O}_Y)$, $X$의 열린 부분 스킴 $(U, {\mathcal O}_X|_U)$, 및 스킴의 사상 $(U, {\mathcal O}_X|_U) \to (Y, {\mathcal O}_Y)$ 이며, 스킴의 사상 $(X, {\mathcal O}_X) \to (Y, {\mathcal O}_Y)$에 연장할 수 없는 사례를 제시하시오. \end{exercise}

\begin{exercise} \label{exercises-exercise-closed-subscheme-does-not-extend} (이것은 꽤 어렵다.) % P12-ERR-0450: source imperative/article grammar.
스킴 $X$, 열린 부분 스킴 $U \subset X$, 및 닫힌 부분 스킴 $Z \subset U$ 이며, $Z$가 $X$의 닫힌 부분 스킴으로 연장될 수 없는 사례를 제시하시오. \end{exercise}

\begin{exercise} \label{exercises-exercise-not-morphism-schemes} 스킴 $X$, 체 $K$, 및 환이 있는 공간의 사상 $\Spec(K) \to X$ 중 스킴의 사상이 아닌 것의 예를 들어라. \end{exercise}

\begin{exercise} \label{exercises-exercise-just-kidding} \cite[Chapter II]{H} 1, 2장의 연습 문제를 모두 풀어라……\ \  농담이다! \end{exercise}

\begin{definition} \label{exercises-definition-integral} 스킴 $X$가 {\it 정}이라는 것은, $X$가 공집합이 아니고 공집합이 아닌 모든 아핀 열린집합 $U \subset X$에 대해 환 $\Gamma(U, \mathcal{O}_X) = \mathcal{O}_X(U)$가 정역인 것을 말한다. \end{definition}

\begin{exercise} \label{exercises-exercise-morphism-integral-schemes-surjective-stalks-not-closed} {\it 정} 스킴의 사상 $f : X \to Y$ 중에서, 유도되는 함수 ${\mathcal O}_{Y, f(x)}
\to {\mathcal O}_{X, x}$가 모든 $x\in X$에 대해 전사이지만, $f$가 닫힌 포함이 아닌 예를 들어라. \end{exercise}

\begin{exercise} \label{exercises-exercise-fibre-product-affines-not-affine} 섬유곱 $X \times_S Y$에서 $X$와 $Y$는 아핀이지만 $X \times_S Y$는 아핀이 아닌 예를 들어라. \end{exercise}

\begin{remark} \label{exercises-remark-separated-base-fibre-product-affines-affine} $S$가 분리 스킴이면 이런 일은 일어날 수 없음이 밝혀진다. 그 이유를 아는가? \end{remark}

\begin{exercise} \label{exercises-exercise-not-geometrically-integral} % P12-ERR-0451: source article/compound-modifier grammar.
${\mathbf Q}$ 위의 유한형인 1차원 정 스킴 $V$로서, $\Spec({\mathbf C}) \times_{\Spec({\mathbf Q})} V$가 정 스킴이 아닌 것의 예를 들어라. \end{exercise}

\begin{exercise} \label{exercises-exercise-not-geometrically-reduced} % P12-ERR-0452: source article/compound-modifier grammar.
체 $k$ 위의 유한형인 1차원 정 스킴 $V$로서, 어떤 유한 체확장 $k'/k$에 대해 $\Spec(k') \times_{\Spec(k)} V$가 축소 스킴이 아닌 것의 예를 들어라. \end{exercise}

\begin{remark} \label{exercises-remark-affine-dimension} 스킴이 아핀이라면, 차원은 대응하는 환의 Krull 차원과 같다. 따라서 지난 학기 결과를 차원 계산에 사용할 수 있다. % P12-ERR-0453: source possessive grammar.
\end{remark}


% Frozen source: exercises.tex lines 2970--3157.
\section{사상} \label{exercises-section-morphisms}

\noindent 사상 $\pi : X \to S$가 주어졌을 때 이 사상이 단면 또는 유리 단면을 갖는지는 중요한 문제이다. 다음은 몇 가지 예시 연습문제이다.

\begin{exercise} \label{exercises-exercise-no-section} 스킴의 사상 $$
\pi :
X = \Spec(\mathbf{C}[x, t, 1/xt])
\longrightarrow
S = \Spec(\mathbf{C}[t]).
$$을 생각하자. \begin{enumerate} \item $\pi \circ \sigma = \text{id}_S$ 되는 사상 $\sigma : S \to X$은 존재하지 않음을 보이시오. \item 공집합이 아닌 열린 집합 $U \subset S$과 사상 $\sigma : U \to X$이 있어서, $\pi \circ \sigma = \text{id}_U$가 되는 것이 존재함을 보이라. \end{enumerate} \end{exercise}

\begin{exercise} \label{exercises-exercise-no-rational-section} 스킴의 사상 $$
\pi :
X = \Spec(\mathbf{C}[x, t]/(x^2 + t))
\longrightarrow
S = \Spec(\mathbf{C}[t]).
$$을 고려한다. 공집합이 아닌 열린집합 $U \subset S$과 사상 $\sigma : U \to X$이 있어서, $\pi \circ \sigma = \text{id}_U$가 되는 것은 존재하지 않음을 보이라. \end{exercise}

\begin{exercise} \label{exercises-exercise-has-rational-section} $A, B, C \in \mathbf{C}[t]$를 0이 아닌 다항식으로 하자. 스킴의 사상 $$
\pi :
X = \Spec(\mathbf{C}[x, y, t]/(A + Bx^2 + Cy^2))
\longrightarrow
S = \Spec(\mathbf{C}[t]).
$$을 생각하라. 비어 있지 않은 열린 집합 $U \subset S$과 사상 $\sigma : U \to X$에 대하여, $\pi \circ \sigma = \text{id}_U$가 되는 것이 존재함을 보이라. (힌트: 형식적으로, 어떤 $d$에 대해 차수 $\leq d$의 $X, Y, Z \in \mathbf{C}[t]$를 사용하여 $x = X/Z$, $y = Y/Z$로 쓴다, 이것이 방정식의 해가 되는 조건을 구하시오. % P12-ERR-0454: preserve literal >> in the frozen formula.
다음으로 차원론을 사용하여, $d >> 0$ 라면 방정식을 만족하는 0이 아닌 $X, Y, Z$를 찾을 수 있음을 보이시오. ) \end{exercise}

\begin{remark} \label{exercises-remark-tsen} 연습문제 \ref{exercises-exercise-has-rational-section}은 ``Tsen의 정리''의 특별한 경우이다. 연습문제 \ref{exercises-exercise-no-section-curve}는 이 방법이 낮은 차수의 방정식에 한정됨을 보여준다(밑과 올의 차원이 1일 때에는 이차 곡선까지이다). \end{remark}

\begin{exercise} \label{exercises-exercise-no-section-curve} 스킴의 사상 $$
\pi :
X = \Spec(\mathbf{C}[x, y, t]
/(1 + t x^3 + t^2 y^3))
\longrightarrow
S = \Spec(\mathbf{C}[t])
$$을 생각한다. 공집합이 아닌 열린 집합 $U \subset S$과 사상 $\sigma : U \to X$에 대해, $\pi \circ \sigma = \text{id}_U$가 되는 것이 존재하지 않음을 보이라. \end{exercise}

\begin{exercise} \label{exercises-exercise-no-section-surface} 스킴 $$
X = \Spec(\mathbf{C}[\{x_i\}_{i = 1}^{8}, s, t]
/(1 + s x_1^3 + s^2 x_2^3 +
t x_3^3 + st x_4^3 + s^2t x_5^3 +
t^2 x_6^3 + st^2 x_7^3 + s^2t^2 x_8^3))
$$ 및 $$
S = \Spec(\mathbf{C}[s, t])
$$ 그리고 스킴의 사상 $$
\pi : X \longrightarrow S
$$을 생각한다. 공집합이 아닌 열린 집합 $U \subset S$과 사상 $\sigma : U \to X$이 있어, $\pi \circ \sigma = \text{id}_U$이 되는 것이 존재하지 않음을 증명하시오. \end{exercise}

\begin{exercise} \label{exercises-exercise-for-number-theorists} (수론을 배우는 사람을 위한. ) 닫힌 부분 스킴 $$
Z \subset \Spec\left({\mathbf Z}[x, \frac{1 }{ x(x-1)(2x-1)}]\right)
$$으로서, 사상 $Z \to \Spec({\mathbf Z})$이 유한하고 전단사인 경우의 예를 들어라. \end{exercise}

\begin{exercise} \label{exercises-exercise-quasi-section} 수론을 선호하지 않는다면, 다음을 보는 변형에 도전해도 좋다. $$
\Spec\left({\mathbf F}_p[t, x, \frac{1 }{ x(x-t)(tx-1)}]\right)
\longrightarrow
\Spec({\mathbf F}_p[t])
$$ % P12-ERR-0455: source adverbial grammar.
위의 스킴의 닫힌 부분 스킴 중에서, 아래 스킴으로 유한 전사 사상으로 사상되는 것을 찾아라. (여기서 기초체가 유한하다는 것에는 이론적인 이유가 있다. 다만, 이 특별한 경우에는 필요하지 않을 수도 있다. ) \end{exercise}

\begin{remark} \label{exercises-remark-interpretation-skolem-noether} 연습문제 \ref{exercises-exercise-for-number-theorists}과 \ref{exercises-exercise-quasi-section}에서는 사상 $f : X \to S$가 주어져 있다. 여기서 $S$는 Dedekind 환 $A$의 스펙트럼이고, $f$는 평탄하고 전사이며 일반 올은 기하학적으로 기약이다. 두 경우 모두 유한 전사 사상 $S' \to S$가 존재하여 밑변환 $X_{S'} \to S'$이 단면을 갖게 된다는 결론을 얻는다. $A$가 excellent이고 극대 아이디얼에서의 잉여체들이 유한체의 대수적 확대이며, $A'$이 $A$의 분수체의 유한 확대 안에서 그 정수적 폐포일 때 $\Pic(A')$이 비틀림군이면 이 결론이 성립함이 밝혀진다. \cite{MB1}을 참조하라. 예를 들어 $A = \mathbf{Z}$ 또는 $A = \mathbf{F}_p[x]$이거나 이들의 유한 확대인 경우이다. 그러나 극대 아이디얼에서의 잉여체들이 유한하면서 Picard 군이 비비틀림 원소를 갖는 Dedekind 환 $A$가 존재함이 알려져 있다. \cite{Goldman}을 참조하라. 그러한 환의 스펙트럼에 대해서는 이 결론이 거짓이다. 연습문제 \ref{exercises-exercise-no-quasi-section}은 이 결론이 성립하지 않는 기하학적 예를 준다. \end{remark}

\begin{exercise} \label{exercises-exercise-no-quasi-section} % P12-ERR-0456: source existential-clause grammar.
어떤 $\alpha \in \mathbf{C}$에 대해서도 $t - \alpha$로 나누어지지 않는 $f \in \mathbf{C}[x, t]$이며, $\Spec(\mathbf{C}[t])$에 유한 전사로 사상되는 $Z \subset \Spec(\mathbf{C}[x, t, 1/f])$가 존재하지 않는 것이 있음을 증명하라. ($f(x, t) = (xt - 2)(x - t + 3)$에서는 잘 될 것 같다. 어떤 후보에 대해서도, 원하는 성질을 가진다는 것을 보이는 것은 그리 쉽지 않을 것 같다. ) \end{exercise}

\begin{exercise} \label{exercises-exercise-finite} $A \to B$를 유한형 환 준동형이라 하자. $\Spec(B) \to \Spec(A)$가 어떤 $n$에 대한 닫힌 몰입 $\Spec(B) \to \mathbf{P}^n_A$을 통해 분해된다고 가정하자. $A \to B$가 유한 환 준동형임을, 즉 $B$가 $A$-가군으로서 유한함을 증명하라. 힌트: $A$가 뇌터 환이라면(그렇게 가정해도 좋다), 모든 닫힌 부분스킴 $Z \subset \mathbf{P}^n_A$와 $i \in \mathbf{Z}$에 대해 $H^i(Z, \mathcal{O}_Z)$가 유한 $A$-가군이라는 사실을 이용하여 논증할 수 있다. \end{exercise}

\begin{exercise} \label{exercises-exercise-projective-finite} $k$를 대수적으로 닫힌 체로 한다. $f : X \to Y$를 사영 다양체의 사상으로 하고, 모든 닫힌 점 $y \in Y$에 대해 $f^{-1}(\{y\})$가 유한하다고 하자. $f$가 스킴의 사상으로서 유한함을 증명하시오. 힌트: (a) 유한함은 국소적인 성질이다. (b) 연습 문제 \ref{exercises-exercise-finite} 로 환원해 보라. (c) 폐포함 $X \to \mathbf{P}^n_k$을 이용하여, $Y$ 위의 폐포함 $X \to \mathbf{P}^n_Y$을 얻어라. \end{exercise}

% Frozen source: exercises.tex lines 3158--3326.
\section{접공간} \label{exercises-section-tangent-space}

\begin{definition} \label{exercises-definition-dual-numbers} 임의의 환 $R$에 대해 {\it 쌍대수}의 환을 $R[\epsilon]$로 나타낸다. 이는 $R$-가군으로서 기저 $1$, $\epsilon$을 갖는 자유 가군이다. 환 구조는 $\epsilon^2 = 0$으로 정의된다. \end{definition}

\begin{exercise} \label{exercises-exercise-tangent-space-Zariski} $f : X \to S$를 스킴의 사상으로 한다. $x \in X$를 점으로 하고, $s = f(x)$로 한다. 실선의 가환 도식 $$
\xymatrix{
\Spec(\kappa(x)) \ar[r] \ar[dr] \ar@/^1pc/[rr] &
\Spec(\kappa(x)[\epsilon]) \ar@{.>}[r] \ar[d]&
X \ar[d] \\
&
\Spec(\kappa(s)) \ar[r] &
S
}
$$를 생각한다. 여기서 굽은 화살표는 $\Spec(\kappa(x))$에서 $X$ 로의 표준적인 사상이다. % P12-ERR-0457: source subject-verb agreement.
$\kappa(x) = \kappa(s)$ 라면, 도식을 가환하게 하는 점선 사상의 집합이 선형 사상의 집합 $$
\Hom_{\kappa(x)}(
\frac{\mathfrak m_x}{\mathfrak m_x^2 + \mathfrak m_s\mathcal{O}_{X, x}},
\kappa(x))
$$와 일대일로 대응함을 보이라. 다시 말하면, 그러한 전단사를 기술하라. (보다 일반적으로, $\kappa(x) \supset \kappa(s)$가 분리대수 확대인 경우에도 성립한다.) \end{exercise}

\begin{definition} \label{exercises-definition-tangent-space} $f : X \to S$를 스킴의 사상이라 하자. $x \in X$라 하자. 연습문제 \ref{exercises-exercise-tangent-space-Zariski}의 점선 사상들의 집합을 {\it $S$ 위에서 $X$의 접공간}이라 부르고 $T_{X/S, x}$로 나타낸다. 이 공간의 원소를 $X/S$의 $x$에서의 {\it 접벡터}라 부른다. \end{definition}

\begin{exercise} \label{exercises-exercise-simple-push-out} 임의의 체 $K$에 대해, 도식 $$
\xymatrix{
\Spec(K) \ar[r] \ar[d] & \Spec(K[\epsilon_1]) \ar[d] \\
\Spec(K[\epsilon_2]) \ar[r] &
\Spec(K[\epsilon_1, \epsilon_2]/(\epsilon_1\epsilon_2))
}
$$이 스킴 범주에서의 pushout 도식임을 증명하시오. (여기서도 이전과 마찬가지로 $\epsilon_i^2 = 0$이다.) \end{exercise}

\begin{exercise} \label{exercises-exercise-tangent-space-vectors-space} $f : X \to S$을 스킴의 사상으로 하자. $x \in X$로 하자. 연습문제 \ref{exercises-exercise-simple-push-out} 과 적절한 사상 $$
\Spec(K[\epsilon])
\longrightarrow
\Spec(K[\epsilon_1, \epsilon_2]/(\epsilon_1\epsilon_2)).
$$을 이용하여, 접벡터의 덧셈을 정의하시오. 마찬가지로 접벡터의 스칼라 배를 정의하라(이쪽이 더 쉽다). 이러한 구성으로 $T_{X/S, x}$가 $\kappa(x)$-벡터 공간이 됨을 보여라. \end{exercise}

\begin{exercise} \label{exercises-exercise-compute-TS} $k$를 체라 하자. 구조 사상 $f : X = \mathbf{A}^1_k \to \Spec(k) = S$를 생각하라. \begin{enumerate} \item $x \in X$를 닫힌 점이라 하자. $T_{X/S, x}$의 차원은 얼마인가? \item $\eta \in X$를 일반점이라 하자. $T_{X/S, \eta}$의 차원은 얼마인가? \item 이제 $X$를 $\Spec(\mathbf{Z})$ 위의 스킴으로 보자. $T_{X/\mathbf{Z}, x}$와 $T_{X/\mathbf{Z}, \eta}$의 차원은 얼마인가? \end{enumerate} \end{exercise}

\begin{remark} \label{exercises-remark-tangent-space-relative} 연습 \ref{exercises-exercise-compute-TS} 은, 좋은 개념을 얻기 위해 $S$ 위의 $X$의 접공간을 고려할 필요가 있는 이유를 설명하고 있다. \end{remark}

\begin{exercise} \label{exercises-exercise-compute-TS-field} 스킴의 사상 $$
f : X = \Spec(\mathbf{F}_p(t))
\longrightarrow
\Spec(\mathbf{F}_p(t^p)) = S
$$을 생각하라. $X$의 유일한 점에서 $X/S$의 접공간을 계산하라. 이상하지 않은가. $\mathbf{F}_p[t^p] \to \mathbf{F}_p[t]$에 대응하는 스킴의 사상을 취하면 어떤 일이 일어날 것 같는가. \end{exercise}

\begin{exercise} \label{exercises-exercise-compute-TS-cusp} $k$를 체로 한다. $X = \Spec(k[x, y]/(x^2 - y^3))$였을 때, 점 $x = (0, 0)$에서의 $X/k$의 접공간을 계산하라. \end{exercise}

\begin{exercise} \label{exercises-exercise-map-tangent-spaces} $f : X \to Y$을 $S$ 위의 스킴 사상으로 하자. $x \in X$를 점으로 하여, $y = f(x)$라 하자. 자연스러운 사상 $\kappa(y) \to \kappa(x)$가 전단사라고 가정하자. 정의를 이용하여, $f$가 자연스러운 선형 사상 $$
\text{d}f : T_{X/S, x} \longrightarrow T_{Y/S, y}.
$$를 유도함을 보이라. $\kappa(x) = \kappa(s)$의 경우에, 연습 \ref{exercises-exercise-tangent-space-Zariski}을 통해, 이것이 국소환 위에서 일어나는 것과 대응함을 확인하라. \end{exercise}

\begin{exercise} \label{exercises-exercise-Jacobian} $k$를 대수적으로 닫힌 체라고 하자. % P12-ERR-0458: preserve the frozen single-symbol coordinate arguments.
\begin{eqnarray*}
f : \mathbf{A}_k^n & \longrightarrow & \mathbf{A}^m_k \\
(x_1, \ldots, x_n) & \longmapsto & (f_1(x_i), \ldots, f_m(x_i))
\end{eqnarray*}을 $k$ 위의 스킴 사상이라고 하자. 이것은 $n$ 개 변수의 $m$ 개 다항식 $f_1, \ldots, f_m$으로 주어진다. % P12-ERR-0459: preserve the frozen Jacobian indexing/convention.
행렬 $$
A = \left( \frac{\partial f_j}{\partial x_i} \right)
$$을 생각하라. $x \in \mathbf{A}^n_k$을 고정점으로 한다. $y =  f(x)$라고 한다. 접공간의 사상 $T_{\mathbf{A}^n_k/k, x} \to T_{\mathbf{A}^m_k/k, y}$은 점 $x$에서의 행렬 $A$의 값으로 주어짐을 보여라. \end{exercise}


% Frozen source: exercises.tex lines 3327--3443.
\section{준연접 층} \label{exercises-section-quasi-coherent}

\begin{definition} \label{exercises-definition-quasi-coherent} $X$을 스킴으로 한다. $\mathcal{O}_X$-가군의 층 $\mathcal{F}$이 {\it 준연접}이라는 것은, 임의의 아핀 열린 집합 $\Spec(R) = U \subset X$에 대해, 제한 $\mathcal{F}|_U$가 어떤 $R$-가군 $M$을 사용하여 $\widetilde M$의 형태가 되는 것을 말한다. \end{definition}

\noindent % P12-ERR-0460: source demonstrative-number agreement.
이 조건은 $X$이 있는 아핀 열림 덮개의 각 요소에 대해 조사하면 충분하다. 더 많은 결과에 대해서는 스킴, 절 \ref{schemes-section-quasi-coherent}를 참조하라.

\begin{definition} \label{exercises-definition-specialization} $X$를 위상 공간으로 한다.$x, x' \in X$로 한다.$x \in \overline{\{x'\}}$일 때, 또한 그때에 한하여, $x$는 $x'$의 {\it 특수화}라고 한다. \end{definition}

\begin{exercise} \label{exercises-exercise-quasi-coherent-specialization-points} $X$를 스킴으로 한다.$x, x' \in X$라고 한다. $\mathcal{F}$를 $\mathcal{O}_X$-가군의 준연접 층으로 한다. (a) $x$은 $x'$의 특수화이며, (b) $\mathcal{F}_{x'} \not = 0$라고 가정한다. $\mathcal{F}_x \not = 0$를 증명하라. \end{exercise}

\begin{exercise} \label{exercises-exercise-O-module-specialization-points} 스킴 $X$, 점 $x, x' \in X$, 그리고 $\mathcal{O}_X$-가군의 층 $\mathcal{F}$으로서, (a) $x$는 $x'$의 특수화이며, (b) $\mathcal{F}_{x'} \not = 0$ 그리고 $\mathcal{F}_x = 0$인 것의 예를 찾아라. \end{exercise}

\begin{definition} \label{exercises-definition-Noetherian-scheme} 스킴 $X$이 {\it 국소 뇌터} 이라는 것은, % P12-ERR-0544: preserve the missing explicit condition x in U.
각 점 $x \in X$에 대해 아핀 열린 집합 $\Spec(R) = U \subset X$이 존재하고, $R$가 뇌터 환이 되는 것을 말한다. 스킴이 국소 뇌터 이고 준콤팩트일 때, {\it 뇌터} 라고 한다. \end{definition}

\noindent $X$가 국소 뇌터라면, $X$의 임의의 아핀 열린 집합은 뇌터 환의 스펙트럼이다. 성질, 보조정리 \ref{properties-lemma-locally-Noetherian} 를 참조.

\begin{definition} \label{exercises-definition-coherent} $X$를 국소 뇌터 스킴으로 하자. $\mathcal{F}$를 $\mathcal{O}_X$-가군의 준연접 층으로 하자. $\mathcal{F}$가 {\it 연접}이라는 것은, % P12-ERR-0544: preserve the missing explicit condition x in U.
각 점 $x \in X$에 대해 아핀 열린 집합 $\Spec(R) = U \subset X$가 존재하고, $\mathcal{F}|_U$이 어떤 유한 $R$-가군 $M$을 사용하여 $\widetilde M$와 동형이 됨을 의미한다. \end{definition}

\begin{exercise} \label{exercises-exercise-extend-quasi-coherent} $X = \Spec(R)$를 아핀 스킴이라고 하자. \begin{enumerate} \item $f \in R$라고 하자. $\mathcal{G}$를 열린 부분 스킴 $D(f)$ 위의 $\mathcal{O}_{D(f)}$-가군의 준연접 층으로 한다. 어떤 $\mathcal{O}_X$-가군의 준연접 층 $\mathcal{F}$에 대하여 $\mathcal{G} = \mathcal{F}|_{D(f)}$가 됨을 보이시오. \item $I \subset R$를 아이디얼로 한다. $i : Z \to X$를 $I$에 대응하는 $X$의 닫힌 부분 스킴으로 한다. $\mathcal{G}$를 닫힌 부분 스킴 $Z$ 위의 $\mathcal{O}_Z$-가군의 준연접 층으로 한다. 어떤 $\mathcal{O}_X$-가군의 준연접 층 $\mathcal{F}$에 대하여 $\mathcal{G} = i^*\mathcal{F}$가 됨을 보여라. (왜 이것이 사소한 문제인지.) \item $R$가 뇌터 환이라고 가정한다. $f \in R$라고 하자. $\mathcal{G}$를 열린 부분 스킴 $D(f)$ 위의 $\mathcal{O}_{D(f)}$-가군의 연접층이라고 하자. 어떤 $\mathcal{O}_X$-가군의 연접층 $\mathcal{F}$에 대하여 $\mathcal{G} = \mathcal{F}|_{D(f)}$가 됨을 보여라. \end{enumerate} \end{exercise}

\begin{remark} \label{exercises-remark-extend-off-open} $U \to X$가 준콤팩트 몰입이면 $U$ 위의 임의의 준연접 층은 $X$ 위의 준연접 층의 제한이다. $X$가 뇌터 스킴이고 $U \subset X$가 열려 있으면 $U$ 위의 임의의 연접층은 $X$ 위의 연접층의 제한이다. 물론 위 연습문제는 더 쉬우므로 이러한 일반적인 사실들을 사용해서는 안 된다. \end{remark}

% Frozen source: exercises.tex lines 3444--3598.
\section{Proj와 사영 스킴} \label{exercises-section-proj}

\begin{exercise} \label{exercises-exercise-graded-ring-specified-result} 다음을 만족하는 등급환 $S$의 예들을 들어라. \begin{enumerate} \item $\text{Proj}(S)$는 아핀이고 공집합이 아니다. \item $\text{Proj}(S)$는 정 스킴이고 공집합이 아니지만, 어떤 $n\geq 0$와 어떤 환 $A$에 대해서도 ${\mathbf P}^n_A$와 동형이 아니다. \end{enumerate} \end{exercise}

\begin{exercise} \label{exercises-exercise-nonconstant-morphism-proj} $\Spec({\mathbf C})$ 위의 스킴의 상수 사상이 아닌 사상 ${\mathbf P}^1_{\mathbf C} \to {\mathbf P}^5_{\mathbf C}$의 예를 들어라. \end{exercise}

\begin{exercise} \label{exercises-exercise-isomorphism-P1-conic} 스킴의 동형 $$
{\mathbf P}^1_{\mathbf C} \to
\text{Proj}({\mathbf C}[X_0, X_1, X_2]/(X_0^2 + X_1^2 + X_2^2))
$$의 예를 제시하라. \end{exercise}

\begin{exercise} \label{exercises-exercise-family-special-fibre-different} 다음 성질을 갖는 스킴의 사상 $f : X \to {\mathbf A}^1_{\mathbf C} = \Spec({\mathbf C}[T])$의 예를 들어라. $t \in {\mathbf A}^1_{\mathbf C}$ 위에서 $f$의 스킴론적 올 $X_t$는 (a) $t$가 $0$이 아닌 닫힌 점이면 ${\mathbf P}^1_{\mathbf C}$와 동형이고, (b) $t = 0$이면 ${\mathbf P}^1_{\mathbf C}$와 동형이 아니다. $X_0$을 이 사상의 {\it 특수 올}이라 부르겠다. 이는 매우 많은 방법으로 할 수 있다. 다음 추가 조건들을 각각 만족하는 예를 (불가능한 경우를 제외하고) 찾아보라. \begin{enumerate} \item 특수 올이 사영이 되게 할 수 있는가? \item 특수 올이 기약이고 사영이 되게 할 수 있는가? \item 특수 올이 정 스킴이고 사영이 되게 할 수 있는가? \item 특수 올이 매끄럽고 사영이 되게 할 수 있는가? \item $f$를 평탄 사상이 되게 할 수 있는가? 이는 모든 아핀 열린집합 $\Spec(A) \subset X$에 대해 유도되는 환 준동형 $\mathbf{C}[t] \to A$가 평탄하다는 뜻이고, 이 경우에는 $t$에 관한 임의의 영이 아닌 다항식이 $A$에서 영인자가 아닌 원소라는 뜻이다. \item $f$가 평탄 사영 사상이 되게 할 수 있는가? \item $f$가 평탄 사영 사상이고 특수 올이 축소 스킴이 되게 할 수 있는가? \item $f$가 평탄 사영 사상이고 특수 올이 기약이 되게 할 수 있는가? \item $f$가 평탄 사영 사상이고 특수 올이 정 스킴이 되게 할 수 있는가? \end{enumerate} ${\mathbf P}^1_{\mathbf C}$을 ${\mathbf C}$ 위의 다른 다양체로 바꾸면 어떻게 될 것 같은가? (위 변형 중 어느 것을 묻느냐에 따라 매우 어려워질 수 있다.) \end{exercise}

\begin{exercise} \label{exercises-exercise-affine-onto-projective-space} $n \geq 1$를 임의의 양의 정수라고 하자. $X$가 아핀인 전사 $X \to {\mathbf P}^n_{\mathbf C}$의 예를 제시하시오. \end{exercise}

\begin{exercise} \label{exercises-exercise-morphism-proj} $\text{Proj}$의 사상.$R$과 $S$를 등급환으로 한다. 환준동형 $$
\psi : R \to S
$$과 정수 $e \geq 1$가 주어지고, 모든 $d \geq 0$에 대해 $\psi(R_d) \subset S_{de}$가 성립한다고 가정한다. (우리의 규약에서는, $e = 1$가 아닌 한 이것은 등급환의 준동형이 아니다. ) \begin{enumerate} \item 어느 원소 $\mathfrak p \in \text{Proj}(S)$에 대해, $\text{Proj}(R)$의 대응하는 점이 올바르게 정해지는가. 바꾸어 말하면, $\psi$가 연속 사상 $r_\psi : U \to \text{Proj}(R)$를 정의하도록 하는, 적절한 열린 집합 $U \subset \text{Proj}(S)$을 구하시오. \item $U \not = \text{Proj}(S)$이 되는 예를 제시하라. \item $U = \text{Proj}(S)$이 되는 예를 제시하라. \item (적어 쓰지 않아도 좋다.) 연속 사상 $U \to \text{Proj}(R)$에는 표준적으로 층의 사상이 부속되며, $r_\psi$가 스킴의 사상 $$
\text{Proj}(S) \supset U \longrightarrow \text{Proj}(R).
$$이 되는 것을 납득하라. \item $R = \bigoplus_{d \geq 0} S_{de}$ ($S_e$, $S_{2e}$ 등등을 차수 $1$, $2$ 등등으로 가지는 등급환으로)에서, $\psi$가 포함 사상일 경우, 이 사상에 대해 무엇을 말할 수 있는가. \end{enumerate} \end{exercise}

\noindent {\bf 표기법.} $R$을 위와 같은 등급환이라 하고 $n \geq 0$을 정수라 하자. $X = \text{Proj}(R)$로 두자. 이때 $X$ 위에는 유일한 준연접 ${\mathcal O}_X$-가군 ${\mathcal O}_X(n)$이 존재하여, 양의 차수인 모든 동차 원소 $f \in R$에 대해, % P12-ERR-0545: preserve the frozen untwisted O_X restriction.
${\mathcal O}_X |_{D_{+}(f)}$는 $R_{(f)} = (R_f)_0$-가군 $(R_f)_n$에 수반되는 준연접 층이다. % P12-ERR-0461: retain equality span, render the explanatory prose readably.
($ = $ $R_f = R[1/f]$에서 차수가 $n$인 동차 원소들.) \cite[page 116+]{H}을 참조하라. 자연스러운 사상 $$
{\mathcal O}_X(n_1) \otimes_{{\mathcal O}_X} {\mathcal O}_X(n_2)
\longrightarrow
{\mathcal O}_X(n_1 + n_2)
$$가 존재함에 주의하라.

\begin{exercise} \label{exercises-exercise-pathologies-proj} $\text{Proj}$ 에서의 병리. 위와 같은 $R$에서, 다음을 만족하는 것의 예를 제시하라. \begin{enumerate} \item ${\mathcal O}_X(1)$는 가역 ${\mathcal O}_X$-가군이 아니다. \item ${\mathcal O}_X(1)$는 가역이지만, 자연스러운 사상 ${\mathcal O}_X(1) \otimes_{{\mathcal O}_X} {\mathcal O}_X(1) \to
{\mathcal O}_X(2)$는 동형이 아니다. \end{enumerate} \end{exercise}

\begin{exercise} \label{exercises-exercise-finitely-many-points-in-affine} $S$를 등급환으로 하자. $X = \text{Proj}(S)$로 하자. $X$의 임의의 유한한 점 집합이 표준 아핀 열린 집합에 포함됨을 보이라. \end{exercise}

\begin{exercise} \label{exercises-exercise-prepare-glueing} $S$를 등급환으로 한다. $X = \text{Proj}(S)$로 한다. $Z, Z' \subset X$를 두 개의 닫힌 부분 스킴으로 한다. $\varphi : Z \to Z'$를 동형으로 한다. $Z \cap Z' = \emptyset$라고 가정한다. 임의의 $z \in Z$에 대해 아핀 열린 집합 $U \subset X$가 존재하고, $z \in U$, $\varphi(z) \in U$, 및 $\varphi(Z \cap U) = Z' \cap U$이 됨을 보여라. (힌트: 연습 \ref{exercises-exercise-finitely-many-points-in-affine}과 스킴, 보조정리 \ref{schemes-lemma-standard-open-two-affines}과 유사한 것을 사용하라.) \end{exercise}

% Frozen source: exercises.tex lines 3599--3682.
\section{사영 직선에서 나오는 사상} \label{exercises-section-from-P1}

\noindent 이 절에서는 $\mathbf{P}^1$에서 사영 스킴으로 가는 사상들을 살펴본다.

\begin{exercise} \label{exercises-exercise-from-generic-point} $k$를 체로 한다. $k[t] \subset k(t)$를 다항식 환에서 그 분수체로의 포함이라고 하자. $X$를 $k$ 위 유한형 스킴이라고 하자. $k$ 위의 임의의 사상 $$
\varphi : \Spec(k(t)) \longrightarrow X
$$에 대해, 0이 아닌 $f \in k[t]$과 $k$ 위의 사상 $\psi : \Spec(k[t, 1/f]) \to X$이 존재하고, $\varphi$가 합성 $$
\Spec(k(t)) \longrightarrow \Spec(k[t, 1/f]) \longrightarrow X
$$가 됨을 보이라. \end{exercise}

\begin{exercise} \label{exercises-exercise-from-generic-point-Pn} $k$를 체로 한다. $k[t] \subset k(t)$를 다항식 환에서 그 분수체로의 포함으로 한다. $k$ 위의 임의의 사상 $$
\varphi : \Spec(k(t)) \longrightarrow \mathbf{P}^n_k
$$에 대해, $k$ 위의 사상 $\psi : \Spec(k[t]) \to \mathbf{P}^n_k$가 존재하며, $\varphi$가 합성 $$
\Spec(k(t)) \longrightarrow \Spec(k[t]) \longrightarrow \mathbf{P}^n_k
$$가 됨을 증명하라. 힌트: $\varphi$의 상은 어떤 $i$에 대해 표준 열린집합 $D_+(T_i)$에 포함된다. 따라서 '분모를 없앨 수 있음'을 보이라. \end{exercise}

\begin{exercise} \label{exercises-exercise-from-generic-point-projective} $k$를 체로 하자. $k[t] \subset k(t)$를 다항식 환에서 그 분수체로의 포함이라고 하자. $X$를 $k$ 위의 사영 스킴이라고 하자. $k$ 위의 임의의 사상 $$
\varphi : \Spec(k(t)) \longrightarrow X
$$에 대해, $k$ 위의 사상 $\psi : \Spec(k[t]) \to X$가 존재하여, $\varphi$가 합성 $$
\Spec(k(t)) \longrightarrow \Spec(k[t]) \longrightarrow X
$$가 됨을 보이시오. 힌트: 연습 \ref{exercises-exercise-from-generic-point-Pn}를 사용하라. \end{exercise}

\begin{exercise} \label{exercises-exercise-from-generic-point-P1-to-projective} $k$를 체라고 하자. $X$을 $k$ 위의 사영 스킴으로 하자. $K$을 $\mathbf{P}^1_k$의 함수체로 하자(아래의 힌트를 참조). $k$ 위의 임의의 사상 $$
\varphi : \Spec(K) \longrightarrow X
$$에 대해, $k$ 위의 사상 $\psi : \mathbf{P}^1_k \to X$가 존재하여, $\varphi$가 합성 $$
\Spec(k(t)) \longrightarrow \mathbf{P}^1_k \longrightarrow X
$$가 됨을 보이라. 힌트: 아핀 열림 덮개 $\mathbf{P}^1_k = D_+(T_0) \cup D_+(T_1)$의 두 부분 각각에 대해, 연습 \ref{exercises-exercise-from-generic-point-projective}을 사용하라. $D_+(T_0)$는 다항식 환의 스펙트럼이며, $K$는 이 다항식 환의 분수체임을 이용하라. \end{exercise}


% Frozen source: exercises.tex lines 3683--3762.
\section{곡면에서 곡선으로의 사상} \label{exercises-section-from-surfaces-to-curves}

\begin{exercise} \label{exercises-exercise-points-projective-space} $R$를 환으로 한다. $R \to k$를 $R$에서 체로 가는 사상이라고 하자. $n \geq 0$라고 하자. $$
\Mor_{\Spec(R)}(\Spec(k), \mathbf{P}^n_R)
=
(k^{n + 1} \setminus \{0\})/k^*
$$를 증명하라. 여기서 $k^*$는 $k^{n + 1}$에 스칼라 곱으로 작용한다. 이후, 원소 $(x_0, \ldots, x_n) \in k^{n + 1} \setminus \{0\}$의 동치류에 대응하는 사상 $\Spec(k) \to \mathbf{P}^n_k$를 $(x_0 : \ldots : x_n)$로 나타낸다. \end{exercise}

\begin{exercise} \label{exercises-exercise-curve-projective-plane} $k$를 체라 하자. $Z \subset \mathbf{P}^2_k$를 기약이고 축소인 닫힌 부분스킴이라 하자. % P12-ERR-0462: source indefinite-article grammar.
(a) $Z$가 닫힌 점이거나, (b) 차수 $> 0$인 동차 기약 다항식 $F \in k[X_0, X_1, X_2]$가 존재하여 $Z = V_{+}(F)$이거나, (c) $Z = \mathbf{P}^2_k$임을 증명하라. (힌트: 표준 아핀 열린집합 위에서 생각하라.) \end{exercise}

\begin{exercise} \label{exercises-exercise-bezout} $k$를 체라 하자. $Z_1, Z_2 \subset \mathbf{P}^2_k$를, % P12-ERR-0463: preserve the frozen unindexed F and indexed F_i.
어떤 차수 $> 0$인 동차 기약 다항식 $F_i \in k[X_0, X_1, X_2]$에 대해 $V_{+}(F)$ 꼴인 기약 닫힌 부분스킴들이라 하자. $Z_1 \cap Z_2$가 공집합이 아님을 증명하라. (힌트: 차원론을 사용하여 $k[X_0, X_1, X_2]/(F_1, F_2)$의 $0$에서의 국소환의 차원을 평가하라.) \end{exercise}

\begin{exercise} \label{exercises-exercise-no-nonconstant-morphism-proj} $\Spec(\mathbf{C})$ 위에서 상수 사상이 아닌 스킴의 사상 $\mathbf{P}^2_{\mathbf{C}} \to \mathbf{P}^1_{\mathbf{C}}$은 존재하지 않음을 보여라. 여기서 {\it 상수 사상}이란 상이 한 점인 사상을 말한다. (힌트: 사상이 상수 사상이 아니면 $0$과 $\infty$ 위의 올들을 생각하고, 이들이 만나야 함을 보여 모순을 얻어라.) \end{exercise}

\begin{exercise} \label{exercises-exercise-nonconstant-morphism} $k$을 체로 하자. $X \subset \mathbf{P}^3_k$을 하나의 동차방정식 $F \in k[X_0, X_1, X_2, X_3]$으로 주어진 닫힌 부분 스킴으로 하자. 다시 말하면, 강의에서 설명한 것처럼 $$
X = \text{Proj}(k[X_0, X_1, X_2, X_3]/(F)) \subset \mathbf{P}^3_k
$$이다. 양의 차수를 갖는 어떤 동차 다항식 $G, H \in k[X_0, X_1, X_2, X_3]$에 대해 $$
F = X_0 G + X_1 H
$$가 된다고 가정하자. $X_0, X_1, G, H$가 공통 영점을 갖지 않는다면, $\Spec(k)$ 위에서 상수 사상이 아닌 스킴의 사상 $$
X \longrightarrow \mathbf{P}^1_k
$$ 중 체 값을 갖는 점들에서(연습문제 \ref{exercises-exercise-points-projective-space} 참조), $x_0$ 또는 $x_1$이 영이 아닐 때 $(x_0 : x_1 : x_2 : x_3) \mapsto (x_0 : x_1)$의 꼴인 것이 존재함을 보여라. \end{exercise}

% Frozen source: exercises.tex lines 3763--3925.
\section{가역층} \label{exercises-section-invertible-sheaves}

\begin{definition} \label{exercises-definition-invertible-sheaf} $X$을 국소환 달린 공간으로 하자. $X$ 위의 {\it 가역 ${\mathcal O}_X$-가군}이란, ${\mathcal O}_X$-가군의 층 ${\mathcal L}$으로서, 각 점이 열린 근방 $U \subset X$을 갖고, ${\mathcal L}|_U$이 ${\mathcal O}_U$-가군으로서 ${\mathcal O}_U$과 동형이 되는 것을 말한다. ${\mathcal L}$가 ${\mathcal O}_X$-가군으로서 ${\mathcal O}_X$와 동형일 때, 자명하다고 한다. \end{definition}

\begin{exercise} \label{exercises-exercise-general-facts-invertible} 일반적인 사실. \begin{enumerate} \item 스킴 $X$ 위의 가역 ${\mathcal O}_X$-가군이 준연접임을 보이시오. % P12-ERR-0464: source predicate omits "is".
\item $X\to Y$를 국소환이 달린 공간의 사상이라고 하고, ${\mathcal L}$를 가역 ${\mathcal O}_Y$-가군이라고 하자. % P12-ERR-0465: preserve the source's unhyphenated module phrase.
$f^\ast {\mathcal L}$가 가역 ${\mathcal O}_X$ 가군임을 증명하라. \end{enumerate} \end{exercise}

\begin{exercise} \label{exercises-exercise-invertible-algebra} 대수. \begin{enumerate} \item 아핀 스킴 $\Spec(A)$ 위의 가역 ${\mathcal O}_X$-가군은 (i) 유한, (ii) 사영적, (iii) 계수 1의 국소 자유이므로, 따라서 (iv) 평탄, (v) 유한 표현인 $A$-가군 $M$에 대응함을 증명하라. % P12-ERR-0466: source possessive grammar.
(지난 학기의 강의나 대수학 책의 결과는 자유롭게 인용해도 된다.) % P12-ERR-0467: preserve the frozen manual item reference (a).
\item $A$을 정역으로 하고, $M$을 (a)와 같은 가군으로 한다. $M$는 $A$-가군으로서, 임의의 소 아이디얼 ${\mathfrak p}$에 대해 $IA_{\mathfrak p}$가 주 아이디얼이 되도록 하는 아이디얼 $I \subset A$과 동형임을 보이라. \end{enumerate} \end{exercise}

\begin{definition} \label{exercises-definition-invertible-module} $R$을 환으로 한다. {\it 가역 가군 $M$}이란, $R$-가군 $M$으로서, $\widetilde M$가 $R$의 스펙트럼 위의 가역층인 것을 말한다. $R$-가군으로서 $M \cong R$일 때, $M$는 {\it 자명}이라고 한다. \end{definition}

\noindent 바꿔 말하면, $M$가 가역적이라는 것은 평탄, 유한 표현, 사영, 계수 1의 국소 자유 조건을 모두 만족하는 것과 동치이다. (물론, 계수 1의 국소 자유이면 충분하다.)

\begin{exercise} \label{exercises-exercise-simple-examples-invertible} 간단한 예. \begin{enumerate} \item \label{exercises-item-affine-line} $k$를 체로 한다.$A = k[x]$라고 하자. $X = \Spec(A)$ 위의 가역 ${\mathcal O}_X$-가군은 모두 자명함을 보이라. 다시 말하면, 임의의 가역 $A$-가군은 계수 1의 자유 가군임을 보이라. \item \label{exercises-item-affine-line-with-0-and-1-identified} $A$를 환 $$
A = \{ f\in k[x] \mid f(0) = f(1) \}.
$$라고 하자. $k = {\mathbf F}_2$가 아닐 경우, 비자명한 가역 $A$-가군이 존재함을 보이라. (힌트: $\Spec(A)$는 ${\mathbf A}^1_k = \Spec(k[x])$의 $0$와 $1$을 동일시한 것으로 생각하라. ) \item \label{exercises-item-affine-line-with-cusp} 환 $A = k[x^2, x^3] \subset k[x]$에 대해서 (\ref{exercises-item-affine-line-with-0-and-1-identified}) 와 같은 질문을 생각하라 (단, 이번에는 $k = {\mathbf F}_2$의 경우에도 성립한다). \end{enumerate} \end{exercise}

\begin{exercise} \label{exercises-exercise-higher-dimension-invertible} 고차원 경우. \begin{enumerate} \item 2차원 아핀 공간 위의 임의의 가역 층이 자명함을 증명하라. 더 정확히 말하면, $k$를 체로 하고, ${\mathbf A}^2_k = \Spec(k[x, y])$로 한다. ${\mathbf A}^2_k$ 위의 임의의 가역 층이 자명함을 증명하라. (힌트: 한 가지 방법은 대응하는 가군 $M$을 생각하고, $M \otimes_{k[x, y]} k(x)[y]$을 조사하며, 연습문제 \ref{exercises-exercise-simple-examples-invertible} (\ref{exercises-item-affine-line})를 통해 그 생성원을 찾는 것이다. 그 후에도 생각할 필요가 있다. % P12-ERR-0468: source duplicated infinitive marker.
다른 방법은 연습문제 \ref{exercises-exercise-invertible-algebra}와 다항식환의 아이디얼에 대해 알고 있는 것, 즉 높이 1의 소 아이디얼은 기약 다항식으로 생성됨을 이용하는 방법이다. 그 후에도 생각할 필요가 있다. ) \item 2차원 아핀 공간의 임의의 열린 부분 스킴 위의 가역층이 모두 자명함을 증명하라. 보다 정확히 말하면, $k$를 체로 하고, $U \subset {\mathbf A}^2_k$를 열린 부분 스킴으로 하자. $U$ 위의 임의의 가역층이 자명함을 보이자. 힌트: $U$ 위의 임의의 가역층이 ${\mathbf A}^2_k$ 위의 가역층으로 연장될 수 있음을 보이라. 쉽지 않지만, \cite{H} 에서 찾을 수 있다. \item 체 위에서 꼭짓점을 제외한 원뿔 위에서 비자명한 가역층의 예를 찾아라. 보다 정확하게는, $k$를 체로 하고, $C = \Spec(k[x, y, z]/(xy-z^2))$로 한다. $U = C \setminus \{ (x, y, z) \}$로 한다. $U$ 위의 비자명한 가역층을 찾아라. 힌트: 하나만 찾는 것보다, $U$ 위의 가역층의 동형류 군을 계산하는 것이 더 쉬울 수 있다. $U$가 "다루기 쉬운" 열린 집합 $\Spec(k[x, y, z, 1/x]/(xy-z^2))$와 $\Spec(k[x, y, z, 1/y]/(xy-z^2))$로 덮여 있다는 것에 주의하라. \end{enumerate} \end{exercise}

\begin{definition} \label{exercises-definition-picard-group} $X$을 국소환 달린 공간이라고 하자. {\it $X$의 Picard 군}은 가역 $\mathcal{O}_X$-가군의 동형류 집합 $\Pic(X)$에 텐서곱으로 주어지는 덧셈을 넣은 것이다. 가군, 정의 \ref{modules-definition-pic}를 참조하라. 환 $R$에 대해서는 $\Pic(R) = \Pic(\Spec(R))$라 하자. \end{definition}

\begin{exercise} \label{exercises-exercise-traverso} $R$을 환으로 한다. \begin{enumerate} \item $R$이 뇌터 정규 정역이라면 $\Pic(R) = \Pic(R[t])$임을 증명하라. [힌트: 사상 $R \to R[t]$의 왼쪽 역이 되는 사상 $R[t] \to R$, $t \mapsto 0$이 존재한다. 따라서, $M/tM \cong R$를 만족하는 임의의 가역 $R[t]$-가군 $M$이 계수 $1$의 자유 가군임을 증명하면 충분하다. $K$를 $R$의 분수체라고 하자. 자명화 $K[t] \to M \otimes_{R[t]} K[t]$를 선택하라. 이는 연습문제 \ref{exercises-exercise-simple-examples-invertible} (\ref{exercises-item-affine-line})에 의해 가능하다. 이를 조정하여 위의 $M/tM$의 자명화와 일치하도록 하라. 실제로 이것이 $R[t]$ 상에서의 $M$의 자명화임을 보이라 (여기서 정규성이 사용된다). \item $k$를 체로 한다. $\Pic(k[x^2, x^3, t]) \not = \Pic(k[x^2, x^3])$를 보이라. \end{enumerate} \end{exercise}


\section{{\v C}ech 코호몰로지} \label{exercises-section-cech-cohomology}

\begin{exercise} \label{exercises-exercise-cech-cohomology} {\v C}ech 코호몰로지. 여기서 $k$는 체로 한다. \begin{enumerate} \item $X$를 열린 덮개 ${\mathcal U} : X = U_1 \cup U_2$을 가지는 스킴으로 한다. 여기서 $U_1 = \Spec(k[x])$, $U_2 =  \Spec(k[y])$, $U_1 \cap U_2 = \Spec(k[z, 1/z])$ 이며, 열린 몰입 $U_1 \cap U_2 \to U_1$ 및 $U_1 \cap U_2 \to U_2$는 각각 $x \mapsto z$ 및 $y \mapsto z$에 의해 정해진다 (여기는 문자 그대로의 의미이다). (강의에서 이러한 $X$가 존재함을 보았다. 이것은 원점을 두 배로 한 아핀 선이다.) ${\check H}^1({\mathcal U}, {\mathcal
O})$를 계산하라. 예를 들어, $k$-벡터 공간으로서의 기저를 제시하라. \item ${\check H}^1({\mathcal U}, {\mathcal O})$의 각 원소에 대하여, ${\mathcal O}_X$-가군의 층의 완전열 $$
0 \to {\mathcal O}_X \to E \to {\mathcal O}_X \to 0
$$로서, 경계 $\delta(1) \in {\check H}^1({\mathcal U},
{\mathcal O})$가 주어진 원소와 같아지도록 하는 것을 구성하라. (이 주장을 의미 있는 것으로 만드는 것도 문제의 일부이다. 아래도 참고하라. 이러한 것이 존재하지 않을 수 없음을 추상적으로 보여주어도 좋다. ) \end{enumerate} \end{exercise}

\begin{definition} \label{exercises-definition-delta} (delta의 정의.)임의의 위상 공간 $X$ 위의 아벨 군 층의 짧은 완전열 $$
0 \to {\mathcal F}_1 \to {\mathcal F}_2 \to {\mathcal F}_3 \to 0
$$이 주어졌다고 하자. 경계 사상 $\delta : H^0(X, {\mathcal F}_3) \to {\check H}^1(X, {\mathcal F}_1)$을 다음과 같이 정의한다. 원소 $\tau \in H^0(X, {\mathcal F}_3)$를 가져온다. 각 $i$에 대해, $\tau$의 $U_i$에 대한 제한을 올리는 단면 $\tilde \tau_i \in {\mathcal F}_2$이 존재하는 열린 덮개 ${\mathcal U} : X = \bigcup_{i\in I} U_i$를 선택한다. 이때, 대응 $$
(i_0, i_1) \longmapsto
\tilde \tau_{i_0}|_{U_{i_0i_1}} - \tilde \tau_{i_1}|_{U_{i_0i_1}}.
$$를 고려한다. 이것은 명백히 {\v C}ech 복체 ${\check C}^\ast({\mathcal U}, {\mathcal F}_2)$의 1-코바운더리이다. 그러나 (${\mathcal F}_1$를 ${\mathcal F}_2$의 부분층으로 간주하면), 우변은 항상 $U_{i_0i_1}$ 위의 ${\mathcal F}_1$의 단면임에 주의한다. % P12-ERR-0469: preserve the frozen F_2 coefficient and cochain wording.
따라서, 이 대응은 복합체 ${\check C}^\ast({\mathcal U}, {\mathcal F}_2)$의 1-코체인을 정한다. 이 1-코체인의 코호몰로지류를 정의에 따라 {\it $\delta(\tau)$} 라고 한다. \end{definition}


\section{코호몰로지} \label{exercises-section-cohomology}

\begin{exercise} \label{exercises-exercise-cohomology-not-zero} $X = \mathbf{R}$에 보통 유클리드 위상을 넣는다. 코호몰로지의 형식적 성질(함자성 및 긴 완전 코호몰로지 열)만을 사용하여, $H^1(X, \mathcal{F})$가 0이 아닌 $X$ 위의 층 $\mathcal{F}$가 존재함을 보여라. \end{exercise}

\begin{exercise} \label{exercises-exercise-mayer-vietoris} 위상 공간 $X = U \cup V$가 두 개의 열린집합의 합으로 쓰여 있다고 가정하자. 이때, 긴 완전 Mayer–Vietoris 열 $$
0 \to
H^0(X, \mathcal{F}) \to
H^0(U, \mathcal{F}) \oplus H^0(V, \mathcal{F}) \to
H^0(U \cap V, \mathcal{F}) \to
H^1(X, \mathcal{F}) \to \ldots
$$이 있다. Mayer–Vietoris 긴 완전열의 구성에서 본질적인 단사층의 성질은 무엇인가. 왜 그 성질이 성립하는가. \end{exercise}

\begin{exercise} \label{exercises-exercise-cohomology-two-point-space} $X$을 위상 공간으로 하자. \begin{enumerate} \item $X$의 점이 $2$ 개 이하이면, $i > 0$에 대하여 $H^i(X, \mathcal{F})$이 0임을 보이라. \item $X$ 점이 $3$ 개라면 어떻게 되는가. \end{enumerate} \end{exercise}

\begin{exercise} \label{exercises-exercise-cohomology-spec-local-ring} $X$를 국소 환의 스펙트럼으로 한다. 임의의 아벨 군 층 $\mathcal{F}$와 $i > 0$에 대해 $H^i(X, \mathcal{F})$이 0임을 보여라. \end{exercise}

\begin{exercise} \label{exercises-exercise-affine-morphism} $f : X \to Y$를 스킴의 아핀 사상으로 한다. 임의의 준연접 $\mathcal{O}_X$-가군 $\mathcal{F}$에 대해 $H^i(X, \mathcal{F}) = H^i(Y, f_*\mathcal{F})$임을 증명하라. 필요하다면 $X$와 $Y$에 추가적으로 몇 가지 조건을 부과하고, 분리 스킴의 아핀 열린 덮개와 준연접 층에 대해, {\v C}ech 코호몰로지와 코호몰로지가 일치함을 이용해도 좋다. \end{exercise}

\begin{exercise} \label{exercises-exercise-affine-morphism-to-Pn} $A$를 환으로 한다. $\mathbf{P}^n_A = \text{Proj}(A[T_0, \ldots, T_n])$를 $A$ 위의 사영 공간으로 한다. $\mathbf{A}^{n + 1}_A = \Spec(A[T_0, \ldots, T_n])$로 하고, $$
U = \bigcup\nolimits_{i = 0, \ldots, n} D(T_i) \subset \mathbf{A}^{n + 1}_A
$$을 닫힌 몰입 $0 : \Spec(A) \to \mathbf{A}^{n + 1}_A$의 상의 여집합으로 하자. 아핀 전사 $$
f : U \longrightarrow \mathbf{P}^n_A
$$를 구성하고, $f_*\mathcal{O}_U =
\bigoplus_{d \in \mathbf{Z}} \mathcal{O}_{\mathbf{P}^n_A}(d)$임을 증명하라. 더 일반적으로, 등급 $A[T_0, \ldots, T_n]$-가군 $M$에 대해 $$
f_*(\widetilde{M}|_U) =
\bigoplus\nolimits_{d \in \mathbf{Z}} \widetilde{M(d)}
$$임을 보이라. 여기서 좌변에서는 $\mathbf{A}^{n + 1}_A$ 위에서 $M$에 부속된 준연접 층 $\widetilde{M}$을 사용하고, 우변에서는 등급 가군 $M(d)$에 부속된 준연접 층 $\widetilde{M(d)}$을 사용한다. \end{exercise}

\begin{exercise} \label{exercises-exercise-compute-Pn} $A$을 환으로 하고, $\mathbf{P}^n_A = \text{Proj}(A[T_0, \ldots, T_n])$를 $A$ 위의 사영 공간으로 한다. $\mathbf{P}^n_A$ 위의 구조층의 Serre 꼬임 $\mathcal{O}_{\mathbf{P}^n_A}(d)$의 코호몰로지를 자세히 계산하라. 분리 스킴의 아핀 열린 덮개와 준연접 층에 대해, {\v C}ech 코호몰로지와 코호몰로지가 일치함, 또한 {\v C}ech 코호몰로지를 사용할 수 있다. \end{exercise}

\begin{exercise} \label{exercises-exercise-cohomology-hypersurface} $A$를 환으로 하고, $\mathbf{P}^n_A = \text{Proj}(A[T_0, \ldots, T_n])$를 $A$ 위의 사영 공간이라고 하자. $F \in A[T_0, \ldots, T_n]$를 차수 $d$의 동차식이라고 하자. $X \subset \mathbf{P}^n_A$를 $A[T_0, \ldots, T_n]$의 등급 아이디얼 $(F)$에 대응하는 닫힌 부분 스킴이라고 하자. $H^i(X, \mathcal{O}_X)$에 대해 무엇을 말할 수 있는가. \end{exercise}

\begin{exercise} \label{exercises-exercise-characterize-finite-product-fields} $R$를, 임의의 좌완전 함자 $F : \text{Mod}_R \to \textit{Ab}$에 대하여 $i > 0$라면 $R^iF = 0$가 되는 환으로 하자. $R$가 유한 개 체의 직접곱임을 보이라. \end{exercise}


\section{추가적인 코호몰로지} \label{exercises-section-more-cohomology}

\begin{exercise} \label{exercises-exercise-cohomology-coordinate-axes} $k$를 체라 하자. $X \subset \mathbf{P}^n_k$를 「좌표 십자」라고 한다.즉, $X$는 동차 방정식 $$
T_i T_j = 0\text{, 여기서 }i > j > 0
$$로 정의된 것으로 한다.여기서 평소와 같이 $\mathbf{P}^n_k = \text{Proj}(k[T_0, \ldots, T_n])$라고 쓴다. 다시 말하면, $X$는 다항식 환의 몫 $k[T_0, \ldots, T_n]/(T_iT_j; i > j > 0)$에 대응하는 닫힌 부분 스킴이다. 모든 $i$에 대해 $H^i(X, \mathcal{O}_X)$를 계산하시오. 힌트: {\v C}ech 코호몰로지를 사용하라. \end{exercise}

\begin{exercise} \label{exercises-exercise-compute-cohomology-punctured} $A$를 환으로 한다. $I = (f_1, \ldots, f_t)$를 $A$의 유한 생성 아이디얼로 한다. $U \subset \Spec(A)$를 $V(I)$의 여집합으로 한다. 임의의 $A$-가군 $M$에 대해, 그 코호몰로지 군이 $H^n(U, \widetilde{M})$를 주는 $A$-가군의 복체를 ($A$, $f_1, \ldots, f_t$, $M$를 사용하여) 적어라. \end{exercise}

\begin{exercise} \label{exercises-exercise-compute-cohomology-affine-space-punctured} $k$를 체로 한다. $U \subset \mathbf{A}^d_k$를 $\mathbf{A}^d_k$의 닫힌 점 $0$의 여집합으로 한다. 모든 $n$에 대해 $H^n(U, \mathcal{O}_U)$를 계산하라. \end{exercise}

\begin{exercise} \label{exercises-exercise-find-curve-genus-one-hundred} $k$을 체로 한다. $k$ 위 사영적이고 차원 $1$의 스킴 $X$으로서, $H^0(X, \mathcal{O}_X) = k$ 및 $\dim_k H^1(X, \mathcal{O}_X) = 100$를 만족하는 것을 구체적으로 찾아라. \end{exercise}

\begin{exercise} \label{exercises-exercise-degree-2-cover} $f : X \to Y$를 차수 $2$의 유한 국소 자유 사상이라고 하자. $X$와 $Y$는 정 스킴이며, $2$는 $Y$의 구조층에서 가역적, 즉 $2 \in \Gamma(Y, \mathcal{O}_Y)$는 가역적이라고 가정한다. $\mathcal{O}_Y$-가군의 사상 $$
f^\sharp : \mathcal{O}_Y \longrightarrow f_*\mathcal{O}_X
$$이 왼쪽 역을 갖는 것, 즉 $\tau \circ f^\sharp = \text{id}$를 만족하는 $\mathcal{O}_Y$-가군의 사상 $\tau : f_*\mathcal{O}_X \to \mathcal{O}_Y$이 존재함을 보여라. $H^n(Y, \mathcal{O}_Y) \to H^n(X, \mathcal{O}_X)$이 단사임을 이끌어라\footnote{차수 $2$인 정 스킴 사이의 유한 국소 자유 사상 $X \to Y$ 중, 사상 $H^1(Y, \mathcal{O}_Y) \to H^1(X, \mathcal{O}_X)$이 단사가 아닌 경우도 실제로 존재한다.}. \end{exercise}

\begin{exercise} \label{exercises-exercise-pic} $X$를 스킴(또는 국소 환을 가진 공간)이라고 하자. 대응 $U \mapsto \mathcal{O}_X(U)^*$는 $\mathcal{O}_X^*$로 쓰이는 군의 층을 정의한다. $X$의 Picard 군(정의 \ref{exercises-definition-picard-group})이 $H^1(X, \mathcal{O}_X^*)$와 같은 이유를 간단히 설명하라. \end{exercise}

\begin{exercise} \label{exercises-exercise-pic-nontrivial} $\Pic(X)$가 자명하지 않은 아핀 스킴 $X$의 예를 들어라. 연습 \ref{exercises-exercise-pic}를 사용하여, 그러한 임의의 $X$에 대해 $H^1(X, -)$가 영 사상이 아님을 이끌어라. \end{exercise}

\begin{exercise} \label{exercises-exercise-kill-cohomology-complement} $A$를 환으로 하자. $I = (f_1, \ldots, f_t)$를 $A$의 유한 생성 아이디얼로 하자. $U \subset \Spec(A)$을 $V(I)$의 여집합으로 한다. 준연접 $\mathcal{O}_{\Spec(A)}$-가군 $\mathcal{F}$과, $p > 0$을 만족하는 $\xi \in H^p(U, \mathcal{F})$가 주어졌을 때, 모든 $i = 1, \ldots, t$에 대하여 $f_i^n \xi = 0$가 되는 $n > 0$가 존재함을 증명하라. 힌트: 연습 \ref{exercises-exercise-compute-cohomology-punctured}에서 얻은 복체를 사용하는 것이 한 가지 방법이다. \end{exercise}

\begin{exercise} \label{exercises-exercise-h0-complement} $A$를 환으로 한다. $I = (f_1, \ldots, f_t)$를 $A$의 유한 생성 아이디얼이라고 하자. $U \subset \Spec(A)$를 $V(I)$의 여집합이라고 하자. $M$를 $I$-비틀림이 0인 $A$-가군이라고 하자, 즉 $0 = \Ker((f_1, \ldots, f_t) : M \to M^{\oplus t})$를 만족하는 것으로 한다. 표준 동형 $$
H^0(U, \widetilde{M}) = \colim \Hom_A(I^n, M).
$$가 존재함을 보여라. 주의: 이는 자명하지 않다. \end{exercise}

\begin{exercise} \label{exercises-exercise-Noetherian} $A$를 뇌터 환이라고 하자. $I$를 $A$의 아이디얼로 한다. $M$를 $A$-가군으로 한다. $M[I^\infty]$를, 다음과 같이 정의되는 $I$-거듭제곱 비틀림 원소들의 집합으로 한다: $$
M[I^\infty] = \{x \in M \mid
\text{어떤 }n \geq 1\text{이 존재하여 }I^nx = 0\}
$$ $M' = M/M[I^\infty]$라고 하자. 이때 $M'$의 $I$-거듭제곱 꼬임은 0이다. $$
\colim \Hom_A(I^n, M) = \colim \Hom_A(I^n, M').
$$임을 보이라. % P12-ERR-0470: preserve the duplicated final hint number (3).
주의: 이는 자명하지 않다. 힌트: (1) $M$가 유한한 경우로 귀결되도록 시도하라, (2) $M$가 유한하고 $N = M[I^\infty]$일 때, $\Ext^1_A(I^n, N)$의 임의의 원소는 어떤 $m > 0$에 대한 $\Ext^1_A(I^{n + m}, N)$로의 사상에서 0이 됨을 보여라, (3) $\Hom_A(I^n, N)$에 대해 (2)와 동일한 것을 보여라, (3) $M$가 유한한 경우에 긴 완전열 $$
0 \to \Hom_A(I^n, M[I^\infty]) \to \Hom_A(I^n, M) \to \Hom_A(I^n, M')
\to \Ext^1_A(I^n, M[I^\infty])
$$를 고려하고, $I^{n + m}$에 대한 열과 비교하여 결론을 도출하라. \end{exercise}


\section{코호몰로지 재방문} \label{exercises-section-more-more-cohomology}

\begin{exercise} \label{exercises-exercise-nonkillable} 체 $k$,$k$ 위의 곡선 $X$,가역 $\mathcal{O}_X$-가군 $\mathcal{L}$,및 코호몰로지 클래스 $\xi \in H^1(X, \mathcal{L})$ 로서, 다음 성질을 가지는 것의 예를 만들어라: 스킴의 임의의 전사 유한 사상 $\pi : Y \to X$에 대해, 원소 $\xi$의 끌어옴은 $H^1(Y, \pi^*\mathcal{L})$의 0이 아닌 원소이다. 힌트: 닫힌 부분 스킴 $Z \subset X$가 $1$개보다 많은 닫힌 점으로 이루어지고, 짧은 완전 열 $0 \to \mathcal{L} \to \mathcal{O}_X \to i_*\mathcal{O}_Z \to 0$이 존재하도록 $X$, $k$, $\mathcal{L}$를 구성하라. % P12-ERR-0471: preserve the frozen pullback-as-verb source locus.
그 후, 당겨올렸을 때 무슨 일이 일어나는지 조사하라. \end{exercise}

\begin{exercise} \label{exercises-exercise-cohomology-cycle-curves} $k$를 대수적으로 닫힌 체로 한다. $X$를 사영적인 $1$차원 스킴으로 한다. $X$이 곡선들의 순환을 포함한다고 가정한다. 즉, $n \geq 2$, 서로 다른 $1$ 차원 정 닫힌 부분스킴 $C_1, \ldots, C_n$, 그리고 서로 다른 닫힌 점 $x_1, \ldots, x_n \in X$가 존재하여, $x_n \in C_n \cap C_1$, 그리고 $i = 1, \ldots, n - 1$에 대하여 $x_i \in C_i \cap C_{i + 1}$이라고 하자. $H^1(X, \mathcal{O}_X)$이 0이 아님을 증명하라. 힌트: $\mathcal{F}$을 사상 % P12-ERR-0472: preserve the frozen unindexed direct sum.
$\mathcal{O}_X \to \bigoplus \mathcal{O}_{C_i}$의 상으로 하고, $\kappa(x_i) = k$ 및 $H^0(C_i, \mathcal{O}_{C_i}) = k$을 이용하여 $H^1(X, \mathcal{F})$가 영이 아님을 증명하라. 또한, Grothendieck의 정리에 따라 $H^2(X, -) = 0$임을 이용하라. \end{exercise}

\begin{exercise} \label{exercises-exercise-surface} $X$을 대수폐포체 $k$ 위의 사영면으로 하자. $H^1(Z, \mathcal{O}_Z)$가 0이 아닌 진 닫힌 부분스킴 $Z \subset X$이 존재함을 증명하시오. 힌트: 연습문제 \ref{exercises-exercise-cohomology-cycle-curves}를 사용하시오. \end{exercise}

\begin{exercise} \label{exercises-exercise-surface-better} $X$을 대수적으로 닫힌 체 $k$ 위의 사영곡면으로 하자. 임의의 $n \geq 0$에 대하여, $\dim_k H^1(Z, \mathcal{O}_Z) > n$가 되는 진 닫힌 부분스킴 $Z \subset X$이 존재함을 보이시오. 연습문제 \ref{exercises-exercise-surface} 및 \ref{exercises-exercise-cohomology-cycle-curves}의 논의를 어떻게 수정하면 좋은지에 대해서만 설명하고, 세부 사항을 모두 제시할 필요는 없다. \end{exercise}

\begin{exercise}\label{exercises-exercise-surface-h2-nonzero}$X$를 대수적으로 닫힌 체 $k$ 위의 사영 곡면이라고 하자. $H^2(X, \mathcal{F})$이 0이 아닌 연접 $\mathcal{O}_X$-가군 $\mathcal{F}$이 존재함을 증명하라. 힌트: 연습문제 \ref{exercises-exercise-surface-better}의 결과와,교묘하게 선택한 완전 열을 사용하라. \end{exercise}

\begin{exercise} \label{exercises-exercise-kunneth} $X$와 $Y$를 체 $k$ 위의 스킴으로 한다 (필요하다면 $X$와 $Y$는 좋은 스킴, 예를 들어 qcqs 또는 $k$ 위 사영적이라고 가정해도 좋다). $X \times Y = X \times_{\Spec(k)} Y$라고 두고, 사영을 $p : X \times Y \to X$ 및 $q : X \times Y \to Y$로 한다. 준연접 $\mathcal{O}_X$-가군 $\mathcal{F}$과 준연접 $\mathcal{O}_Y$-가군 $\mathcal{G}$에 대해 $$
H^n(X \times Y,
p^*\mathcal{F} \otimes_{\mathcal{O}_{X \times Y}} q^*\mathcal{G}) =
\bigoplus\nolimits_{a + b = n}
H^a(X, \mathcal{F}) \otimes_k H^b(Y, \mathcal{G})
$$를 증명하시오. 또는, 차원을 취했을 때 이것이 성립함만을 보여도 좋다. '간결한' 답변에는 추가 점수를 부여한다. \end{exercise}

\begin{exercise} \label{exercises-exercise-Oab} $k$를 체로 한다. % P12-ERR-0473: preserve the frozen malformed first projective-line factor.
$X = \mathbf{P}|^1 \times \mathbf{P}^1$을 $k$ 위의 사영 직선과 그것 자체의 곱으로 하고, 사영을 $p : X \to \mathbf{P}^1_k$ 및 $q : X \to \mathbf{P}^1_k$로 한다. $$
\mathcal{O}(a, b) = p^*\mathcal{O}_{\mathbf{P}^1_k}(a)
\otimes_{\mathcal{O}_X} q^*\mathcal{O}_{\mathbf{P}^1_k}(b)
$$라고 한다. 모든 $i, a, b$에 대해 $H^i(X, \mathcal{O}(a, b))$의 차원을 계산하라. 힌트: 연습문제 \ref{exercises-exercise-kunneth}을 사용하라. \end{exercise}


\section{코호몰로지와 Hilbert 다항식} \label{exercises-section-cohomology-hilbert-polynomials}

\begin{situation} \label{exercises-situation-hilbert-polynomial} $k$을 체로 한다. $X = \mathbf{P}^n_k$을 $n$차원 사영 공간으로 한다. $\mathcal{F}$을 연접 $\mathcal{O}_X$-가군으로 한다. $$
\chi(X, \mathcal{F}) =
\sum\nolimits_{i = 0}^n (-1)^i \dim_k H^i(X, \mathcal{F})
$$였음을 상기하라. $\mathcal{F}$의 {\it Hilbert 다항식}이란, 함수 $$
t \longmapsto \chi(X, \mathcal{F}(t))
$$였던 것이다. 또한, $\mathcal{F}(t) = \mathcal{F} \otimes_{\mathcal{O}_X} \mathcal{O}_X(t)$이며, 여기서 $\mathcal{O}_X(t)$은 구성, 정의 \ref{constructions-definition-twist}에서의 구조 층의 $t$ 차 꼬임이었다. 다양체, \ref{varieties-subsection-hilbert} 절에서, Hilbert 다항식이 $t$의 다항식임을 증명하였다. \end{situation}

\begin{exercise} \label{exercises-exercise-hilbert-pol-easy} 설정 \ref{exercises-situation-hilbert-polynomial} 하에서 고려한다. \begin{enumerate} \item $P(t)$가 $\mathcal{F}$의 Hilbert 다항식이라면, $\mathcal{F}(-13)$의 Hilbert 다항식은 무엇인가. \item $P_i$가 $\mathcal{F}_i$의 Hilbert 다항식이라면, $\mathcal{F}_1 \oplus \mathcal{F}_2$의 Hilbert 다항식은 무엇인가. \item $P_i$가 $\mathcal{F}_i$의 Hilbert 다항식이고, $\mathcal{F}$가 전사 $\mathcal{F}_1 \to \mathcal{F}_2$의 핵일 때, $\mathcal{F}$의 Hilbert 다항식은 무엇인가. \end{enumerate} \end{exercise}

\begin{exercise} \label{exercises-exercise-find-given-hilbert-pol-dim-1} 설정 \ref{exercises-situation-hilbert-polynomial}에서 $n \geq 1$라고 가정하자. Hilbert 다항식이 $t - 101$인 연접층을 찾아라. \end{exercise}

\begin{exercise} \label{exercises-exercise-find-given-hilbert-pol-dim-2} 설정 \ref{exercises-situation-hilbert-polynomial} 에서 $n \geq 2$라고 가정한다. 힐베르트 다항식이 $t^2/2 + t/2 - 1$인 연접층을 찾아라. (이는 조금 어렵다. 그러한 연접층이 존재함만 보이면 충분하다.) \end{exercise}

\begin{exercise} \label{exercises-exercise-bound-genus-in-degree} 설정 \ref{exercises-situation-hilbert-polynomial} 에서 $n \geq 2$ 및 $k$은 대수적으로 닫힌 체라고 가정한다. $C \subset X$을 차원 $1$의 정 닫힌 부분스킴으로 한다. 다시 말해, $C$은 사영 곡선이다. $X$ 위의 연접층으로 간주한 $\mathcal{O}_C$의 Hilbert 다항식을 $d t + e$이라고 한다. \begin{enumerate} \item $e$의 상한을 제시하라. (힌트: $\mathcal{O}_C(t)$은 차수 $0$와 $1$에 대해서만 코호몰로지를 가지는 것을 이용하여, $H^0(C, \mathcal{O}_C)$을 조사하라. ) \item $C$와 횡단적으로 교차하는 $\mathcal{O}_X(1)$의 대역 단면 $s$, 즉 서로 다른 닫힌 점 $c_1, \ldots, c_r \in C$과 짧은 완전열 $$
0 \to \mathcal{O}_C \xrightarrow{s} \mathcal{O}_C(1) \to
\bigoplus\nolimits_{i = 1, \ldots, r} k_{c_i} \to 0
$$이 존재하는 것을 선택하라. 여기서 $k_{c_i}$은 $c_i$에서 값이 $k$인 마천루 층이다. (이러한 $s$는 존재한다. 여기서는 그 사실을 사용해도 좋다.) $r = d$임을 증명하라. (힌트: 이 열을 비틀어서 무엇을 얻을 수 있는지 살펴보라. ) \item 짧은 완전열을 꼬면, 임의의 $t$에 대해 $k$-선형 사상 $\varphi_t : \Gamma(C, \mathcal{O}_C(t)) \to \bigoplus_{i = 1, \ldots, d} k$을 얻을 수 있다. % P12-ERR-0474: preserve the frozen extraneous-if source locus.
$t \geq d - 1$이라면 이 사상이 전사임을 보이라. \item $d$을 사용하여 $e$의 하한을 제시하라. (힌트: (3)의 결과를 사용하여 $H^1(C, \mathcal{O}_C(d - 2)) = 0$임을 보이고, 소멸을 이용하라. ) \end{enumerate} \end{exercise}

\begin{exercise} \label{exercises-exercise-three-quadrics-in-plane} 설정 \ref{exercises-situation-hilbert-polynomial} 에서 $n = 2$라고 가정한다. $s_1, s_2, s_3 \in \Gamma(X, \mathcal{O}_X(2))$를 세 개의 이차방정식으로 한다. 연접층 $$
\mathcal{F} = \Coker\left(\mathcal{O}_X(-2)^{\oplus 3}
\xrightarrow{s_1, s_2, s_3} \mathcal{O}_X\right)
$$를 생각한다. 이러한 $\mathcal{F}$가 가질 수 있는 Hilbert 다항식을 나열하라. (사영평면 내의 이차곡선의 교점을 시각화해보라.) \end{exercise}


\section{곡선}
\label{exercises-section-curves}

\begin{exercise}
\label{exercises-exercise-closed-subset-vanshing}
$k$를 대수적으로 닫힌 체라 하자. $X$를 $k$ 위의 사영 곡선이라 하고,
$\mathcal{L}$을 가역 $\mathcal{O}_X$-가군이라 하자.
$s_0, \ldots, s_n \in H^0(X, \mathcal{L})$을 $\mathcal{L}$의 대역 단면이라
하자. 닫힌 점 $((\lambda_0 : \ldots : \lambda_n), x)$가 $Z$에 속할
필요충분조건이 단면 $\lambda_0 s_0 + \ldots + \lambda_n s_n$가 $x$에서
영이 되는 것인, 자연스러운 닫힌 부분스킴
$$
Z \subset
\mathbf{P}^n \times X
$$
이 존재함을 증명하라. (힌트: $Z$를 아핀 국소적으로 기술하라.)
\end{exercise}

\begin{exercise}
\label{exercises-exercise-diagonal-cartier}
$k$를 대수적으로 닫힌 체라 하자. $X$를 $k$ 위의 매끄러운 곡선이라 하자.
$r \geq 1$이라 하자. 닫힌 점들이 어떤 $i$에 대하여 $x = x_i$를 만족하는
튜플 $(x, x_1, \ldots, x_r)$인 닫힌 부분집합
$$
D \subset X \times X^r = X^{r + 1}
$$
이 가역 아이디얼층을 가짐을 보여라. 달리 말해 $D$가 유효 Cartier
인자임을 보여라. 힌트: $r = 1$로 귀착하고, $X$가 매끄러운 곡선이라는
사실을 써서 대각선에 관하여 무언가를 말하라(이에 대해서는 책을 찾아보라).
\end{exercise}

\begin{exercise}
\label{exercises-exercise-one-divisor-in-another}
$k$를 대수적으로 닫힌 체라 하자. $X$를 $k$ 위의 매끄러운 사영 곡선이라
하자. $T$를 $k$ 위의 유한형 스킴이라 하고,
$$
D_1 \subset X \times T
\quad\text{and}\quad
D_2 \subset X \times T
$$
를 두 유효 Cartier 인자라 하자. 각 $t \in T$에 대하여 올
$D_{i, t} \subset X_t$가 조밀하지 않다고 하자(즉, 곡선 $X_t$의 일반점을
포함하지 않는다고 하자). 닫힌 점 $t \in T$가 $Z$에 속할 필요충분조건이
$D_1$, $D_2$의 스킴론적 올 $D_{1, t}$, $D_{2, t}$가 $X_t$의 닫힌
부분스킴으로서
$$
D_{1, t} \subset D_{2, t}
$$
를 만족하는 것인 표준적인 닫힌 부분스킴 $Z \subset T$가 존재함을
증명하라. 힌트: 필요하면 $T$를 축소하여 $T = \Spec(A)$가 아핀이고
$D_i \subset U \times T$가 되는 아핀 열린집합 $U \subset X$가 있다고
가정해도 됨을 보여라. 그런 다음
$M_1 = \Gamma(D_1, \mathcal{O}_{D_1})$이 유한 국소 자유 $A$-가군임을
보여라(여기서는 자명하지 않은 대수학이 필요하므로 친구들에게 물어보라).
$T$를 축소하여 $M_1$이 자유 $A$-가군이라고 가정해도 된다. 이제
$$
\Gamma(U \times T, \mathcal{I}_{D_2}) \to M_1 = A^{\oplus N}
$$
을 살펴보고, 이 사상의 상에 속하는 벡터들의 계수들이 생성하는
아이디얼로 잘라 낸 닫힌 부분스킴으로 $Z$를 정의하라. 이것이 성립하는
이유를 설명하라(아마 가환대수학이 조금 더 필요할 것이다).
\end{exercise}

\begin{exercise}
\label{exercises-exercise-closed-subset-vanshing-r}
$k$를 대수적으로 닫힌 체라 하자. $X$를 $k$ 위의 매끄러운 사영 곡선이라
하고, $\mathcal{L}$을 가역 $\mathcal{O}_X$-가군이라 하자.
$s_0, \ldots, s_n \in H^0(X, \mathcal{L})$을 $\mathcal{L}$의 대역 단면이라
하자. $r \geq 1$이라 하자. 닫힌 점
$((\lambda_0 : \ldots : \lambda_n), x_1, \ldots, x_r)$가 $Z$에 속할
필요충분조건이 단면
$s_\lambda = \lambda_0 s_0 + \ldots + \lambda_n s_n$가 인자
$D = x_1 + \ldots + x_r$ 위에서 영이 되는 것, 즉 단면 $s_\lambda$가
$\mathcal{L}(-D)$에 속하는 것인 자연스러운 닫힌 부분스킴
$$
Z \subset
\mathbf{P}^n \times X \times \ldots \times X = \mathbf{P}^n \times X^r
$$
이 존재함을 증명하라. 힌트: 연습문제
\ref{exercises-exercise-diagonal-cartier}와
\ref{exercises-exercise-one-divisor-in-another}의 결과를 결합하면 이를 얻는다는 것을
설명하라.
\end{exercise}

\begin{exercise}
\label{exercises-exercise-effective}
$k$를 대수적으로 닫힌 체라 하자. $X$를 $k$ 위의 매끄러운 사영 곡선이라
하고, $\mathcal{L}$을 가역 $\mathcal{O}_X$-가군이라 하자. $X^r$의 닫힌 점
$(x_1, \ldots, x_r)$가 $Z$에 속할 필요충분조건이
$\mathcal{L}(-x_1 - \ldots -x_r)$이 영이 아닌 대역 단면을 갖는 것인
자연스러운 닫힌 부분집합
$$
Z \subset X^r
$$
이 존재함을 보여라. 힌트: 연습문제
\ref{exercises-exercise-closed-subset-vanshing-r}을 사용하라.
\end{exercise}

\begin{exercise}
\label{exercises-exercise-difference-effective}
$k$를 대수적으로 닫힌 체라 하자. $X$를 $k$ 위의 매끄러운 사영 곡선이라
하자. $r \geq s$를 정수라 하자. $X^r \times X^s$의 닫힌 점
$(x_1, \ldots, x_r, y_1, \ldots, y_s)$가 $Z$에 속할 필요충분조건이
$x_1 + \ldots + x_r - y_1 - \ldots - y_s$가 유효 인자와 선형 동치인
것인 자연스러운 닫힌 부분집합
$$
Z \subset X^r \times X^s
$$
이 존재함을 보여라. 힌트: 차수가 매우 큰 보조 가역 가군
$\mathcal{L}$을 택하여, 차수 $r$인 모든 유효 인자 $D$에 대하여
$\mathcal{L}(-D)$가 영이 되지 않는 단면을 갖게 하라. 그런 다음
연습문제 \ref{exercises-exercise-effective}의 결과를 두 번 사용하라.
\end{exercise}

\begin{exercise}
\label{exercises-exercise-genus-7}
좋아하는 대수적으로 닫힌 체 $k$를 택하라. 종수 $7$인 어떤 곡선 위에
존재할 수 있는 모든 $\mathfrak g^r_d$를 가능한 한 잘 결정하라. 이 과정에서
다음도 해 보라.
\begin{enumerate}
\item 어떤 경우에 $\mathfrak g^r_d$가 밑점을 갖지 않는지 결정하라.
\item 어떤 경우에 $\mathfrak g^r_d$가 $\mathbf{P}^r$로의 닫힌 몰입을
주는지 결정하라.
\end{enumerate}
곡선이 ``일반적''이라고 가정한 경우에도 같은 일을 하라(일반적이라는
개념은 스스로 정해도 되며, 이는 위 문제보다 쉬울 수 있다). 곡선이
초타원 곡선이라고 가정한 경우에도 같은 일을 하라. 곡선이 삼각 곡선이고
(초타원 곡선은 아니라고) 가정한 경우에도 같은 일을 하라. 기타 등등.
\end{exercise}






\section{모듈라이} \label{exercises-section-moduli}

\noindent 이 절에서는, 대수기하학적 대상의 모듈라이에 대한 소박한 취급을 몇 가지 생각해본다.

\medskip\noindent $k$를 대수적으로 닫힌 체라고 하자.$M$를 $k$ 위의 모듈라이 문제라고 하자. 이것이 정확히 무엇을 의미하는지는 여기서 정의하지 않지만, 각 연습에서 의도는 명확할 것이다. 이하를 이해하려면, $k$ 위의 $M$ 대상, $k$ 위의 $M$ 대상 간의 동형, 그리고 다양체 위의 % P12-ERR-0478: preserve the frozen singular-object source locus.
$M$ 대상의 족이 무엇인지 알면 충분하다. 이때, 다음 조건이 성립할 경우, {\it $M$의 모듈라이 수}는 $d \geq 0$라고 한다: \begin{enumerate} \item 유한 개의 족 $X_i \to V_i$, $i = 1, \ldots, n$가 존재하며, $k$ 위의 $M$의 임의의 대상은 그것들 중 어느 하나의 올와 동형이며, 또한 $\max \dim(V_i) = d$이다. \item 이것을 더 작은 $d$으로 수행할 수 없다. \end{enumerate} 이는 실제로 문헌에 있는 더 나은 개념의 매우 거친 근사에 지나지 않는다.

\begin{exercise} \label{exercises-exercise-number-moduli-lines-conics} $k$을 대수적으로 닫힌 체로 하자. $d \geq 1$ 및 $n \geq 1$으로 한다. $P^n$ 내의 차수 $d$의 초곡면의 모듈라이를 다음과 같이 정한다: \begin{enumerate} \item 대상은 차수 $d$의 초곡면 $X \subset \mathbf{P}^n_k$이다. \item 두 대상 $X \subset \mathbf{P}^n_k$과 $Y \subset \mathbf{P}^n_k$ 사이의 동형은, % P12-ERR-0479: preserve the frozen PGL_n formula.
$g(X) = Y$을 만족하는 원소 $g \in \text{PGL}_n(k)$이다. \item 다양체 $V$ 위의 초곡면들의 족은 닫힌 부분 스킴 $X \subset \mathbf{P}^n_V$이며, 임의의 $v \in V$에 대해, % P12-ERR-0480: preserve the frozen singular/plural source locus.
$X \to V$의 스킴론적 올 $X_v$가 $\mathbf{P}^n_v$ 내의 초곡면인 것이다. \end{enumerate} 다음 모듈라이 수를 계산하라 (가능하다면). 나중으로 갈수록 어려워진다)： \begin{enumerate} \item $n = 1$에서 $d$가 임의인 경우, \item $n = 2$ 및 $d = 1$의 경우, \item $n = 2$ 및 $d = 2$의 경우, \item $n \geq 1$ 및 $d = 2$의 경우, \item $n = 2$ 및 $d = 3$의 경우, \item $n = 3$ 및 $d = 3$의 경우, 및 \item $n = 2$ 그리고 $d = 4$의 경우. \end{enumerate} \end{exercise}

\begin{exercise} \label{exercises-exercise-number-moduli-hyperelliptic} $k$를 대수적으로 닫힌 체로 한다. $g \geq 2$라고 한다. 종수 $g$의 초타원곡선의 모듈라이를 다음과 같이 정의한다: \begin{enumerate} \item 대상은 종수 $g$의 매끄러운 사영 초타원곡선 $X$이다. \item 두 대상 $X$과 $Y$ 사이의 동형은, $k$ 위의 스킴 동형 $X \to Y$이다. \item 다양체 $V$ 위의 종수 $g$ 초타원 곡선의 가족은, % P12-ERR-0481: preserve the frozen X-to-Y family morphism.
고유 평탄 \footnote{이 가정을 제거해도 문제의 답은 변하지 않는다.} 사상 $X \to Y$으로서, $X \to V$의 모든 스킴론적 올가 종수 $g$의 매끄러운 사영 초타원 곡선인 것이다. \end{enumerate} 모듈라이 수가 $2g - 1$ 임을 보여라. \end{exercise}


\section{대역 Ext} \label{exercises-section-global-exts}

\begin{exercise} \label{exercises-exercise-ext-line-plane-p3} $k$를 체라 하고 $X = \mathbf{P}^3_k$라 하자. $L \subset X$와 $P \subset X$를 각각 $1$개와 $2$개의 일차방정식으로 잘라낸 닫힌 부분스킴으로 본 직선과 평면이라 하자. 모든 $i$에 대해 $$
\Ext^i_X(\mathcal{O}_L, \mathcal{O}_P)
$$의 차원을 계산하라. $L$가 $P$에 포함되는 경우와, $L$가 $P$에 포함되지 않는 경우의 양쪽을 반드시 다루어야 한다. \end{exercise}

\begin{exercise} \label{exercises-exercise-ext-point} $k$를 체로 한다. $X = \mathbf{P}^n_k$로 한다. $Z \subset X$를 닫힌 부분 스킴으로 간주한 닫힌 $k$-유리점으로 한다. 예를 들어 동차 좌표 $(1 : 0 : \ldots : 0)$를 가진 점이다. 모든 $i$에 대해 $$
\Ext^i_X(\mathcal{O}_Z, \mathcal{O}_Z)
$$의 차원을 계산하라. \end{exercise}

\begin{exercise} \label{exercises-exercise-ext-pairings} $X$를 환 달린 공간이라 하자. $\mathcal{O}_X$-가군 $\mathcal{F}, \mathcal{G}, \mathcal{H}$에 대해 컵곱 사상 $$
\Ext^i_X(\mathcal{G}, \mathcal{H})
\times
\Ext^j_X(\mathcal{F}, \mathcal{G})
\longrightarrow
\Ext^{i + j}_X(\mathcal{F}, \mathcal{H})
$$을 정의하라. (힌트: 이것은 매우 일반적인 구성이다.) \end{exercise}

\begin{exercise} \label{exercises-exercise-move-out-locally-free} $X$를 환 달린 공간이라 하자. $\mathcal{E}$를 유한 국소 자유 $\mathcal{O}_X$-가군이라 하고, 그 쌍대를 $\mathcal{E}^\vee = \SheafHom_{\mathcal{O}_X}(\mathcal{E}, \mathcal{O}_X)$라 하자. 다음을 증명하라. \begin{enumerate} \item $\SheafExt^i_{\mathcal{O}_X}(
\mathcal{F} \otimes_{\mathcal{O}_X} \mathcal{E}, \mathcal{G}) =
\SheafExt^i_{\mathcal{O}_X}(
\mathcal{F},
\mathcal{E}^\vee \otimes_{\mathcal{O}_X} \mathcal{G}) =
\SheafExt^i_{\mathcal{O}_X}(
\mathcal{F}, \mathcal{G}) \otimes_{\mathcal{O}_X} \mathcal{E}^\vee$, 그리고 \item $\Ext^i_X(
\mathcal{F} \otimes_{\mathcal{O}_X} \mathcal{E}, \mathcal{G}) =
\Ext^i_X(\mathcal{F},
\mathcal{E}^\vee \otimes_{\mathcal{O}_X} \mathcal{G})$. \end{enumerate} 여기서 $\mathcal{F}$과 $\mathcal{G}$은 $\mathcal{O}_X$-가군이다. 다음을 이끌어라: $$
\Ext^i_X(\mathcal{E}, \mathcal{G}) =
H^i(X, \mathcal{E}^\vee \otimes_{\mathcal{O}_X} \mathcal{G})
$$ \end{exercise}

\begin{exercise} \label{exercises-exercise-trace-on-ext} $X$를 환이 있는 공간으로 하자. $\mathcal{E}$을 유한 국소 자유 $\mathcal{O}_X$-가군으로 하자. 모든 $i$에 대해 트레이스 사상 $$
\Ext^i_X(\mathcal{E}, \mathcal{E}) \to H^i(X, \mathcal{O}_X)
$$을 구성하라. 임의의 $\mathcal{O}_X$-가군 $\mathcal{F}$에 대한 트레이스 사상 $$
\Ext^i_X(\mathcal{E}, \mathcal{E} \otimes_{\mathcal{O}_X} \mathcal{F})
\to H^i(X, \mathcal{F})
$$으로 일반화하라. \end{exercise}

\begin{exercise} \label{exercises-exercise-duality} $k$를 체로 한다. $X = \mathbf{P}^d_k$라고 한다. $\omega_{X/k} = \mathcal{O}_X(-d - 1)$라고 둔다. 유한 국소 자유 가군 $\mathcal{E}$, $\mathcal{F}$에 대해, Ext 위의 컵 곱과 Ext 위의 트레이스 사상을 결합한 사상 $$
\Ext^i_X(\mathcal{E}, \mathcal{F} \otimes_{\mathcal{O}_X} \omega_{X/k})
\times
\Ext^{d - i}_X(\mathcal{F}, \mathcal{E})
\to
\Ext_X^d(\mathcal{F}, \mathcal{F} \otimes_{\mathcal{O}_X} \omega_{X/k})
\to
H^d(X, \omega_{X/k}) = k
$$이 비퇴화 쌍을 준다는 것을 증명하라. 힌트: 이 설정에서 쌍대성을 새로이 증명해도 좋고, 층의 코호몰로지로 환원하여 강의에서 증명한 Serre 쌍대 정리를 적용해도 좋다. \end{exercise}

\section{인자} \label{exercises-section-divisors}

\noindent 편의를 위해, 관련된 모든 정의를 여기 한 곳에 모은다.

\begin{definition} \label{exercises-definition-divisor} 이하 $S$를 임의의 스킴이라 하고 $X$를 뇌터 정 스킴이라 하자. \begin{enumerate} \item $X$ 위의 {\it Weil 인자}란 소인자 $Z_i$들의 정수 계수 형식적 선형 결합 $\Sigma n_i[Z_i]$을 말한다. \item {\it 소인자}란 정 스킴인 닫힌 부분스킴 $Z \subset X$로서, 그 일반점 $\xi \in Z$에서 ${\mathcal O}_{X, \xi}$의 차원이 $1$인 것을 말한다. 위와 같은 $\xi \in Z \subset X$에 대해 기호 ${\mathcal O}_{X, Z} = {\mathcal O}_{X, \xi}$를 사용한다. ${\mathcal O}_{X, Z} \subset K(X)$는 $X$의 함수체의 부분환임에 주의하라. \item {\it 유리 함수 $f \in K(X)^\ast$에 수반되는 Weil 인자}란 합 $\Sigma v_Z(f)[Z]$를 말한다. 여기서 $v_Z(f)$를 다음과 같이 정의한다. \begin{enumerate} \item $f \in {\mathcal O}_{X, Z}^\ast$이면 $v_Z(f) = 0$으로 둔다. \item $f \in {\mathcal O}_{X, Z}$이면 $$
v_Z(f) = \text{length}_{{\mathcal O}_{X, Z}}({\mathcal O}_{X, Z}/(f)).
$$로 한다. \item $a, b \in {\mathcal O}_{X, Z}$를 사용하여 $f = \frac{a}{b}$로 쓸 수 있다면 $$
v_Z(f) = \text{length}_{{\mathcal O}_{X, Z}}({\mathcal O}_{X, Z}/(a)) -
\text{length}_{{\mathcal O}_{X, Z}}({\mathcal O}_{X, Z}/(b)).
$$로 한다. \end{enumerate} \item 스킴 $S$ 위의 {\it 유효 Cartier 인자}란, 임의의 점 $d\in D$이 $S$ 내에 아핀 열린 근방 % P12-ERR-0482: preserve the frozen duplicated ambient-space locus.
$\Spec(A) = U \subset S$을 가지며, 그곳에서 $D \cap U = \Spec(A/(f))$이 영인자가 아닌 원소 $f \in A$에 의해 정해지는 닫힌 부분 스킴 $D \subset S$을 말한다. \item 뇌터 정 스킴 $X$의 유효 Cartier 인자 $D \subset X$에 {\it 수반하는 Weil 인자 $[D]$}란, 합 $\Sigma v_Z(D)[Z]$을 의미한다. 여기서 $v_Z(D)$을 다음과 같이 정의한다: \begin{enumerate} \item $Z$의 일반점 $\xi$이 $D$에 속하지 않는다면 $v_Z(D) = 0$로 한다. \item $Z$의 일반 점 $\xi$이 $D$에 속한다면 $$
v_Z(D) = \text{length}_{{\mathcal O}_{X, Z}}({\mathcal O}_{X, Z}/(f))
$$라고 한다. 여기서 $f \in {\mathcal O}_{X, Z} = {\mathcal O}_{X, \xi}$는 $\xi$의 아핀 근방에서 $D$를 정의하는 영인자가 아닌 원소(위의 (4) 의 의미에서의 것)이다. \end{enumerate} \item $S$를 스킴으로 한다. {\it 전분수환층 ${\mathcal K}_S$}이란, 다음과 같이 정의되는 전층 ${\mathcal K}'$을 층화하여 얻어지는 ${\mathcal O}_S$-대수의 층을 말한다. 열린 집합 $U \subset S$에 대해 ${\mathcal K}'(U) = S_U^{-1}{\mathcal O}_S(U)$라고 하자. 여기서 $S_U \subset {\mathcal O}_S(U)$는, 단면 $f \in {\mathcal O}_S(U)$으로서 모든 $u\in U$에 대해 ${\mathcal O}_{S, u}$에서의 $f$의 싹이 영인자가 아닌 원소가 되는 곱셈적 부분집합으로 구성된다. 특히, $S_U$의 모든 원소는 영인자가 아닌 원소이다. 따라서 ${\mathcal O}_S$는 ${\mathcal K}_S$의 부분층이며, 짧은 완전열 $$
0 \to {\mathcal O}_S^\ast \to {\mathcal K}_S^\ast \to
{\mathcal K}_S^\ast/{\mathcal O}_S^\ast \to 0.
$$을 얻는다. \item 스킴 $S$ 위의 {\it Cartier 인자}란, 몫층 ${\mathcal K}_S^\ast/{\mathcal O}_S^\ast$의 전역 단면을 의미한다. \item 뇌터 정 스킴 $X$ 위의 Cartier 인자 $\tau \in \Gamma(X, {\mathcal K}_X^\ast/{\mathcal O}_X^\ast)$에 {\it 부속된 Weil 인자}란, 합 $\Sigma v_Z(\tau)[Z]$을 의미한다. % P12-ERR-0483: preserve the frozen extraneous-as source locus.
여기서 $v_Z(\tau)$는 다음 절차로 정의된다: \begin{enumerate} \item $Z$의 일반점 $\xi$에서 $\tau$의 싹이 0인 경우, 다시 말해 줄기 $({\mathcal K}^\ast/{\mathcal O}^\ast)_\xi$에서 $\tau$의 상이 '0'이라면, $v_Z(\tau) = 0$라고 한다. \item $\tau|_U$이 단면 $f \in {\mathcal K}(U)$의 상이며, 게다가 $a, b \in A$을 사용하여 $f = a/b$로 쓸 수 있는 아핀 열린 이웃 $\Spec(A) = U \subset X$을 찾아라. 이때 $$
v_Z(f) = \text{length}_{{\mathcal O}_{X, Z}}({\mathcal O}_{X, Z}/(a)) -
\text{length}_{{\mathcal O}_{X, Z}}({\mathcal O}_{X, Z}/(b)).
$$이라 하자. \end{enumerate} \end{enumerate} \end{definition}

\begin{remarks} \label{exercises-remarks-divisors} 자명한 주의를 몇 가지 언급한다. \begin{enumerate} \item 뇌터 정 스킴 $X$ 위에서는, 층 ${\mathcal K}_X$은 함수체 $K(X)$를 값으로 가지는 상수층이다. \item 위의 정의를 의미 있게 만들기 위해서는, 1차원 뇌터 국소 정역 ${\mathcal O}$의 임의의 0이 아닌 원소의 쌍 $(a, b)$에 대해 $$
\text{length}_{\mathcal O}({\mathcal O}/(ab)) =
\text{length}_{\mathcal O}({\mathcal O}/(a)) +
\text{length}_{\mathcal O}({\mathcal O}/(b))
$$를 보여야 한다. 이는 강의에서 다룬다. \end{enumerate} \end{remarks}

\begin{exercise} \label{exercises-exercise-effective-cartier-cartier} (임의의 스킴 위에서.) 유효한 Cartier 인자에 Cartier 인자를 대응시키는 방법을 기술하라. \end{exercise}

\begin{exercise} \label{exercises-exercise-rational-function-cartier} (정 스킴 위에서.) 유리함수 $f$에 Cartier 인자 $D$를 대응시키고, $D$와 $f$에 수반되는 Weil 인자가 일치하도록 하는 방법을 기술하라. (이는 시시한 것이다. ) \end{exercise}

\begin{exercise} \label{exercises-exercise-weil-not-cartier} 어떤 Cartier 인자에도 부속되지 않는, 다양체 위의 Weil 인자의 예를 들어라. \end{exercise}

\begin{exercise} \label{exercises-exercise-weil-Q-cartier} 어떤 Cartier 인자에도 부속되지 않는 다양체 위의 Weil 인자 $D$ 로서, 어떤 $n > 1$에 대해 $nD$가 Cartier 인자에 부속되는 Weil 인자가 되는 예를 들어라. \end{exercise}

\begin{exercise} \label{exercises-exercise-weil-not-Q-cartier} 어떤 Cartier 인자에도 수반되지 않고, 또한 임의의 $n > 1$에 대해 $nD$도 Cartier 인자에 수반되지 않는 Weil 인자인, 다양체 위의 Weil 인자 $D$의 예를 들어라. (힌트: 원뿔, 예를 들어 $\mathbf{A}^4_k$ 내의 $X : xy - zw = 0$를 생각해 보라. $D = [x = 0, z = 0]$이 조건을 만족함을 보여 보라. ) \end{exercise}

\begin{exercise} \label{exercises-exercise-cartier-not-difference-effective-cartier} 체상 유한형 분리 스킴 $X$에 대하여: 두 개의 유효한 Cartier 인자의 차이가 아닌 Cartier 인자의 예를 들어라. 힌트: 공집합이 아닌 유효한 Cartier 인자를 하나도 가지지 않는 $X$, 예를 들어 \cite[III Exercise 5.9]{H}로 구성된 스킴을 찾아라. $X$가 다양체인 예시도 존재한다. 즉, 연습 \ref{exercises-exercise-nonprojective}의 다양체이다. \end{exercise}

\begin{exercise} \label{exercises-exercise-nonprojective} 비사영적 고유 다양체의 예. $k$를 체로 한다. $L \subset \mathbf{P}^3_k$를 직선으로, $C \subset \mathbf{P}^3_k$를 비특이 이차 곡선으로 한다. $C \cap L = \emptyset$라고 가정한다. 동형 $\varphi : L \to C$를 선택한다. $C$와 $L$를 $\varphi$에 의해 붙여서 얻어지는 $k$-다양체를 $X$라고 한다. 다시 말하면, 전사 고유 쌍유리 사상 $$
\pi : \mathbf{P}^3_k \longrightarrow X
$$과 열린 집합 $U \subset X$이 존재하며, $\pi : \pi^{-1}(U) \to U$은 동형이고, $\pi^{-1}(U) = \mathbf{P}^3_k \setminus (L \cup C)$이며, $\pi|_L = \pi|_C \circ \varphi$을 만족한다. (이 조건들만으로는 아직 $X$을 유일하게 정하지 않는다. 유일하게 정하려면, $Z = X \setminus U$의 점을 따라 $X$의 구조층을 지정할 필요가 있다.) $X$이 존재하며, 고유 다양체이지만 사영이 아님을 증명하라. (힌트: 존재에 대해서는 연습문제 \ref{exercises-exercise-prepare-glueing}의 결과를 사용하라. 비사영성에 대해서는 $\Pic(\mathbf{P}^3_k) = \mathbf{Z}$을 사용하여, $X$가 풍부한 가역 층을 가질 수 없음을 보일 수 있다.) \end{exercise}

\section{미분} \label{exercises-section-differentials}

\noindent {\bf 정의와 결과.} K\"ahler 미분. \begin{enumerate} \item $R \to A$를 환 준동형으로 한다. {\it $R$ 위의 $A$의 K\"ahler 미분 가군}을 $\Omega_{A/R}$라고 쓴다. 이것은 원래 $\text{d}a$,$a \in A$에서 생성되었으며, 다음 관계를 만족한다： $$
\text{d}(a_1 + a_2) = \text{d}a_1 + \text{d}a_2,\quad
\text{d}(a_1a_2) = a_1\text{d}a_2 + a_2\text{d}a_1,\quad
\text{d}r = 0
$$ 표준적인 보편 $R$-도함 $\text{d} : A \to \Omega_{A/R}$는 $a\mapsto \text{d}a$로 사상된다. \item 아이디얼 $I$를 정의하는 짧은 완전열 $$
0 \to I \to A \otimes_R A \to A \to 0
$$을 생각한다. $a$를 $a \otimes 1 - 1 \otimes a$의 동치류로 사상하는 표준적인 도함 $\text{d} : A \to I/I^2$가 있다. % P12-ERR-0546: preserve the frozen module-of-derivations wording.
이것은 $R$ 위의 $A$ 도함 가군의 또 다른 표현이다. 다시 말하면 $$
(I/I^2, \text{d}) \cong (\Omega_{A/R}, \text{d}).
$$ \item $S_R$의 상이 $S_A$에 포함되는 곱셈적 부분집합 $S_R \subset R$과 $S_A \subset A$에 대해 $$
\Omega_{S_A^{-1}A / S_R^{-1}R} =
S_A^{-1}\Omega_{A/R}.
$$가 성립한다. \item $A$이 유한 표현 $R$-대수라면, $\Omega_{A/R}$은 유한 표현 $A$-가군이다. 따라서 이 경우, $\Omega_{A/R}$의 {\it Fitting} 아이디얼이 정의된다. \item $f : X \to S$를 스킴의 사상으로 한다. 준연접 ${\mathcal O}_X$-가군의 층 $\Omega_{X/S}$과 % P12-ERR-0484: preserve the frozen O_S-linear wording.
${\mathcal O}_S$-선형 도함 $$
\text{d} : {\mathcal O}_X \longrightarrow \Omega_{X/S}
$$이 존재하며, $f(U) \subset V$을 만족하는 임의의 아핀 열린 집합 $\Spec(A) = U \subset X$, $\Spec(R) = V \subset S$에 대해 $$
\Gamma(\Spec(A), \Omega_{X/S}) = \Omega_{A/R}
$$이며, 이는 $\text{d}$와 양립한다. \end{enumerate}

\begin{exercise} \label{exercises-exercise-dual-numbers} $k[\epsilon]$를 체 $k$ 상의 쌍대수의 환, 즉 $\epsilon^2 = 0$로 하자. \begin{enumerate} \item 환 준동형 $$
R = k[\epsilon] \to A = k[x, \epsilon]/(\epsilon x)
$$를 생각하자. $\Omega_{A/R}$의 Fitting 아이디얼이 (제로 번째 Fitting 아이디얼부터 시작하여) $$
(\epsilon), A, A, \ldots
$$ 임을 보이라. \item 사상 $R = k[t] \to
A = k[x, y, t]/(x(y-t)(y-1), x(x-t))$을 생각하자. 간단하게 하기 위해 $k$의 표수는 0이라고 가정하고, $A$에서의 $\Omega_{A/R}$의 Fitting 아이디얼이 $$
x(2x-t)(2y-t-1)A, \ (x, y, t)\cap (x, y-1, t), \ A, \ A, \ldots
$$임을 보여라. 따라서, % P12-ERR-0485: preserve the frozen 0-the ordinal source locus.
제 $0$ Fitting 아이디얼은 $A$의 한 원소로 잘라지며, 제 $1$ Fitting 아이디얼은 $\Spec(A)$의 두 닫힌 점을 정의하며, 나머지는 모두 자명하다. \item 사상 $R = k[t] \to A = k[x, y, t]/(xy-t^n)$을 고려한다. $\Omega_{A/R}$의 Fitting 아이디얼을 계산하라. \end{enumerate} \end{exercise}

\begin{remark} \label{exercises-remark-fitting-omega-not-sings} % P12-ERR-0486: preserve the frozen alphabetic references to numeric items.
$X$가 $S$ 위에서 상대 차원 $k$을 가질 때, $\Omega_{X/S}$의 제 $k$ Fitting 아이디얼은 사상 $X \to S$의 특이 스킴을 정의하기 위해 일반적으로 사용된다. 그러나 (a)에서 보이듯이, 족의 올 차원이 '일정'하지 않은 경우, 예를 들어 평탄하지 않은 경우에는 이를 수행할 때 주의가 필요하다. (b)에서 보이듯이, 평탄성도 이것이 잘 될 것을 보장하지 않는다 (실제로, 이것은 평탄한 족이다). (c)와 같은 '좋은 경우'에는, 곡선의 족에 대해, 제 $0$ Fitting 아이디얼은 0이며, 제 $1$ Fitting 아이디얼은 특이점 집합을 (스킴론적으로) 정의할 것으로 기대된다. \end{remark}

\begin{exercise} \label{exercises-exercise-formally-smooth} $R$을 환이라 하고, $$
A = R[x_1, \ldots, x_n]/(f_1, \ldots, f_n).
$$로 한다. $A$은 변수와 같은 수의 방정식으로 표현되어 있음을 주의하라. 따라서 편도함수의 행렬 $$
( \partial f_i / \partial x_j )
$$은 $n \times n$ 행렬, 즉 정방행렬이다. 그 행렬식이 $A$의 원소로서 가역적이라고 가정한다. 이것은, $n$ 개의 생성원과 $n$ 개의 관계식을 갖는 이 경우에 $\Omega_{A/R} = (0)$ 임을 의미하는 조건이다. $\pi : B' \to B$를 $R$-대수의 전사로 하고, 그 핵 $J$은 ($B'$의 아이디얼로서) 제곱이 0이라고 하자. $\varphi : A \to B$을 $R$-대수의 준동형으로 하자. $\varphi = \pi \circ \varphi'$을 만족하는 $R$-대수의 준동형 $\varphi' : A \to B'$가 유일하게 존재함을 증명하라. \end{exercise}

\begin{exercise} \label{exercises-exercise-formally-smooth-one-equation} $A = R[x, y]/(f)$의 경우에, 연습 \ref{exercises-exercise-formally-smooth}의 결과를 일반화하라. \end{exercise}

\begin{exercise} \label{exercises-exercise-Jacobian-criterion} $k$를 체로 하고, $f_1, \ldots, f_c \in k[x_1, \ldots, x_n]$로 하며, $A = k[x_1, \ldots, x_n]/(f_1, \ldots, f_c)$로 한다. $f_j(0, \ldots, 0) = 0$라고 가정한다. 이는 $\mathfrak m = (x_1, \ldots, x_n)A$가 극대 아이디얼임을 의미한다. 행렬 $$
(\partial f_j/ \partial x_i)|_{(x_1, \ldots, x_n) = (0, \ldots, 0)}
$$의 계수가 $c$라면, 국소환 $A_\mathfrak m$가 정칙임을 증명하라. 이 경우의 $A_\mathfrak m$의 차원은 무엇인가. $A_\mathfrak m$는 정칙이지만 계수가 $c$ 미만이 되는 예를 들어, 역이 거짓임을 보여라. 그 예에서 $A_\mathfrak m$의 차원은 무엇인가. \end{exercise}


\section{스킴론, 2007년 가을 학기 기말시험} \label{exercises-section-final-exam-fall-2007}

\noindent % P12-ERR-0487: preserve the frozen Fall-heading/Spring-prose mismatch.
다음은 Columbia 대학교에서 2007년 봄에 진행된 스킴론 강의의 기말시험 문제이다.

\begin{exercise}[정의] \label{exercises-exercise-definitions} 다음 개념을 정의하라. \begin{enumerate} \item $X$는 {\it 스킴}이다. \item 스킴의 사상 $f : X \to Y$은 {\it 유한}이다. % P12-ERR-0488: preserve the frozen plural-subject grammar locus.
\item 스킴의 사상 $f : X \to Y$은 {\it 유한형}이다. \item 스킴 $X$은 {\it 뇌터}이다. \item 스킴 $X$ 위의 ${\mathcal O}_X$-가군 ${\mathcal L}$은 {\it 가역}이다. \item 대수적 폐체 위의 비특이 사영 곡선의 {\it 종수}. \end{enumerate} \end{exercise}

\begin{exercise} \label{exercises-exercise-kill-global-sections} $X = \Spec({\mathbf Z}[x, y])$으로 하며, ${\mathcal F}$를 준연접 ${\mathcal O}_X$-가군으로 한다. ${\mathcal F}$의 표준 아핀 열린 집합 $D(x)$에 대한 제한이 0이라고 가정한다. \begin{enumerate} \item ${\mathcal F}$의 임의의 대역 단면 $s$은 $x$의 어떤 거듭제곱으로 사라짐을, 즉, 어떤 $n\in {\mathbf N}$에 대해 $x^ns = 0$가 됨을 보이라. \item ${\mathcal F}$이 준연접이라고 가정하지 않는다면 동일한 주장이 성립한다고 생각하는가. \end{enumerate} \end{exercise}

\begin{exercise} \label{exercises-exercise-empty-fibre-empty} $X \to \Spec(R)$를 고유 사상으로 하고, $R$를 잉여체 $k$를 갖는 이산 값매김환으로 한다. $X \times_{\Spec(R)} \Spec(k)$가 공집합 스킴이라고 가정한다. $X$가 공집합 스킴임을 보이라. \end{exercise}

\begin{exercise} \label{exercises-exercise-curve-p1-p1} 복소수체 ${\mathbf C}$ 위의 사영 다양체 \footnote{사영 임베딩은 $((X_0, X_1), (Y_0, Y_1))\mapsto (X_0Y_0, X_0Y_1, X_1Y_0, X_1Y_1)$, 다시 말하면 $(x, y)\mapsto (1, y, x, xy)$이다.} $$
{\mathbf P}^1 \times {\mathbf P}^1 = {\mathbf P}^1_{{\mathbf C}}
\times_{\Spec({\mathbf C})} {\mathbf P}^1_{\mathbf C}
$$를 생각한다. 이것은, ${\mathbf P}^1_{\mathbf C}$의 표준 아핀 부분들의 조합에 대응하는 네 개의 아핀 부분으로 덮여 있다. 예를 들어, 첫 번째 인자에 동차 좌표 $X_0, X_1$, 두 번째 인자에 $Y_0, Y_1$을 사용한다고 하자. $x = X_1/X_0$, $y = Y_1/Y_0$라고 둔다. 이때 네 개의 아핀 열린 부분은 다음 환들의 스펙트럼이다: $$
{\mathbf C}[x, y], \quad
{\mathbf C}[x^{-1}, y], \quad
{\mathbf C}[x, y^{-1}], \quad
{\mathbf C}[x^{-1}, y^{-1}].
$$ $X \subset {\mathbf P}^1 \times {\mathbf P}^1$, 첫 번째 아핀 부분 내의 방정식 $$
y^3(x^4 + 1) = x^4 -1.
$$으로 주어지는 닫힌 부분집합의 폐포인 닫힌 부분 스킴으로 하자. \begin{enumerate} \item $X$가 네 개의 아핀 열린 부분 중 첫 번째와 마지막 합에 포함됨을 보이라. \item $X$가 비특이 사영 곡선임을 보이라. % P12-ERR-0489: preserve the frozen pr_2/first-factor mismatch.
\item 사상 $pr_2 : X \to {\mathbf P}^1$(첫 번째 인자로의 사상)을 고려하라. 첫 번째 아핀 부분에서는, 이것은 사상 $(x, y) \mapsto x$이다. 이것이 차수 $3$를 가지는 이유를 간략히 설명하라. \item 사상 $pr_2 : X \to {\mathbf P}^1$의 분기점과 분기 지수를 계산하라. \item $X$의 종수를 계산하라. \end{enumerate} \end{exercise}

\begin{exercise} \label{exercises-exercise-finite-type-over-Z} $X \to \Spec({\mathbf Z})$를 유한형 사상으로 한다. $X \times_{\Spec({\mathbf Z})} \Spec({\mathbf F}_p)$이 공집합이 아닌 소수 $p$가 무한히 존재한다고 가정한다. \begin{enumerate} \item $X \times_{\Spec({\mathbf Z})}\Spec(\mathbf{Q})$가 비어 있지 않음을 보여라. \item 조건 '유한형'을 조건 '국소유한형'으로 대체한 경우에도, 같은 주장이 성립한다고 생각하는가. \end{enumerate} \end{exercise}


\section{스킴론, 2009년 봄 학기말 시험} \label{exercises-section-final-exam-spring-2009}

\noindent 아래는 Columbia 대학교에서 2009년 봄에 진행된 스킴론 강의의 학기말 시험 문제이다.

\begin{exercise} \label{exercises-exercise-Noetherian-coherent} $X$를 뇌터 스킴이라고 한다. $\mathcal{F}$를 $X$ 위의 연접층이라고 한다. $x \in X$을 한 점이라고 한다. $\text{Supp}(\mathcal{F}) = \{ x \}$이라고 가정한다. \begin{enumerate} \item $x$가 $X$의 닫힌 점임을 보여라. \item $H^0(X, \mathcal{F})$가 0이 아님을 보여라. \item $\mathcal{F}$가 대역 단면으로 생성됨을 보여라. \item $p > 0$에 대해 $H^p(X, \mathcal{F}) = 0$임을 보여라. \end{enumerate} \end{exercise}

\begin{remark} \label{exercises-remark-invertible-projective-space} $k$를 체로 한다. $\mathbf{P}^2_k = \text{Proj}(k[X_0, X_1, X_2])$라 하자. $\mathbf{P}^2_k$ 위의 임의의 가역층은 어떤 $n \in \mathbf{Z}$에 대한 $\mathcal{O}_{\mathbf{P}^2_k}(n)$와 동형이다. 다음을 기억하자: $$
\Gamma(\mathbf{P}^2_k, \mathcal{O}_{\mathbf{P}^2_k}(n)) =
k[X_0, X_1, X_2]_n
$$는 다항식환의 차수 $n$의 부분이다. $\mathbf{P}^2_k$ 위의 준연접 층 $\mathcal{F}$에 대해, 통상대로 $\mathcal{F}(n) =
\mathcal{F}
\otimes_{\mathcal{O}_{\mathbf{P}^2_k}}
\mathcal{O}_{\mathbf{P}^2_k}(n)$라고 놓자. \end{remark}

\begin{exercise} \label{exercises-exercise-nonsplit-vectorbundle} $k$를 체로 한다. $\mathcal{E}$를 $\mathbf{P}^2_k$ 위의 벡터 다발, 즉 유한 국소 자유 $\mathcal{O}_{\mathbf{P}^2_k}$-가군이라고 하자. % P12-ERR-0490: preserve the frozen missing-of source locus.
$\mathcal{E}$가 가역 $\mathcal{O}_{\mathbf{P}^2_k}$-가군의 직합과 동형일 때, $\mathcal{E}$를 {\it 분해된다}고 한다. \begin{enumerate} \item $\mathcal{E}$가 분해됨과 $\mathcal{E}(n)$가 분해됨이 동치임을 보이라. \item $\mathcal{E}$이 분해된다면, 모든 $n \in \mathbf{Z}$에 대해 $H^1({\mathbf{P}^2_k}, \mathcal{E}(n)) = 0$임을 보여라. \item 선형형식 $L_0, L_1, L_2 \in \Gamma(\mathbf{P}^2_k, \mathcal{O}_{\mathbf{P}^2_k}(1))$에 의해 주어지는 $$
\varphi :
\mathcal{O}_{\mathbf{P}^2_k}
\longrightarrow
\mathcal{O}_{\mathbf{P}^2_k}(1) \oplus
\mathcal{O}_{\mathbf{P}^2_k}(1) \oplus
\mathcal{O}_{\mathbf{P}^2_k}(1)
$$를 고려한다. 어떤 $i$에 대해 $L_i \not = 0$라고 가정한다. $\varphi$의 여핵이 벡터다발이 되기 위한 $L_0, L_1, L_2$에 관한 조건은 무엇인가. 또한, 그것은 왜인가. % P12-ERR-0491: preserve the frozen Given/Give imperative locus.
\item 그러한 $\varphi$의 예를 하나 들어라. \item $\Coker(\varphi)$이 벡터 다발이라면, 그것이 분해되지 않음을 보이시오. \end{enumerate} \end{exercise}

\begin{remark} \label{exercises-remark-recall-dimension-theory}다음 차원론의 사실들을 자유롭게 사용해도 좋다 (필요하다면, 추가적인 사실을 더해도 좋다). \begin{enumerate} \item 스킴의 차원은, 기약 닫힌 부분집합 사슬의 길이의 상한이다. \item 체 위 유한 형식 스킴의 차원은, 그것의 아핀 열린집합의 차원의 최댓값이다. \item 뇌터 스킴의 차원은 그 기약 성분들의 차원의 최댓값이다. \item 아핀 스킴의 차원은 대응하는 환의 차원과 일치한다. \item $k$를 체로 하고, $A$을 유한형 $k$-대수라고 하자. % P12-ERR-0492: preserve the frozen unbound-x source locus.
$A$가 정역이고 $x \not = 0$라면, $\dim(A) = \dim(A/xA) + 1$이다. \end{enumerate} \end{remark}

\begin{exercise} \label{exercises-exercise-irreducible-fibres-same-dimension-irreducible} $k$를 체로 한다. $X$을 $k$ 위의 사영 축소 스킴이라고 한다. $f : X \to \mathbf{P}^1_k$을 $k$ 위의 스킴의 사상이라고 한다. 모든 점 $t \in \mathbf{P}^1_k$에 대해, 올 $X_t = f^{-1}(t)$가 차원 $d$인 기약 공간이 되도록 하는 정수 $d \geq 0$가 존재한다고 가정한다 (기약 공간은 공집합이 아님을 기억하자). \begin{enumerate} \item $\dim(X) = d + 1$임을 보여라. \item $X_0 \subset X$를 차원 $d + 1$의 $X$의 기약 성분이라고 하자. 모든 $t \in \mathbf{P}^1_k$에 대해, 올 $X_{0, t}$의 차원이 $d$임을 증명하라. \item 위의 결과로부터 $X_t$와 $X_{0, t}$에 대해 어떤 결론을 내릴 수 있는가. \item $X$가 기약임을 보여라. \end{enumerate} \end{exercise}

\begin{remark} \label{exercises-remark-chi} 체 $k$ 위의 사영 스킴 $X$과 $X$ 위의 연접층 $\mathcal{F}$에 대해 $$
\chi(X, \mathcal{F}) =
\sum\nolimits_{i \geq 0} (-1)^i\dim_k H^i(X, \mathcal{F}).
$$로 둔다. \end{remark}

\begin{exercise} \label{exercises-exercise-complete-intersection} $k$를 체로 한다. $\mathbf{P}^3_k = \text{Proj}(k[X_0, X_1, X_2, X_3])$라고 쓴다. $C \subset \mathbf{P}^3_k$을 {\it 형 $(5, 6)$의 완전 교차 곡선}으로 한다. 이것은 $F \in k[X_0, X_1, X_2, X_3]_5$과 $G \in k[X_0, X_1, X_2, X_3]_6$이 존재하며, $$
C = \text{Proj}(k[X_0, X_1, X_2, X_3]/(F, G))
$$가 차원 $1$의 다양체가 됨을 의미한다 (다양체이므로 축소이고 기약하지만, 필요하다면 $C$가 비특이적이라고 가정할 수 있다). $i : C \to \mathbf{P}^3_k$을 대응하는 닫힌 몰입으로 한다. 완전 교차이므로, 추가로 층의 완전 열 $$
\xymatrix{
0 \ar[r] &
\mathcal{O}_{\mathbf{P}^3_k}(-11)
\ar[r]^-{
\left(
\begin{matrix}
-G \\
F
\end{matrix}
\right)
} &
\mathcal{O}_{\mathbf{P}^3_k}(-5) \oplus \mathcal{O}_{\mathbf{P}^3_k}(-6)
\ar[r]^-{(F, G)} &
\mathcal{O}_{\mathbf{P}^3_k} \ar[r] &
i_*\mathcal{O}_C \ar[r] &
0
}
$$도 얻을 수 있다. 이 사실들을 사용하여 다음을 수행하라: \begin{enumerate} \item 임의의 $n \in \mathbf{Z}$에 대해 $\chi(C, i^*\mathcal{O}_{\mathbf{P}^3_k}(n))$를 계산하라; \item $H^1(C, \mathcal{O}_C)$의 차원을 계산하라. \end{enumerate} \end{exercise}

\begin{exercise} \label{exercises-exercise-glueing} $k$를 체로 하자. 환 \begin{align*}
A & = k[x, y]/(xy) \\
B & = k[u, v]/(uv) \\
C & = k[t, t^{-1}] \times k[s, s^{-1}]
\end{align*}과, $k$-대수 준동형 $$
\begin{matrix}
A \longrightarrow C, &
x \mapsto (t, 0), &
y \mapsto (0, s) \\
B \longrightarrow C, &
u \mapsto (t^{-1}, 0), &
v \mapsto (0, s^{-1})
\end{matrix}
$$를 고려하라. 이들 사상은 동형 $A_{x + y} \to C$ 및 $B_{u + v} \to C$를 유도한다. 따라서, 사상 $A \to C$ 및 $B \to C$는 $\Spec(C)$를 $\Spec(A)$ 및 $\Spec(B)$의 열린 부분집합으로 동일시한다. $\Spec(A)$과 $\Spec(B)$를 $\Spec(C)$에 따라 붙여 얻어지는 스킴을 $X$ 라 하자: $$
X = \Spec(A) \amalg_{\Spec(C)} \Spec(B).
$$ 강의에서 본 것처럼, 이러한 스킴은 존재하며, 거기에는 아핀 열린 집합 $\Spec(A) \subset X$와 $\Spec(B) \subset X$가 있고, 그것들의 공통 부분은 정확히 $\Spec(C)$이며, 위의 사상을 사용하여 각각의 열린 부분과 동일시된다. \begin{enumerate} \item 왜 $X$는 분리되어 있는가. \item 왜 $X$는 $k$ 위에서 유한형인가. \item $H^1(X, \mathcal{O}_X)$를 계산하라. 또는, 그 차원은 몇인가. \item $X$를 기술하는, 더 기하학적인 방법은 무엇인가. \end{enumerate} \end{exercise}

\section{스킴 이론, 2010년 가을 학기 기말시험} \label{exercises-section-final-exam-fall-2010}

\noindent 아래는 콜럼비아 대학교에서 2010년 가을에 진행된 스킴 이론 강의의 기말시험 문제이다.

\begin{exercise}[정의] \label{exercises-exercise-definitions-fall-2010} 다음 개념을 정의하라. \begin{enumerate} \item 분리 스킴, \item 스킴의 준콤팩트 사상, \item 스킴의 아핀 사상, \item 환의 곱셈적 부분집합, \item 뇌터 스킴, \item 다양체. \end{enumerate} \end{exercise}

\begin{exercise} \label{exercises-exercise-prime-avoidance} 소 아이디얼 회피. \begin{enumerate} \item $A$을 환으로 한다. $I \subset A$를 아이디얼로 하고, $\mathfrak q_1$, $\mathfrak q_2$를 $I \not \subset \mathfrak q_i$를 만족하는 소아이디얼로 하자. $I \not \subset \mathfrak q_1 \cup \mathfrak q_2$임을 보이라. \item (1)의 기하학적 해석은 무엇인가. \item $X = \text{Proj}(S)$라 하고, $S$를 어떤 등급환이라고 하자. $x_1, x_2 \in X$라 하자. $x_1$와 $x_2$를 모두 포함하는 표준 열린 집합 $D_{+}(F)$가 존재함을 보이시오. \end{enumerate} \end{exercise}

\begin{exercise} \label{exercises-exercise-compose-affine} 아핀 사상의 합성이 아핀인 이유는 무엇인가. \end{exercise}

\begin{exercise}[예] \label{exercises-exercise-examples} 다음 예를 들어라: \begin{enumerate} \item 체 $k$ 위의 가약 사영 스킴. \item 100개의 점을 가진 스킴. \item 스킴의 아핀이 아닌 사상. \end{enumerate} \end{exercise}

\begin{exercise} \label{exercises-exercise-chevalley-hilbert-nullstellensatz} Chevalley의 정리와 Hilbert의 영점 정리. \begin{enumerate} \item $\mathfrak p \subset \mathbf{Z}[x_1, \ldots, x_n]$를 극대 아이디얼이라고 하자. 셰발리의 정리는 $\mathfrak p \cap \mathbf{Z}$에 대해 무엇을 의미하는가. \item 또한, 힐베르트의 영점 정리는 $\kappa(\mathfrak p)$에 대해 무엇을 의미하는가. \end{enumerate} \end{exercise}

\begin{exercise} \label{exercises-exercise-P0} $A$을 환으로 한다. $X$의 차수를 $1$라고 하여 등급 $A$-대수로서 $S = A[X]$라고 둔다. $A$ 위의 스킴으로서 $\text{Proj}(S) \cong \Spec(A)$임을 보여라. \end{exercise}

\begin{exercise} \label{exercises-exercise-finite-is-projective} $A \to B$을 유한한 환 준동형으로 하자. $\Spec(B)$가 $\Spec(A)$ 위의 H-프로젝티브 스킴임을 보여라. \end{exercise}

\begin{exercise} \label{exercises-exercise-not-geometrically-irreducible} 체 $k$ 위의 스킴 $X$이고, $X$는 기약이지만, 어떤 유한 확장 $k'/k$에 대한 밑변환 $X_{k'} = X \times_{\Spec(k)} \Spec(k')$이 연결이고 가약이 되는 예를 제시하라. \end{exercise}


\section{스킴 이론, 2011년 봄 학기말 시험} \label{exercises-section-final-exam-spring-2011}

\noindent 이하 내용은 컬럼비아 대학교에서 2011년 봄에 진행된 스킴 이론 강의의 학기말 시험 문제이다.

\begin{exercise}[정의] \label{exercises-exercise-definitions-spring-2011} 이탤릭체로 표시된 개념을 정의하시오. \begin{enumerate} \item {\it 분리된} 스킴, \item 스킴의 {\it 보편적 폐포} 사영, \item 공통의 체에 포함된 국소 환 $A, B$에 대한 {\it $A$은 $B$을 지배하는}, \item 스킴 $X$의 {\it 차원}, \item 스킴 $X$의 기존 닫힌 부분 스킴 $Y$의 {\it 여차원}. \end{enumerate} \end{exercise}

\begin{exercise}[결과] \label{exercises-exercise-results-spring-2011} 강의에서 논의한 사실과 형식적으로 동등한 명제를 서술하라. \begin{enumerate} \item 예를 들어, 다양체의 사상 $X \to Y$에 대한 고유성의 부여 판정법. \item 체 $k$ 위의 다양체 $X$에 대한 $\dim(X)$과 $X$의 함수체 $k(X)$의 관계. \item 빈칸을 채우세요: $k$ 위의 비특이 사영 곡선과 비상수 사상의 범주는, $\ldots\ldots\ldots$와 반동등이다. \item 뇌터 정규화. \item Jacobian 판정법. \end{enumerate} \end{exercise}

\begin{exercise} \label{exercises-exercise-genus-plane-curve} $k$를 체로 한다. $F \in k[X_0, X_1, X_2]$를 차수 $d$의 동차 형식으로 한다. $C = V_{+}(F) \subset \mathbf{P}^2_k$를 $k$ 위의 매끄러운 곡선이라고 가정한다. % P12-ERR-0493: preserve the frozen missing-by declaration locus.
$i : C \to \mathbf{P}^2_k$를 대응하는 닫힌 매입으로 한다. \begin{enumerate} \item $\mathbf{P}^2_k$ 위의 연접층의 짧은 완전 열 $$
0 \to \mathcal{O}_{\mathbf{P}^2_k}(-d)
\to \mathcal{O}_{\mathbf{P}^2_k} \to i_*\mathcal{O}_C \to 0
$$이 존재함을 보여라: 사상이 무엇인지 말하고, 왜 완전한지 간단히 설명하라. \item $H^0(C, \mathcal{O}_C) = k$임을 결론지어라. \item $C$의 종 수를 계산하라. \item 이번에는 $P = (0 : 0 : 1)$가 $C$ 위에 없다고 가정한다. $(a_0 : a_1 : a_2) \mapsto (a_0 : a_1)$로 주어지는 $\pi : C \to \mathbf{P}^1_k$의 차수가 $d$임을 증명하라. \item $k$은 대수적으로 닫혀 있고, 모든 분기 지수(「$e_i$」)는 $1$ 또는 $2$이며, $k$의 특성은 $2$가 아니라고 가정한다. $\pi : C \to \mathbf{P}^1_k$은 분기점을 몇 개 가지고 있는가. \item $F$를 사용하여, $\pi$의 분기점 집합을 정하는 방정식계는 무엇이라고 생각하는가. \item $K_C$를 추측할 수 있는가. \end{enumerate} \end{exercise}

\begin{exercise} \label{exercises-exercise-Pic-triangle} $k$를 몸으로 한다. $X$를 $k$ 위의 '삼각형'으로 한다. 즉, $\mathbf{A}^1_k$의 세 개의 복사본을, 첫 번째 복사본의 $0$와 두 번째 복사본의 $1$, 두 번째 복사본의 $0$과 세 번째 복사본의 $1$, 세 번째 복사본의 $0$과 첫 번째 복사본의 $1$을 각각 동일시하여 서로 맞붙임으로써 $X$를 얻는다. $X$는 $\Spec(k[x, y]/(xy(x + y + 1)))$와 동형이 된다; 이는 자유롭게 사용할 수 있다. $X$의 Picard 군을 계산하라. \end{exercise}

\begin{exercise} \label{exercises-exercise-birational-morphism-curves-ample} $k$를 몸체로 한다. $\pi : X \to Y$를 곡선의 유한 쌍유리사상으로 하고, $X$는 $k$ 위의 비특이 사영 곡선으로 한다. 강의 내용으로부터, $Y$는 고유곡선이며, $\pi$는 $Y$의 정규화 사상이다. 또한, 조밀 열린 집합 $V \subset Y$가 존재하며, $U = \pi^{-1}(V)$은 $X$의 조밀 열린 집합이고 $\pi : U \to V$는 동형이 되는 것이 존재한다는 것도 강의에서 보았다. \begin{enumerate} \item 유효 Cartier 인자 $D \subset X$이고, $D \subset U$ 및 $\mathcal{O}_X(D)$가 $X$ 위에서 풍부한 것이 존재함을 보이라. \item $D$를 (1)과 같은 것으로 하자. $E = \pi(D)$가 $Y$ 위의 유효 Cartier 인자임을 보이라.\item 다음 이유를 간단히 서술하라： \begin{enumerate} \item 사상 $\mathcal{O}_Y \to \pi_*\mathcal{O}_X$은 $Y \setminus V$에 지지를 가진 연결된 여핵 $Q$을 가진다；\item 모든 $n$에 대하여, 여핵이 $Q$와 동형이 되는 대응 사상 $\mathcal{O}_Y(nE) \to \pi_*\mathcal{O}_X(nD)$이 존재한다. \end{enumerate} \item $\dim_k H^0(X, \mathcal{O}_X(nD)) - \dim_k H^0(Y, \mathcal{O}_Y(nE))$가 유계임을 (무엇에 의하여) 보이고, $n$가 충분히 클 때 가역층 $\mathcal{O}_Y(nE)$은 여러 단절을 가지고 있음을 결론지어라(왜인가). \end{enumerate} \end{exercise}



\section{스킴론, 2011년 가을 학기말 시험} \label{exercises-section-final-exam-fall-2011}

\noindent % P12-ERR-0547: preserve the frozen Schemes-heading/Commutative-Algebra-prose mismatch.
다음은 콜럼비아 대학교에서 2011년 가을에 진행된 가환대수 강의의 학기말 시험 문제이다.

\begin{exercise}[정의] \label{exercises-exercise-definitions-fall-2011} 이탤릭체로 표시된 개념을 정의하시오. \begin{enumerate} \item {\it 뇌터} 환, \item {\it 뇌터} 스킴, \item {\it 유한}한 환 준동형, \item 스킴의 {\it 유한} 사상, \item 환의 {\it 차원}. \end{enumerate} \end{exercise}

\begin{exercise}[결과] \label{exercises-exercise-results-fall-2011} 강의에서 논한 사실과 형식적으로 동치인 명제를 서술하시오. \begin{enumerate} \item Zariski의 주정리. \item 뇌터 정규화. \item 중국 나머지 정리. \item 유한한 환 준동형에 대한 상승 정리. \end{enumerate} \end{exercise}

\begin{exercise} \label{exercises-exercise-dimension-of-ring} $(A, \mathfrak m, \kappa)$를, 나머지 체의 특징이 $2$이 아닌 뇌터 국소환으로 한다. $\mathfrak m$는 세 개의 원소 $x, y, z$로 생성되며, $A$에서 $x^2 + y^2 + z^2 = 0$이라고 가정한다. \begin{enumerate} \item $\dim(A)$의 가능한 값은 무엇인가. \item 각 값이 가능함을 보여주는 예를 제시하라. \item $\dim(A) = 2$라면 $A$이 정역(domain)임을 보여라. (힌트: $\bigoplus_{n \geq 0} \mathfrak m^n/\mathfrak m^{n + 1}$를 조사하라. ) \end{enumerate} \end{exercise}

\begin{exercise} \label{exercises-exercise-localization} $A$를 환으로 한다. $S \subset T \subset A$를 곱셈적 부분집합으로 한다. $$
\{\mathfrak q \mid \mathfrak q \cap S = \emptyset\} =
\{\mathfrak q \mid \mathfrak q \cap T = \emptyset\}.
$$라고 가정한다. $S^{-1}A \to T^{-1}A$가 동형임을 보이라. \end{exercise}

\begin{exercise} \label{exercises-exercise-locus-of-rank-1} $k$를 대수적 폐체로 한다. $$
V_0 = \{ A \in \text{Mat}(3 \times 3, k) \mid \text{rank}(A) = 1\}
\subset \text{Mat}(3 \times 3, k) = k^9.
$$라 하자. \begin{enumerate} \item $V_0$가, (Zariski) 국소폐포집합 $V \subset \mathbf{A}^9_k$의 폐점 전체임을 보여라. \item $V$은(는) 이미약한가. \item $\dim(V)$은 무엇인가. \end{enumerate} \end{exercise}

\begin{exercise} \label{exercises-exercise-not-complete-intersection} $\mathbf{C}[x, y]$의 아이디얼 $(x^2, xy, y^2)$이 $2$ 개의 원소로 생성될 수 없음을 증명하라. \end{exercise}

\begin{exercise} \label{exercises-exercise-finite-projection} $f \in \mathbf{C}[x, y]$를 비정수 다항식이라고 하자. 어떤 $\alpha, \beta \in \mathbf{C}$에 대하여 $\mathbf{C}$-대수 준동형 $$
\mathbf{C}[t] \longrightarrow \mathbf{C}[x, y]/(f),\quad
t \longmapsto \alpha x + \beta y
$$가 유한함을 보여라. \end{exercise}

\begin{exercise} \label{exercises-exercise-union-of-two-affines} 유한한 수의 점 $p_1, \ldots, p_n \in \mathbf{C}^2$이 주어졌을 때, 스킴 $\mathbf{A}^2_\mathbf{C} \setminus \{p_1, \ldots, p_n\}$이 두 개의 아핀 열린 집합의 합임을 보여라. \end{exercise}

\begin{exercise} \label{exercises-exercise-surjection-A1-P1} 스킴의 전사 $\mathbf{A}^1_\mathbf{C} \to \mathbf{P}^1_\mathbf{C}$가 존재함을 보이시오. (전사란 단순히 기반이 되는 점 집합 위에서 전사임을 의미한다.) \end{exercise}

\begin{exercise} \label{exercises-exercise-almost-surjective} $k$를 대수적으로 닫힌 체라고 하자. $A \subset B$를, 둘 다 유한형 $k$-대수인 정역의 확장으로 하자. 다음 절차를 사용하여, $\Spec(B) \to \Spec(A)$의 상이 $\Spec(A)$의 공집합이 아닌 열린 부분집합을 포함함을 증명하라: \begin{enumerate} \item $A \to B$도 유한한 경우에 증명하라. \item $B$의 분수체가 $A$의 분수체의 유한 확장인 경우에 증명하라. \item 주장을 이전 경우로 귀착시켜라. \end{enumerate} \end{exercise}



\section{스킴, 2013년 가을학기 기말시험}
\label{exercises-section-final-exam-fall-2013}

\noindent
% P12-ERR-0548: preserve the frozen Schemes-heading/Commutative-Algebra-prose mismatch.
다음은 Columbia 대학교에서 2013년 가을에 열린 가환대수학 강의의
기말시험 문제들이다.

\begin{exercise}[정의]
\label{exercises-exercise-definitions-fall-2013}
이탤릭체로 표시한 개념들을 정의하라.
\begin{enumerate}
\item 환의 {\it 근기 아이디얼},
\item {\it 유한형} 환 준동형,
\item {\it Weil 미분},
\item {\it 스킴}.
\end{enumerate}
\end{exercise}

\begin{exercise}[결과]
\label{exercises-exercise-results-fall-2013}
강의에서 논의한 각 사실과 형식적으로 동치인 명제를 진술하라.
\begin{enumerate}
\item 등급 가군의 Hilbert 다항식에 관한 결과.
\item 뇌터 국소환 $(R, \mathfrak m)$과
$\bigoplus_{n \geq 0} \mathfrak m^n/\mathfrak m^{n + 1}$의 차원.
\item Riemann--Roch 정리.
\item Clifford 정리.
\item Chevalley 정리.
\end{enumerate}
\end{exercise}

\begin{exercise}
\label{exercises-exercise-surjective-after-localization}
$A \to B$를 환 사상이라 하고, $S \subset A$를 곱셈적 부분집합이라 하자.
$A \to B$가 유한형이고 $S^{-1}A \to S^{-1}B$가 전사라고 가정하자.
$A_f \to B_f$가 전사가 되는 $f \in S$가 존재함을 보여라.
\end{exercise}

\begin{exercise}
\label{exercises-exercise-injective-local-ring-map-not-surjective}
국소환들의 단사 국소 준동형 $A \to B$로서
$\Spec(B) \to \Spec(A)$가 전사가 아닌 예를 들어라.
\end{exercise}

\begin{situation}[평면 곡선의 표기]
\label{exercises-situation-curve-in-the-plane}
$k$를 대수적으로 닫힌 체라 하자.
$F(X_0, X_1, X_2) \in k[X_0, X_1, X_2]$를 차수 $d$인 기약 동차
다항식이라 하자.
$$
D = V(F) \subset \mathbf{P}^2
$$
를 $F$가 영이 되는 것으로 주어지는 사영 평면 곡선이라 하자.
$x = X_1/X_0$, $y = X_2/X_0$라 놓고
$f(x, y) = X_0^{-d}F(X_0, X_1, X_2) = F(1, x, y)$라 놓자.
% P12-ERR-0494: preserve the frozen missing-by declaration locus.
$K$를 정역 $k[x, y]/(f)$의 분수체라 하자. $C$를 $K$에 대응하는 추상
곡선이라 하자. 강의에서 보았듯이 전사 사상 $C \to D$가 존재하며,
이는 $D$의 비특이 궤적 위에서 전단사이고 $D$가 비특이면 동형이다.
$f_x = \partial f/\partial x$, $f_y = \partial f/\partial y$라 놓자.
% P12-ERR-0549: preserve the frozen repeated missing-by denotation loci.
마지막으로 강의에서 논의한 $\Omega_{K/k}$의 원소
$\omega = \text{d}x/f_y = - \text{d}y/f_x$를 생각하자. $K_C$를
$\omega$의 영점과 극점의 인자라 하자.
\end{situation}

\begin{exercise}
\label{exercises-exercise-node-in-the-plane}
상황 \ref{exercises-situation-curve-in-the-plane}에서 $d \geq 3$이고, 곡선 $D$가
특이점 하나, 즉 $P = (1 : 0 : 0)$만을 갖는다고 하자. 더 나아가
$P = (0, 0)$ 근방에서 전개식
$$
f(x, y) = xy + h.o.t
$$
를 갖는다고 하자. 그러면 $C$는 $P$ 위에 두 점 $v$, $w$를 가지며,
이들은
$$
v(x) = 1, v(y) > 1
\quad\text{and}\quad
w(x) > 1, w(y) = 1
$$
로 특징지어진다.
\begin{enumerate}
\item $\Omega_{K/k}$의 원소
$\omega = \text{d}x/f_y = - \text{d}y/f_x$가 $v$와 $w$ 모두에서
1차 극을 가짐을 보여라. ($\omega$가 비특이점에서 보이는 거동은
강의에서 논의한 바와 같다.)
\item 강의에서 $C \to D$를 따라 $C$로 당겨 온 인자 $X_0 = 0$에서
$\omega$가 차수 $d - 3$으로 영이 됨을 보였다. 이 사실과 (1)의 정보를
합치면 $C$ 위에서 $\omega$의 영점과 극점으로 이루어진 인자의 차수는
얼마인가?
\item 곡선 $C$의 종수는 얼마인가?
\end{enumerate}
\end{exercise}

\begin{exercise}
\label{exercises-exercise-smooth-plane-curve-linear-system}
상황 \ref{exercises-situation-curve-in-the-plane}에서 $d = 5$이고 곡선 $C = D$가
비특이라고 하자. 강의에서 $C$의 종수가 $6$이고 선형계 $K_C$가
$$
L(K_C) = \{h\omega \mid h \in k[x, y],\ \deg(h) \leq 2\}
$$
로 주어짐을 보였다. 여기서 $\deg$는 전체 차수를 뜻한다
\footnote{$d - 3 = 5 - 3 = 2$이므로 $\leq 2$를 얻는다.}.
$P_1, P_2, P_3, P_4, P_5 \in D$를 아핀 열린집합 $X_0 \not = 0$에 속하는
서로 다른 점들이라 하자.
% P12-ERR-0550: preserve the frozen missing-by divisor denotation locus.
$\sum P_i = P_1 + P_2 + P_3 + P_4 + P_5$를 $C$의 대응하는 인자라 하자.
\begin{enumerate}
\item $L(K_C - \sum P_i)$를 다항식으로 기술하라.
\item $l(\sum P_i)$에는 어떤 가능성이 있는가?
\end{enumerate}
\end{exercise}

\begin{exercise}
\label{exercises-exercise-rational-curve-high-degree}
상황 \ref{exercises-situation-curve-in-the-plane}에서와 같은 $F$ 중에서 $d = 100$이고
$C$의 종수가 $0$이 되는 것을 하나 써라.
\end{exercise}

\begin{exercise}
\label{exercises-exercise-high-degree-curve-quadratic-equation}
$k$를 대수적으로 닫힌 체라 하자.
% P12-ERR-0495: preserve the frozen missing-article source locus.
$K/k$를 초월차수가 $1$인 유한 생성 체 확대라 하자. $C$를 $K$에 대응하는
추상 곡선이라 하자. $V \subset K$를 $g^r_d$라 하고,
$\Phi : C \to \mathbf{P}^r$를 이에 대응하는 사상이라 하자. $V$가 완전
선형계이고 $d$가 $C$의 종수에 비하여 충분히 크면 $C$의 상이 이차곡면
\footnote{이차곡면은 차수 $2$인 초곡면, 즉 $\mathbf{P}^r$에서 차수 $2$인
동차 다항식의 영점 집합이다.}에 포함됨을 보여라.
(추가 점수: 필요한 차수에 대한 좋은 상계를 구하라.)
\end{exercise}

\begin{exercise}
\label{exercises-exercise-surjective-map-a2-p2}
상황 \ref{exercises-situation-curve-in-the-plane}의 표기를 사용하자.
$U \subset \mathbf{P}^2_k$를 여집합이 $D$인 열린 부분스킴이라 하자.
$k$-대수 $A = \mathcal{O}_{\mathbf{P}^2_k}(U)$를 기술하라. $k$-대수로서
$A$를 생성하는 데 필요한 생성원 개수의 상계를 주어라.
\end{exercise}



\section{스킴 논, 2014년 봄 학기 기말 시험} \label{exercises-section-final-exam-spring-2014}

\noindent 아래는 콜럼비아 대학교에서 2014년 봄에 진행된 스킴 논 강의의 기말 시험 문제이다.

\begin{exercise}[정의] \label{exercises-exercise-definitions-spring-2014} $(X, \mathcal{O}_X)$를 스킴으로 한다. 이탤릭체로 표시된 개념을 정의하시오. \begin{enumerate} \item {\it 점 $x$에서의 $X$ 국소환}, \item $\mathcal{O}_X$-가군의 {\it 준연접}층, % P12-ERR-0496: preserve the frozen unclosed-parenthesis source locus.
\item $\mathcal{O}_X$-가군의 {\it 연결}층 ($X$은 국소 뇌터인 것으로 가정해도 좋다), \item $X$의 {\it 아핀 열린 집합}, \item {\it 스킴의 유한 사상} $X \to Y$. \end{enumerate} \end{exercise}

\begin{exercise}[정리] \label{exercises-exercise-results-spring-2014} 각 항목과 관련된, 강의에서 논의한 비자명한 사실을 정확히 서술하라. \begin{enumerate} \item 다양체의 다중종 수의 쌍유리 불변성에 대하여, \item 아핀 사상임이 국소적 성질임에 대하여, \item 사상의 스킴론적 섬유의 위상에 대하여, \item 고유성의 값 판단법. \end{enumerate} \end{exercise}

\begin{exercise} \label{exercises-exercise-miss-curve} $\mathbf{C}$을 복소수체로 하고, $X = \mathbf{A}^2_\mathbf{C}$라 하자. {\it 직선}이란, 어떤 $a, b, c \in \mathbf{C}$에 대해 $(a, b) \not = (0, 0)$를 만족하는 것에 대해, 하나의 일차 방정식 $ax + by + c = 0$로 정의되는 $X$의 닫힌 부분 스킴을 의미하는 것으로 하자. {\it 곡선}이란, 차원 $1$의 기약(따라서 비어 있지 않은) 닫힌 부분 스킴 $C \subset X$을 의미하는 것으로 하자. {\it 이차 곡선}이란, 어떤 $a, b, c, d, e, f \in \mathbf{C}$에서 $(a, b, c) \not = (0, 0, 0)$을 만족하는 것에 대해, 하나의 이차 방정식 $ax^2 + bxy + cy^2 + dx + ey + f = 0$으로 정의되는 곡선을 의미하는 것으로 하자. \begin{enumerate} \item 모든 직선과 $C$의 공통 부분이 비어 있지 않은 곡선 $C$을 찾아라. \item 모든 직선 및 모든 이차곡선과 $C$의 공통 부분이 공집합이 아닌 곡선 $C$를 찾아라. \item 모든 곡선 $C$에 대하여, $C \cap C' = \emptyset$가 되는 다른 곡선이 존재함을 증명하라. \end{enumerate} \end{exercise}

\begin{exercise} \label{exercises-exercise-normal-bundle-exceptional-curve} $k$를 체로 한다. $b : X \to \mathbf{A}^2_k$를 아핀 평면의 원점에서의 블로우업으로 한다. 다시 말하면, $\mathbf{A}^2_k = \Spec(k[x, y])$라면, $X = \text{Proj}(\bigoplus_{n \geq 0} \mathfrak m^n)$이다. 여기서 $\mathfrak m = (x, y) \subset k[x, y]$라 하자. 다음 주장을 증명하라: \begin{enumerate} \item 원점 위의 $b$의 스킴 이론적 섬 $E$은 $\mathbf{P}^1_k$와 동형이다; \item $E$은 $X$ 위의 유효한 Cartier 인자이다; \item $\mathcal{O}_X(-E)$의 $E$에 대한 제한은 차수 $1$의 직선 다발이다. \end{enumerate} ($\mathcal{O}_X(-E)$가 $X$에서 $E$의 아이디어 층임을 상기하자.) \end{exercise}

\begin{exercise} \label{exercises-exercise-surjective-map-affine-variety-projective-variety} $k$를 체라고 하고, $X$를 $k$ 위의 사영 다양체라고 하자. $k$ 위의 아핀 다양체 $U$와 다양체의 전사 $U \to X$가 존재함을 보이라. \end{exercise}

\begin{exercise} \label{exercises-exercise-vandermonde} $k$을 표수 $p > 0$의 체로 하고, 이는 $2,3$와 다르다고 한다. $\mathbf{P}^n_k$의 닫힌 부분 스킴 $X$을 $$
\sum\nolimits_{i = 0, \ldots, n} X_i = 0,\quad
\sum\nolimits_{i = 0, \ldots, n} X_i^2 = 0,\quad
\sum\nolimits_{i = 0, \ldots, n} X_i^3 = 0
$$에 의해 정의한다. % P12-ERR-0497: preserve the frozen closed-subsheme/variety terminology locus.
의 쌍 $(n, p)$에 대해 이 닫힌 부분 스킴은 특이한가. \end{exercise}



\section{가환대수, 2016년 가을 학기 기말시험} \label{exercises-section-final-exam-fall-2016}

\noindent 이하 내용은 Columbia 대학에서 2016년 가을에 진행된 가환대수 강의의 기말시험 문제이다.

\begin{exercise}[정의] \label{exercises-exercise-definitions-fall-2016} $R$를 환으로 하자. 이탤릭체로 표시된 개념을 정의하시오. \begin{enumerate} \item 소이데알 $\mathfrak p$에서의 $R$의 {\it 국소환}, \item {\it 유한} $R$-가군, \item {\it 유한 표현} $R$-가군, \item $R$는 {\it 정칙}이다, \item $R$는 {\it 사슬형}이다, \item $R$는 {\it Cohen--Macaulay} 이다. \end{enumerate} \end{exercise}

\begin{exercise}[정리] \label{exercises-exercise-results-fall-2016} 각 항목과 관련된, 강의에서 논의한 비자명한 사실을 정확하게 서술하시오. \begin{enumerate} \item 정칙환, \item Cohen--Macaulay 가군의 수반소 아이디얼, \item 체 위 유한형 적분영역의 차원, \item Chevalley의 정리. \end{enumerate} \end{exercise}

\begin{exercise} \label{exercises-exercise-finite-injective} $A \to B$을 다음 조건을 만족하는 환 준동형으로 한다: \begin{enumerate} \item $A$은 극대 아이디얼 $\mathfrak m$을 가진 국소환이고; \item $A \to B$은 유한한 \footnote{이는, $B$이 $A$-가군으로서 유한함을 의미한다. } 환경 준동형이다； \item $A \to B$는 단사이다 ($A$를 $B$의 부분환으로 간주한다). \end{enumerate} $\mathfrak m = A \cap \mathfrak q$가 되는 소아이디얼 $\mathfrak q \subset B$가 존재함을 보이라. \end{exercise}

\begin{exercise} \label{exercises-exercise-find-components} $k$를 체로 하고, $R = k[x, y, z, w]$라고 하자. 아이디얼 $I = (xy, xz, xw)$를 생각한다. $V(I) \subset \Spec(R)$의 기약 성분은 무엇인가. 또한 그것들의 차원은 얼마인가. \end{exercise}

\begin{exercise} \label{exercises-exercise-no-nonconstant-morphism} $k$를 체로 하고, $A = k[x, x^{-1}]$ 및 $B = k[y]$라고 하자. 임의의 $k$-대수 준동형 $\varphi : A \to B$이 $x$를 상수로 보낸다는 것을 보이라. \end{exercise}

\begin{exercise} \label{exercises-exercise-regular-over-Z} 환 $R = \mathbf{Z}[x, y]/(xy - 7)$를 생각하자. $R$이 정칙임을 증명하라. \end{exercise}

\noindent 뇌터 국소환 $(R, \mathfrak m, \kappa)$가 주어졌을 때, $n \geq 0$에 대해 $\varphi_R(n) = \dim_\kappa(\mathfrak m^n/\mathfrak m^{n + 1})$라고 둔다.

\begin{exercise} \label{exercises-exercise-hilbert-function-allowed} 모든 $n \geq 0$에 대해 $\varphi_R(n) = n + 1$가 되는 뇌터 국소환 $R$은 존재하는가. \end{exercise}

\begin{exercise} \label{exercises-exercise-hilbert-function-prohibited} $R$를 뇌터 국소환으로 한다. $\varphi_R(0) = 1$,$\varphi_R(1) = 3$, $\varphi_R(2) = 5$라고 가정하자. $\varphi_R(3) \leq 7$임을 보여라. \end{exercise}



\section{스킴론, 2017년 봄학기 기말시험} \label{exercises-section-final-exam-spring-2017}

\noindent 다음은 컬럼비아 대학교에서 2017년 봄에 진행된 스킴론 강의의 기말시험 문제이다.

\begin{exercise}[정의] \label{exercises-exercise-definitions-spring-2017} $f : X \to Y$를 스킴의 사상으로 한다. 이탤릭체로 표시된 개념을 간단히 정의하라. \begin{enumerate} \item $y \in Y$에서 $f$의 {\it 스킴론적 섬유}, \item $f$는 {\it 유한사상}이다, \item {\it 준연접} $\mathcal{O}_X$-가군, % P12-ERR-0498: preserve the frozen missing-article variety locus.
\item $X$는 {\it 다양체}이다, \item $f$는 {\it 매끄러운 사상}이다, \item $f$는 {\it 고유 사상}이다. \end{enumerate} \end{exercise}

\begin{exercise}[정리] \label{exercises-exercise-results-spring-2017} 각 항목과 관련된, 강의에서 논의한 자명하지 않은 사실을 정확하고 간결하게 서술하라. \begin{enumerate} \item 준연접 층의 순상, \item 사영 다양체 위의 연접층의 코호몰로지, \item 체 위의 사영 스킴에 대한 Serre 쌍대, \item Riemann--Hurwitz 공식. \end{enumerate} \end{exercise}

\begin{exercise} \label{exercises-exercise-compute-degree} $k$를 대수적으로 닫힌 체로 하자. $\ell > 100$을, $k$의 특성과 다른 소수로 하자. $X$을 $\mathbf{A}^2_k$에서의 방정식 $$
y^\ell = x(x - 1)^3
$$로 주어지는 아핀 곡선의 비특이 사영 모델로 하자. 다음 질문에 답하라: \begin{enumerate} \item $X$의 종수는 무엇인가. \item $X$의 구너리티 \footnote{구너리티란, $X$에서 $\mathbf{P}^1_k$으로 가는 비정상 사상의 차수의 최소값이다.}의 상한을 하나 제시하라. \end{enumerate} \end{exercise}

\begin{exercise} \label{exercises-exercise-count-components} $k$을 대수폐쇄체라고 하자. $X$을 $k$ 위의 가환사상 사영 스킴라고 하고, 그 모든 가환 성분은 같은 차원 $1$을 가진다고 하자. $\omega_{X/k}$을 상대 쌍대화 모듈이라고 하자. $\dim_k H^1(X, \omega_{X/k}) > 1$라면 $X$이 비연결임을 보여라. \end{exercise}

\begin{exercise} \label{exercises-exercise-sections-L-L-inverse} 스킴 $X$과 비자명한 가역 $\mathcal{O}_X$-가군 $\mathcal{L}$이고, $H^0(X, \mathcal{L})$과 $H^0(X, \mathcal{L}^{\otimes -1})$ 모두가 0이 아닌 것의 예를 들어라. \end{exercise}

\begin{exercise} \label{exercises-exercise-in-product} $k$를 대수적으로 닫힌 체로 하고, $g \geq 3$으로 한다. $X$과 $X'$을 각각 종수 $g$과 $g + 1$을 가진 $k$ 위의 매끄러운 사영 곡선으로 하자. $Y \subset X \times X'$을 곡선이라고 하고, 사영 $Y \to X$과 $Y \to X'$는 비상수라고 하자. $Y$의 비특이 사영 모델의 종수가 $\geq 2g + 1$임을 증명하라. \end{exercise}

\begin{exercise} \label{exercises-exercise-finitely-many} $k$을 유한체라고 하고, $g > 1$라고 하자. 다음 주장 증명의 개요를 서술하라: $k$ 위의 종수 $g$의 매끄러운 사영 곡선의 동형류는 유한개뿐이다. (시도만 해도 점수가 주어진다.) \end{exercise}



\section{가환대수, 2017년 가을 학기 기말시험} \label{exercises-section-final-exam-fall-2017}

\noindent 다음은 콜럼비아 대학교에서 2017년 가을에 진행된 가환대수 강의의 기말시험 문제이다.

\begin{exercise}[정의] \label{exercises-exercise-definitions-fall-2017} 이탤릭체로 표시된 개념을 간결하게 정의하시오. \begin{enumerate} \item 준동형 $F : \mathcal{A} \to \mathcal{B}$의 {\it 왼쪽 수반}, \item 체의 확대 $L/K$의 {\it 초월 차수}, \item 고전적 아핀 다양체 $X \subset k^n$ 위의 {\it 정칙 함수}, \item 위상 공간 위의 {\it 층}, \item {\it 국소환}, \item 스킴 사상 $f : X \to Y$이 {\it 아파인}임. \end{enumerate} \end{exercise}

\begin{exercise}[정리] \label{exercises-exercise-results-fall-2017} 각 항목에 관련된, 강의에서 논의한 비자명한 사실을 정확하고 간결하게 서술하라 (여러 개가 있는 경우에는 그 중 하나를 선택하면 된다). \begin{enumerate} \item Yoneda 보조정리, \item Mayer--Vietoris, \item 차원과 코호몰로지, \item Hilbert 다항식, \item 사영 공간에 대한 쌍대성. \end{enumerate} \end{exercise}

\begin{exercise} \label{exercises-exercise-compute-dimension} $k$를 대수적으로 닫힌 체라고 하자. Zariski 위상과 좌표 $x_1, x_2, x_3, x_4, x_5$를 가진 $k^5$의 닫힌 부분집합 $X$을 방정식 $$
x_1^2 - x_4 = 0,\quad
x_2^5 - x_5 = 0,\quad
x_3^2 + x_3 + x_4 + x_5 = 0
$$으로 정의한다. $X$의 차원은 무엇인가. 또한, 그것은 왜 그런가. \end{exercise}

\begin{exercise} \label{exercises-exercise-can-there-be} $k$를 체로 한다. $X = \mathbf{P}^1_k$을 $k$ 위의 차원 $1$의 사영 공간으로 하자. $\mathcal{E}$을 유한 국소 자유 $\mathcal{O}_X$-가군으로 하자. % P12-ERR-0551: preserve the frozen missing-by Serre-twist declaration locus.
$d \in \mathbf{Z}$에 대하여, $\mathcal{E}(d) = \mathcal{E} \otimes_{\mathcal{O}_X} \mathcal{O}_X(d)$을 $\mathcal{E}$의 제 $d$ Serre 꼬임이라 하고, $h^i(X, \mathcal{E}(d)) = \dim_k H^i(X, \mathcal{E}(d))$이라 하자. \begin{enumerate} \item $h^0(X, \mathcal{E}) = 5$ 및 $h^0(X, \mathcal{E}(1)) = 4$를 만족하는 $\mathcal{E}$이 존재하지 않는 이유는 무엇인가. \item $h^1(X, \mathcal{E}(1)) = 5$ 그리고 $h^1(X, \mathcal{E}) = 4$을 만족하는 $\mathcal{E}$가 존재하지 않는 이유는 무엇인가. \item 어떤 $a \in \mathbf{Z}$에 대해, $X$ 위의 벡터 다발 $\mathcal{E}$ 중 $$
\begin{matrix}
h^0(X, \mathcal{E})\phantom{(1)} = 1 &
h^1(X, \mathcal{E})\phantom{(1)} = 1 \\
h^0(X, \mathcal{E}(1)) = 2 &
h^1(X, \mathcal{E}(1)) = 0 \\
h^0(X, \mathcal{E}(2)) = 4 &
h^1(X, \mathcal{E}(2)) = a
\end{matrix}
$$을 만족하는 것이 존재할 수 있는가. \end{enumerate} 부분적인 답변도 환영하며, 권장한다. \end{exercise}

\begin{exercise} \label{exercises-exercise-banana} $X$을, 두 개의 닫힌 부분집합 $Y$과 $Z$의 합 $X = Y \cup Z$인 위상 공간으로 하고, 그들의 공통 부분을 $W = Y \cap Z$라고 쓴다. % P12-ERR-0552: preserve the frozen compressed "Denote ... the inclusion maps" locus.
포함 사상을 $i : Y \to X$, $j : Z \to X$, $k : W \to X$라고 쓴다. \begin{enumerate}\item 층의 짧은 완전열 $$
0 \to \underline{\mathbf{Z}}_X \to
i_*(\underline{\mathbf{Z}}_Y) \oplus
j_*(\underline{\mathbf{Z}}_Z) \to
k_*(\underline{\mathbf{Z}}_W) \to 0
$$가 존재함을 보이시오. 여기서 $\underline{\mathbf{Z}}_X$는 $X$ 위의 값 $\mathbf{Z}$의 상수층을 나타내며, 다른 것도 마찬가지이다. \item $\mathbf{Z}$-계수의 $X$, $Y$, $Z$, $W$의 코호몰로지 군의 관계에 대해 무엇을 결론지을 수 있는가. \end{enumerate} \end{exercise}

\begin{exercise} \label{exercises-exercise-cohomology-infinite-punctured} $k$를 체로 한다. $A = k[x_1, x_2, x_3, \ldots]$을 무한 개의 변수의 다항식환으로 한다. % P12-ERR-0553: preserve the frozen missing-by maximal-ideal declaration locus.
$\mathfrak m$를, 모든 변수에서 생성되는 $A$의 극대 아이디어로 하자. $X = \Spec(A)$, $U = X \setminus \{\mathfrak m\}$라고 두자. \begin{enumerate} \item $H^1(U, \mathcal{O}_U) = 0$임을 보이라. 힌트: {\v C}ech 코호몰로지에 의한 계산. \item $i \geq 1$에 대한 $H^i(U, \mathcal{O}_U)$를 무엇으로 예상하는가. \end{enumerate} \end{exercise}

\begin{exercise} \label{exercises-exercise-principal} $A$을 국소환으로 하고, $a \in A$를 영이 아닌 인자로 한다. $I, J \subset A$를 $IJ = (a)$를 만족하는 아이디어로 한다. 아이디어 $I$가 단항, 즉 하나의 원소로 생성됨을 보이라 (그 원소는 영이 아닌 인자가 된다). \end{exercise}



\section{스킴론, 2018년 봄 학기 기말 시험} \label{exercises-section-final-exam-spring-2018}

\noindent 다음은 콜럼비아 대학교에서 2018년 봄에 진행된 스킴론 강의의 기말 시험 문제이다.

\begin{exercise}[정의] \label{exercises-exercise-definitions-spring-2018} 이탤릭체로 표시된 개념을 간결하게 정의하시오. $k$를 대수적으로 닫힌 체라고 하자. $X$를 $k$ 위의 사영 곡선이라고 하자. \begin{enumerate} \item $k$ 위의 {\it 매끄러운} 대수, \item $X$ 위의 가역 $\mathcal{O}_X$-가군의 {\it 횟수}, \item $X$의 {\it 종수}, \item $X$의 {\it Weil 인자류군}, \item $X$은 {\it 초타원적}이며, \item $k$ 위의 매끄러운 사상 곡면 위의 두 곡선의 {\it 교점 수}. \end{enumerate} \end{exercise}

\begin{exercise}[정리] \label{exercises-exercise-results-spring-2018} 각 항목과 관련된, 강의에서 논의한 자명하지 않은 사실을 정확하고 간결하게 서술하라 (여러 개가 있는 경우에는 그 중 하나를 선택하면 된다). \begin{enumerate} \item 리만–허비츠 정리, \item 클리포드의 정리, \item 매끄러운 사영 곡면 사이의 사상의 분해, \item 호지의 지수 정리, 그리고 \item 유한체 위의 곡선에 대한 Riemann 추측. \end{enumerate} \end{exercise}

\begin{exercise} \label{exercises-exercise-hyperelliptic} $k$을 대수적 폐체로 하자. $X \subset \mathbf{P}^3_k$를 차수 $d$, 종수 $\geq 2$의 매끄러운 곡선으로 하자. $X$은 평면에 포함되지 않으며, $\mathbf{P}^3_k$ 내에서 $X$과 $d - 2$ 개의 점에서 교차하는 직선 $\ell$이 존재한다고 가정하자. $X$가 초타원적임을 보이라. \end{exercise}

\begin{exercise} \label{exercises-exercise-singular} $k$을 대수적으로 닫힌 체라고 하자. $X$를 서로 다른 특이점 $p_1, \ldots, p_n$를 가지는 사영 곡선이라고 하자. $X$의 정규화의 종수가 많아야 $-n + \dim_k H^1(X, \mathcal{O}_X)$인 이유를 설명하라. \end{exercise}

\begin{exercise} \label{exercises-exercise-how-many-blowups} $k$를 몸체로 한다.$X = \Spec(k[x, y])$를 아핀 $2$ 차원 공간으로 한다. $$
I = (x^3, x^2y, xy^2, y^3) \subset k[x, y].
$$ 라 하자. $Y \subset X$를 $I$에 대응하는 닫힌 부분 스킴으로 한다. $b : X' \to X$을 아이디얼 $(x, y)$의 블로업, 즉 원점에서의 아핀 공간의 블로업으로 한다. \begin{enumerate} \item 스킴론적 역상 $b^{-1}Y \subset X'$이 유효 카르티에르 인자임을 보이라. % P12-ERR-0499: preserve the frozen "Given an example" imperative locus.
\item $I \subset J \subset (x, y)$을 만족하는 아이디얼 $J \subset k[x, y]$로서, $J$에 대응하는 닫힌 부분 스킴을 $Z \subset X$라고 하면, 스킴론적 역상 $b^{-1}Z$가 유효 카르티에 인자가 되지 않는 예를 들어라. \end{enumerate} \end{exercise}

\begin{exercise} \label{exercises-exercise-rational-surface} $k$를 대수폐체로 하자. 다음과 같은 형태의 곡면을 고려한다: \begin{enumerate} \item $C_1$과 $C_2$가 매끄러운 사영 곡선인 $S = C_1 \times C_2$, \item $C_1$ 그리고 $C_2$이 매끄러운 사영 곡선이며, 또한 $C_1$의 종수가 $> 0$인 $S = C_1 \times C_2$, \item차수 $4$의 초곡면 $S \subset \mathbf{P}^3_k$, 그리고 \item차수 $4$의 매끄러운 초곡면 $S \subset \mathbf{P}^3_k$. \end{enumerate} 각 형에 대해, 그 형의 곡면류가 유리 곡면을 포함하는지 여부와 그 이유를 간단히 서술하라. \end{exercise}

\begin{exercise} \label{exercises-exercise-self-square-line} $k$를 대수적으로 닫힌 체로 한다. $S \subset \mathbf{P}^3_k$을 차수 $d$의 매끄러운 초곡면으로 한다. $S$이 직선 $\ell$을 포함한다고 가정한다. % P12-ERR-0500: preserve the frozen nonstandard "self square" source phrase.
$S$ 위의 인수로 본 $\ell$의 자가 교차점 수는 얼마인가. \end{exercise}



\section{가환대수, 2019년 가을 학기 기말 시험} \label{exercises-section-final-exam-fall-2019}

\noindent 다음은 콜럼비아 대학교에서 2019년 가을에 진행된 가환대수 강의의 기말 시험 문제이다.

\begin{exercise}[정의] \label{exercises-exercise-definitions-fall-2019} 이탤릭체로 표시된 개념을 간결하게 정의하라. \begin{enumerate} \item 뇌터 정리 공간의 {\it 구성 가능한 부분집합}, \item 소아이디얼 $\mathfrak p$에서의 $R$-가군 $M$의 {\it 국소화}, \item 뇌터 국소환 $(A, \mathfrak m, \kappa)$ 위의 가군의 {\it 길이}, \item 환 $R$ 위의 {\it 사영 가군}, 그리고 \item 뇌터 국소환 $(A, \mathfrak m, \kappa)$ 위의 {\it Cohen--Macaulay} 가군. \end{enumerate} \end{exercise}

\begin{exercise}[정리] \label{exercises-exercise-results-fall-2019} 각 항목과 관련된, 강의에서 논의한 비자명한 사실을 정확하고 간결하게 서술하라 (복수인 경우에는 그 중 하나를 선택하면 된다). \begin{enumerate} \item 구성 가능한 집합의 상, \item Hilbert의 영점 정리, \item 체 위의 유한형 대수의 차원, \item 뇌터의 정규화, 그리고 \item 정칙 국소환. \end{enumerate} \end{exercise}

\noindent 환 $R$과 아이디얼 $I \subset R$에 대하여, $V(I)$이 $I \subset \mathfrak p$을 만족하는 $\mathfrak p \in \Spec(R)$의 집합을 나타내는 것을 상기하자.

\begin{exercise}[소아이디얼의 구성] \label{exercises-exercise-infinitely-many-primes} $V(\mathfrak p)$이 $(x, y)$과 $(x - 1, y - 1)$를 포함하도록 서로 다른 소아이디얼 $\mathfrak p \subset \mathbf{C}[x, y]$을 무한히 구성하라. \end{exercise}

\begin{exercise}[소아이디얼 없음] \label{exercises-exercise-no-prime} $R = \mathbf{C}[x, y, z]/(xy)$라고 하자. $V(\mathfrak p)$이 $(x, y - 1, z - 5)$과 $(x - 1, y, z - 7)$를 포함하도록 하는 소아이디얼 $\mathfrak p \subset R$은 존재하지 않음을 간결하게 논하라. \end{exercise}

\begin{exercise}[Frobenius] \label{exercises-exercise-frobenius} $p$를 소수라고 하자 (식을 단순하게 하기 위해 $p = 2$라고 가정해도 좋다). 환 $R$를 취하고, $R$에서 $p = 0$라고 가정한다. \begin{enumerate} \item 사상 $F : R \to R$, $x \mapsto x^p$가 환 준동형임을 보여라. \item $\Spec(F) : \Spec(R) \to \Spec(R)$이 항등 사상임을 보여라. \end{enumerate} \end{exercise}

\noindent 위상 공간의 점의 {\it 특수화} $x \leadsto y$란 단순히 $y$가 $x$의 폐포에 속함을 의미한다는 것을 기억하자. $x$ 및 $y$와 다른 $z$에서, $x \leadsto z$이고 $z \leadsto y$인 것이 존재하지 않을 때, $x \leadsto y$를 {\it 직전 특수화}라고 한다.

\begin{exercise}[차원] \label{exercises-exercise-dimension} $5$ 개의 서로 다른 점 $x, y, z, u, v$을 포함하는 sober 위상 공간 $X$이 있고, 그 특수화가 다음과 같다고 하자: $$
\xymatrix{
x \ar[d] \ar[r] & u & v \ar[l] \ar[dl] \\
y \ar[r] & z
}
$$ 이러한 $X$가 가질 수 있는 최소 차원은 얼마인가. $X$가 체 위에서 유한형인 대수의 스펙트럼이고, $x \leadsto u$가 바로 이전 특수화라면, 특수화 $v \leadsto z$에 대해 무엇을 말할 수 있는가. \end{exercise}

\begin{exercise}[Tor의 계산] \label{exercises-exercise-tor-computation} $R = \mathbf{C}[x, y, z]$로 하자. $M = R/(x, z)$ 및 $N = R/(y, z)$로 두자. 어떤 $i \in \mathbf{Z}$에 대해 $\text{Tor}_i^R(M, N)$은 0이 아닌가. \end{exercise}

\begin{exercise} \label{exercises-exercise-depth-goes-up} $A \to B$를 뇌터 국소환의 평탄한 국소 준동형으로 하자. $A$의 깊이가 $k$라면, $B$의 깊이는 $k$ 이상임을 보여라. \end{exercise}



\section{대수기하, 2020년 봄 학기말 시험} \label{exercises-section-final-exam-spring-2020}

\noindent 아래는 콜럼비아 대학교에서 2020년 봄에 진행된 대수기하 강의의 학기말 시험 문제이다.

\begin{exercise}[정의] \label{exercises-exercise-definitions-spring-2020} 이탤릭체로 표시된 개념을 간단히 정의하라. \begin{enumerate} \item {\it 스킴}, \item {\it 스킴의 사상}, \item 스킴 위의 {\it 준연접 가군}, \item 체 $k$ 위의 {\it 다양체}, \item 체 $k$ 위의 {\it 곡선}, \item 스킴의 {\it 유한 사상}, \item 위상 공간 $X$ 위의 Abel 군의 층 $\mathcal{F}$의 {\it 코호몰로지}, \item 체 $k$ 위 특이적이고 차원 $d$의 스킴 $X$ 위의 {\it 쌍대화 층}, 및 \item 다양체 $X$에서 다양체 $Y$로의 {\it 유리 사상}. \end{enumerate} \end{exercise}

\begin{exercise}[정리] \label{exercises-exercise-results-spring-2020} 각 항목과 관련된, 강의에서 논의한 비자명한 사실을 정확하고 간결하게 서술하라 (복수인 경우에는 그 중 하나를 선택하면 된다). \begin{enumerate} \item 차원 $d$의 뇌터 위상 공간 $X$ 위의 Abel 층의 코호몰로지, \item 체 $k$ 위의 부드러운 다양체의 미분층 $\Omega^1_{X/k}$, \item 체 $k$ 위의 매끄러운 사영 다양체 $X$의 쌍대화 층 $\omega_X$, \item 대수적 닫힌 체 $k$ 위의 매끄럽고 고유한 종수 $0$의 곡선, 그리고 \item 차수 $d$의 평면 곡선의 종수. \end{enumerate} \end{exercise}

\begin{exercise} \label{exercises-exercise-coh-P1-infty-doubled} $k$를 체로 한다. $X$를 $k$ 위의 스킴으로 한다. $X = X_1 \cup X_2$는 개피복이며, $X_1$과 $X_2$는 모두 $\mathbf{P}^1_k$와 동형이고, $X_1 \cap X_2$는 $\mathbf{A}^1_k$와 동형이라고 가정한다 (이러한 스킴은 존재한다. 예를 들어 $\mathbf{P}^1_k$의 $\infty$을 이중화하면 된다). $\dim_k H^1(X, \mathcal{O}_X)$가 무한함을 증명하라. \end{exercise}

\begin{exercise} \label{exercises-exercise-no-morphism} $k$를 대수적으로 닫힌 체라고 하자. $Y$를 종수 $10$의 매끄러운 사영 곡선이라고 하자. 동형이 아닌 비상수 사상 $f : X \to Y$이 존재하는 매끄러운 사영 곡선 $X$의 종수에 대해, 좋은 하한을 구하라. \end{exercise}

\begin{exercise} \label{exercises-exercise-element-order-n-in-pic} $k$을 특성 $0$의 대수적으로 닫힌 체라고 하자. $$
X : T_0^d + T_1^d - T_2^d = 0 \subset \mathbf{P}^2_k
$$을 차수 $d \geq 3$의 페르마 곡선이라고 하자. $X$ 위의 닫힌 점 $p = [1 : 0 : 1]$과 $q = [0 : 1 : 1]$를 생각하고, $D = [p] - [q]$라고 두자. \begin{enumerate} \item $D$가 Weil 유수군에서 비자명함을 증명하시오. \item $d D$가 Weil 유수군에서 자명함을 증명하시오. (힌트: $d[p]$과 $d[q]$가 모두 $X$와 평면 내의 어떤 직선과의 교점을 이루는 것을 증명해 보시오.) ) \end{enumerate} \end{exercise}

\begin{exercise} \label{exercises-exercise-number-equations} $k$를 대수적으로 닫힌 체로 하자. $2$ 다음 Veronese 매핑 $$
\varphi : \mathbf{P}^2 \longrightarrow \mathbf{P}^5
$$을 고려한다. 강의 내용과 표기법에 따르면, 이는 사상 $$
\varphi = \varphi_{\mathcal{O}_{\mathbf{P}^2}(2)} :
\mathbf{P}^2
\longrightarrow
\mathbf{P}(\Gamma(\mathbf{P}^2, \mathcal{O}_{\mathbf{P}^2}(2)))
$$이다. 동차 좌표로는, 이는 $$
[a_0 : a_1 : a_2] \longmapsto
[a_0^2 : a_0a_1 : a_0a_2 : a_1^2 : a_1a_2 : a_2^2]
$$로 주어진다. 이는 닫힌 매핑이다(이 사실은 그대로 사용해도 된다). $I \subset k[T_0, \ldots, T_5]$를 $\varphi(\mathbf{P}^2)$의 동차 아이디얼이라 하자. 즉, 동차 부분 $I_d$의 원소는, 닫힌 부분 스킴 $\varphi(\mathbf{P}^2)$ 위에서 0으로 제한되는 차수 $d$의 동차 다항식 $F(T_0, \ldots, T_5)$이다. $\dim_k I_d$를 $d$의 함수로 계산하라. \end{exercise}

\begin{exercise} \label{exercises-exercise-dualizing-sheaf-rank-1} $k$를 대수 닫힌 체로 한다. $X$를, $k$ 위의 차원 $d$의 유한 스킴으로, 쌍대화 모듈 $\omega_X$을 갖는 것으로 한다. 다음 정보가 주어져 있다: \begin{enumerate} \item 임의의 $i$ 및 임의의 연접 $\mathcal{O}_X$-가군 $\mathcal{F}$에 대해, $\text{Ext}^i_X(\mathcal{F}, \omega_X) \times
H^{d - i}(X, \mathcal{F}) \to H^d(X, \omega_X) \xrightarrow{t} k$는 비퇴화이며, 그리고 \item $\omega_X$는 어떤 계수 $r$를 갖는 유한 국소 자유 가군이다. \end{enumerate} $r = 1$임을 보여라 (힌트: 닫힌 점에 지지점을 갖는 적절한 가군 $\mathcal{F}$을 선택했을 때 무슨 일이 일어나는지 조사하라). \end{exercise}



\section{가환대수, 2021년 가을 학기 기말 시험} \label{exercises-section-final-exam-fall-2021}

\noindent 다음은 콜럼비아 대학교에서 2021년 가을에 진행된 가환대수 강의의 기말 시험 문제이다.

\begin{exercise}[정의] \label{exercises-exercise-definitions-fall-2021} 이탤릭체로 표시된 개념을 간결하게 정의하시오. \begin{enumerate} \item 환 $A$의 {\it 곱셈적 부분집합}, \item {\it Artin 환} $A$, \item 위상 공간으로서의 환 $A$의 {\it 스펙트럼}, \item {\it 평탄한 환 준동형} $A \to B$, \item $A$의 소 아이디얼 $\mathfrak p$의 {\it 높이}, 및 \item 환 $A$ 상의 함자 {\it $\text{Tor}^A_i(-, -)$}. \end{enumerate} \end{exercise}

\begin{exercise}[정리] \label{exercises-exercise-results-fall-2021} 각 항목에 관련된, 강의에서 논의한 비자명한 사실을 정확하고 간결하게 서술하시오 (여러 가지가 있는 경우에는 그중 하나를 선택하면 된다). \begin{enumerate} \item Artin 환, \item 평탄성과 소 아이디얼, \item 뇌터 국소환 $(A, \mathfrak m)$에 대한 $A/\mathfrak m^n$의 길이, \item 보편 사슬형 뇌터 환에 대한 차원 공식, \item 뇌터 국소환의 완비화, 그리고 \item Artin 국소환에 대한 Matlis 쌍대성. \end{enumerate} \end{exercise}

\begin{exercise}[단원] \label{exercises-exercise-simple-units} Abel 군으로서, $\mathbf{Z}[x, 1/x]$의 단원군은 어떤 구조를 가지는가. 설명은 필요 없다. \end{exercise}

\begin{exercise}[아이디얼] \label{exercises-exercise-ideals} $A = \mathbf{F}_2[x, y]/(x^2, xy, y^2)$라고 하자. % P12-ERR-0554: preserve the frozen missing-by notation clause.
$\overline{x}$과 $\overline{y}$을 각각 $A$에서 $x$과 $y$의 상으로 하자. $A$의 아이디얼을 나열하라. 설명은 필요 없다. \end{exercise}

\begin{exercise}[Tor와 Ext] \label{exercises-exercise-compute-tor-ext} $(A, \mathfrak m, \kappa)$를 뇌터 지역환으로 하자. $\varphi(n) = \dim_\kappa \mathfrak m^n/\mathfrak m^{n + 1}$라고 하자. \begin{enumerate} \item $\text{Tor}_1^A(A/\mathfrak m^n, \kappa)$가 $\kappa$-벡터 공간으로서 차원 $\varphi(n)$을 갖는 것을 보이라. \item $\text{Ext}^1_A(A/\mathfrak m^n, \kappa)$이 $\kappa$-벡터 공간으로서 차원 $\varphi(n)$을 가짐을 보이라. \end{enumerate} \end{exercise}

\begin{exercise}[두 벡터] \label{exercises-exercise-two-vectors} % P12-ERR-0501: preserve the frozen b3 generator token.
$A = \mathbf{Z}[a_1, a_2, a_3, b_1, b_2, b3]$라고 하자. $A^{\oplus 3}$에서 $a = (a_1, a_2, a_3)$ 및 $b = (b_1, b_2, b_3)$라고 둔다. 집합 $$
Z = \{\mathfrak p \in \Spec(A) \mid a, b
\text{의 상들이 다음 벡터공간에서 선형 종속이다: }
\kappa(\mathfrak p)^{\oplus 3}\}
$$를 고려하자. \begin{enumerate} % P12-ERR-0502: preserve the frozen extraneous-article source locus.
\item $Z$가 $\Spec(A)$의 닫힌 부분집합임을 증명하라. \item $Z$의 차원 $\dim(Z)$은 얼마인가. \item $\mathbf{Z}$를 체로 바꾸면, $\dim(Z)$는 어떻게 되는가. \end{enumerate} \end{exercise}

\begin{exercise}[단사 가군] \label{exercises-exercise-injective-artinian-local} $(A, \mathfrak m, \kappa)$를 Artin 국소환으로 하자. $A$는 $A$-가군으로서 사상적임을 가정하자. % P12-ERR-0503: preserve the frozen malformed as/has source phrase.
$\Hom_A(\kappa, A)$가 $\kappa$-벡터 공간으로서 차원 $1$을 갖는다는 것을 보이시오. \end{exercise}


\section{대수기하, 2022년 봄학기 기말시험} \label{exercises-section-final-exam-spring-2022}

\noindent 다음은 콜럼비아 대학에서 2022년 봄에 진행된 대수기하 강의의 기말시험 문제이다.

\begin{exercise}[정의] \label{exercises-exercise-definitions-spring-2022} 이탤릭체로 표시된 개념을 간단히 정의하라. \begin{enumerate} \item {\it 스킴}, \item 스킴 $X$ 위의 {\it 준연접 가군}, \item 스킴의 {\it 평탄 사상} $X \to Y$, \item 스킴의 {\it 유한 사상} $X \to Y$, \item 기저 스킴 $S$ 위의 {\it 군 스킴} $G$, \item 기저 스킴 $S$ 위의 {\it 다양체족}, \item 체 $k$ 위의 다양체 $X$의 닫힌 점 $x$의 {\it 차수}, \item $\mathbf{P}^n(\mathbf{Q})$의 점 $p = (a_0 : \ldots : a_n)$의 일반적인 {\it 로그 높이}, 그리고 \item {\it $C_i$ 체}. \end{enumerate} \end{exercise}

\begin{exercise}[정리] \label{exercises-exercise-results-spring-2022} 각 항목과 관련된, 강의에서 논의한 비자명한 사실을 정확하고 간결하게 서술하라 (여러 가지가 있을 경우에는 그 중 하나를 선택하면 된다). \begin{enumerate} \item 스킴 $X$에서 아핀 스킴 $\Spec(A)$으로의 사상, \item 아핀 스킴 $X$ 위의 준연접 가군 $\mathcal{F}$의 코호몰로지, \item $k$이 체일 때의 $\mathbf{P}^1_k$의 Picard 군, \item $S$가 뇌터일 때 평탄하고 고유한 사상 $X \to S$의 섬유의 차원, \item 스킴 $S$ 위의 $\mathbf{G}_m$-변환 가군, 그리고 \item 교차에 관한 Bezout의 정리 (원하면 특별한 경우에 한정해도 좋다). \end{enumerate} \end{exercise}

\begin{exercise}[삼차 초곡면] \label{exercises-exercise-cubics} $F \in \mathbf{C}[T_0, \ldots, T_n]$을 차수 $3$의 동차식으로 한다. $3$ 개의 벡터 $x, y, z \in \mathbf{C}^{n + 1}$가 주어졌을 때, 조건 $$
(*)\quad
F(\lambda x + \mu y + \nu z) = 0 \text{이 다음 다항식환에서 성립한다: } \mathbf{C}[\lambda, \mu, \nu]
$$을 고려한다. \begin{enumerate} \item $x, y, z$의 선택 방법 전체 공간의 차원은 몇인가. % P12-ERR-0555: preserve the frozen subject-verb agreement locus.
\item 조건 (*)은 $x$, $y$, $z$의 좌표에 몇 개의 방정식을 부과하는가. \item (*)가 성립하는 세 쌍 $x, y, z$ 전체 공간의 기대 차원은 얼마인가. \item $x, y, z$가 선형종속인 세 쌍 전체 공간의 차원은 얼마인가. \item $\mathbf{P}^n$ 내의 차수 $3$의 초곡면에는, $n \geq a$라면 차원 $2$의 선형 부분공간이 존재할 것으로 기대된다는 것을 결론내리시오. 여기서 $a$는 스스로 구하시오. \end{enumerate} \end{exercise}

\begin{exercise}[높이] \label{exercises-exercise-heights} $K$를 체로 하자. 강의에서 논한 $2$개의 공리를 만족하는 함수의 집합 $h_n : \mathbf{P}^n(K) \to \mathbf{R}$, $n \geq 0$를 취하자. $X$를 $K$ 위의 사영 다양체라고 하자. $\mathcal{L}$를 가역 $\mathcal{O}_X$-가군이라고 하고, 강의에서는 연관된 높이 함수 $h_\mathcal{L} : X(K) \to \mathbf{R}$를 구성했던 것을 상기하자. $\alpha : X \to X$를 $K$ 위의 $X$의 자기동형이라고 하자. \begin{enumerate} \item $P \mapsto h_\mathcal{L}(\alpha(P))$가 함수 $h_{\alpha^*\mathcal{L}}$와 유계한 양만큼만 다른 것을 증명하라. (힌트: $\mathcal{L} = \varphi^*\mathcal{O}_{\mathbf{P}^n}(1)$를 만족하는 사상 $\varphi : X \to \mathbf{P}^n$이 존재할 때, 구성으로 인해 $h_\mathcal{L}(P) = h_n(\varphi(P))$임을 기억하고, 이것을 사용해 식을 여러 가지로 변형하라. 일반적으로 $\mathcal{L}$를 이 형태의 두 차이로 쓸 수 있다.) \item $h_\mathcal{L}(P) - h_\mathcal{L}(\alpha(P))$가 $X(K)$ 위에서 무한함을 가진다고 가정하자. $\mathcal{N} = \mathcal{L} \otimes \alpha^*\mathcal{L}^{\otimes -1}$에 대한 $h_\mathcal{N}$가 $X(K)$ 위에서 무한함을 가진다는 것을 증명하라. \item $X$는 타원곡선이며, $\mathcal{L}$는 $X$ 위의 대칭적이고 풍부한 가역 가군이고, $h_\mathcal{L}$가 $X(K)$ 위에서 비유계적이라고 가정하자. 차수 $0$의 가역 가군 $\mathcal{N}$가 존재하여 $h_\mathcal{N}$가 비유계적이 되도록 함을 보이라. (힌트: $X$는 차원 $1$의 아벨 다양체임을 상기하자. 따라서 강의 결과에 따라, $h_\mathcal{L}$은 상수 차이를 제외하고 이차적이다. 적절한 점 $P_0 \in X(K)$를 선택하라. $\alpha : X \to X$를 $P_0$에 의한 평행 이동으로 한다. $P \mapsto h_\mathcal{L}(P) - h_\mathcal{L}(P + P_0)$를 고려하고, 위에서 증명한 결과를 적용하라.) \end{enumerate} \end{exercise}

\begin{exercise}[모노사상] \label{exercises-exercise-monomorphisms} $f : X \to Y$를 스킴 범주에서의 모노사상으로 한다. 즉, 임의의 두 스킴 사상 $a, b : T \to X$에 대해, $f \circ a = f \circ b$이면 $a = b$라고 가정한다. $f$가 점 위에서 단사임을 증명하라. % P12-ERR-0504: preserve the frozen "Does you argument" source locus.
이 논의로부터 다른 무엇을 알 수 있는가. \end{exercise}

\begin{exercise}[불변점] \label{exercises-exercise-fix-points} $k$을 대수적으로 닫힌 체라고 하자. \begin{enumerate} \item $G = \mathbf{G}_{m, k}$라고 하자. $G$가 $k$ 위의 사영 다양체 $X$에 작용한다면, 그 작용은 불변점을 갖는 것을 보이시오. 즉, 모든 $g \in G(k)$에 대해 $a(g, x) = x$가 되는 점 $x \in X(k)$가 존재함을 증명하시오. \item $G = (\mathbf{G}_{m, k})^n$이 곱셈군의 $n \geq 1$개의 직적곱인 경우에도 같은 것을 증명하라. \item 매끄러운 사영 다양체 $X$ 위의 연결군 스킴 $G$의 작용 중, 불변점을 가지지 않는 예를 들어라. \end{enumerate} \end{exercise}


\section{대수기하, 2025년 봄 학기말 시험} \label{exercises-section-final-exam-spring-2025}

\noindent 아래는 콜럼비아 대학교에서 2025년 봄에 진행된 대수기하 강의의 학기말 시험 문제이다.

\begin{exercise}[정의] \label{exercises-exercise-definitions-spring-2025} 이탤릭체로 표시된 개념을 간단히 정의하라. \begin{enumerate} \item {\it 스킴} $X$, \item 스킴 $X$ 위의 {\it 준연접 가군}, \item 스킴 $X$의 {\it Picard 군}, % P12-ERR-0505: preserve the frozen "give a morphism" source locus.
\item 스킴의 사상 $f : X \to Y$과 $\mathcal{O}_X$-가군 $\mathcal{F}$가 주어졌을 때의 {\it 순상} $f_*\mathcal{F}$, \item 스킴의 {\it 닫힌 삽입}, 및 \item 주어진 체 $k$ 위의 {\it 다양체}. \end{enumerate} \end{exercise}

\begin{exercise}[정리] \label{exercises-exercise-results-spring-2025} % P12-ERR-0506: preserve the frozen "succintly" spelling locus.
각 항목과 관련하여, 강의에서 논의된 비자명한 사실을 하나 정확하고 간결하게 서술하라 (여러 개가 있는 경우 그 중 하나를 선택하면 된다). \begin{enumerate} \item 스킴 $X$의 위상 공간, \item 스킴의 범주 위의 함자에 대한 표현가능성 판정법, \item 대수적 폐체 $k$ 위의 매끄러운 사영 곡선 $X$의 Picard 함자, \item 대수적 폐체 $k$ 위의 매끄러운 사영 곡선 $X$의 Picard 군의 꼬임, 및 \item 대수적 폐체 $k$ 위의 매끄러운 사영 다양체 $X$,$Y$,$Z$의 곱 $X \times Y \times Z$의 Picard 군에 관한 정리. \end{enumerate} \end{exercise}

\begin{exercise} \label{exercises-exercise-affine-line-infinite-nr-points} $k$를 체로 한다.$\Spec(k[x])$가 무한개의 점을 가지는 이유를 설명하라. \end{exercise}

\begin{exercise} \label{exercises-exercise-hypersurface-variety} $k$를 체로 하고,$x_1, \ldots, x_n$를 변수로 한다. $f \in k[x_1, \ldots, x_n]$를 비상수 다항식으로 한다. $X = \Spec(k[x_1, \ldots, x_n]/(f))$ 멀리. $X$가 $k$ 위의 다양체가 되기 위해서는, $f$은 어떤 성질을 가져야 하는가. \end{exercise}

\begin{exercise} \label{exercises-exercise-find-singular-point} $\mathbf{C}$ 위의 다양체로 본 $\mathbf{C}[x, y]/(x^5 - 5xy^4 + 20y - 16)$의 스펙트럼상의 특이점을 하나 구하라(모두 구할 필요는 없다). \end{exercise}

\begin{exercise} \label{exercises-exercise-kernel-restriction-pic-curve} $X$ 을(를) 대수적 닫힌 체 $k$ 위의 매끄러운 사영 곡선이라고 하자. $x \in X(k)$ 을(를) $X$ 위의 닫힌 점이라고 하고, $U = X \setminus \{x\}$이라고 하자. 제한 사상 $\Pic(X) \to \Pic(U)$의 핵에 대하여 무엇을 말할 수 있는가. \end{exercise}

\begin{exercise} \label{exercises-exercise-sections-degree-2g-minus-4} $X$ 을(를), 종수 $g \geq 100$를 가지는 대수적으로 닫힌 체 $k$ 위의 매끄러운 사영 곡선이라고 하자. % P12-ERR-0507: preserve the frozen "for L be" source locus.
$X$ 위의 차수 $2g - 4$의 가역 $\mathcal{O}_X$-가군 $\mathcal{L}$에 대해, $h^0(\mathcal{L}) = \dim_k H^0(X, \mathcal{L})$는 어떤 값을 가질 수 있는가. 가능한 한 완전히 답하라. \end{exercise}
